% Leçon 34 — Vocabulaire et compréhension de texte : internet et les réseaux sociaux — Chinois Tle A — Cours signé OnBuch+
% Programme officiel MINESEC. Compiler avec : tectonic ZH34-vocabulaire-internet-et-les-reseaux-sociaux.tex
\newcommand{\DOCMATIERE}{Chinois}
\newcommand{\DOCNIVEAU}{Tle A}
\newcommand{\DOCMODULE}{Module 5 : Mass médias et télécommunication}
\newcommand{\DOCLECON}{ZH34}
\newcommand{\DOCTITRE}{Leçon 34 — Vocabulaire et compréhension de texte : internet et les réseaux sociaux}
% OnBuch+ — préambule « cours signés » ENRICHI (compilé avec tectonic / XeLaTeX)
% Système visuel « Pop » d'OnBuch+ : fond crème (jamais blanc), bordures encre
% épaisses, ombres « dures » décalées, titres Archivo Black en capitales.
% Version MATHÉMATIQUES : graphes (pgfplots/tikz), boîtes pédagogiques
% (définition, propriété, méthode, attention, le savais-tu, exercice résolu,
% exercice, TP, bilan), page de garde et signature OnBuch+ ancrée en bas.
%
% Macros attendues dans main.tex AVANT \input{preamble} :
%   \DOCMATIERE  \DOCNIVEAU  \DOCMODULE  \DOCLECON (numéro)  \DOCTITRE
\documentclass[11pt]{article}

\usepackage{fontspec}
\usepackage[french]{babel} % français = langue principale ; le chinois via \zh{..} / textebox (xeCJK)
\usepackage{xeCJK}
\usepackage{pifont} % \ding{51} \ding{55} utilisés par les modèles
\usepackage[a4paper,margin=2cm,top=3.4cm,bottom=2.8cm,headheight=1.4cm,footskip=1.3cm]{geometry}
\usepackage{amsmath,amssymb,amsfonts}
\usepackage{mathtools} % \xmapsto, \coloneqq... souvent utilisés par les modèles
\usepackage{cancel} % simplifications barrées (bilans de forces)
% \abs{z} (module, valeur absolue) et \norm{u} (norme), utilisés par les modèles
\providecommand{\abs}{}\let\abs\relax\DeclarePairedDelimiter{\abs}{\lvert}{\rvert}
\providecommand{\norm}{}\let\norm\relax\DeclarePairedDelimiter{\norm}{\lVert}{\rVert}
\usepackage[table]{xcolor} % [table] : fournit aussi \rowcolors (alternance de lignes)
\usepackage{pagecolor}
\usepackage{tikz}
\usetikzlibrary{arrows.meta,positioning,calc,shapes.geometric,shapes.misc,shapes.arrows,decorations.pathmorphing,decorations.pathreplacing,decorations.markings,patterns,shadows,mindmap,trees,fit,backgrounds,matrix,intersections,angles,quotes}
\usepackage{pgfplots}
\usepackage[version=4]{mhchem} % équations nucléaires \ce{^{235}_{92}U}
\usepackage[european,siunitx]{circuitikz} % schémas électriques
\pgfplotsset{compat=1.18}
\usepgfplotslibrary{fillbetween,groupplots} % aires entre courbes (calcul intégral)
\usepackage{enumitem}
\usepackage{fancyhdr}
% [most] charge toutes les bibliothèques tcolorbox (listings, minted, raster,
% documentation...) et gonfle inutilement la mémoire TeX sur un document long
% (~300 boîtes) : on ne charge que ce qui est réellement utilisé.
\usepackage[skins,breakable]{tcolorbox}
\usepackage{graphicx}
\usepackage{colortbl}
\usepackage{booktabs}
\usepackage{tabularx}
\usepackage{diagbox} % case d'en-tête diagonale des tableaux à double entrée
% Colonnes extensibles centrée / gauche / droite (C, L, R), souvent utilisées par les modèles
\newcolumntype{C}{>{\centering\arraybackslash}X}
\newcolumntype{L}{>{\raggedright\arraybackslash}X}
\newcolumntype{R}{>{\raggedleft\arraybackslash}X}
\usepackage{array}
\usepackage{multicol}
\usepackage{multirow}
\usepackage{siunitx}
% parse-numbers=false : accepte aussi \SI{2,5 \times 10^{-3}}{...}
\sisetup{locale=FR,output-decimal-marker={,},group-minimum-digits=4,parse-numbers=false}
\DeclareSIUnit\ohm{\ensuremath{\symup{\Omega}}} % U+2126 absent de la police
\DeclareSIUnit\division{div} % réglages d'oscilloscope (ms/div, V/div)
\DeclareSIUnit\tr{tr} % tours (vitesse de rotation, ex. \SI{1500}{\tr\per\minute})
\DeclareSIUnit\molar{M} % concentration molaire (ex. \SI{3}{\molar}, \SI{1}{\micro\molar})
\DeclareSIUnit\an{an} % année (le modèle confond parfois avec \year, un compteur TeX) — ex. \SI{5}{\kilo\meter\per\an}
\DeclareSIUnit\FCFA{FCFA} % franc CFA (ex. \SI{500}{\FCFA\per\kilo\gram})
\DeclareSIUnit\degreeBrix{\textdegree Brix} % teneur en sucre d'un sirop/jus (réfractométrie)
\DeclareSIUnit\month{mois} % le modèle utilise parfois \month, un compteur TeX, comme unité de durée
\usepackage{qrcode}
% \degree : commande fréquemment utilisée par les modèles pour un angle ou une
% température en dehors de \SI{}{} (non fournie par les paquets chargés).
\providecommand{\degree}{^{\circ}}
% \qed (fin de démonstration), utilisé par les modèles sans amsthm
\providecommand{\qed}{\hfill$\square$}
% \barre : double barre des valeurs interdites dans un tableau de variations (tkz-tab)
\providecommand{\barre}{\ensuremath{\|}}
% environnement proof (amsthm non chargé) utilisé par les modèles
\ifdefined\proof\else\newenvironment{proof}[1][Démonstration]{\par\noindent\textit{#1.}\ }{\qed\par}\fi
\usepackage{lastpage}
\usepackage{hyperref}
\hypersetup{hidelinks}

% ============================================================
% Palette « Pop » OnBuch+ — mêmes hex que la classe P (plus_ui.dart)
% ============================================================
\definecolor{poporange}{HTML}{F59321}
\definecolor{popdark}{HTML}{E07A0C}
\definecolor{popcream}{HTML}{FFF6E8}
\definecolor{popink}{HTML}{1C1714}
\definecolor{poppurple}{HTML}{7A5AE0}
\definecolor{popgreen}{HTML}{2E9E6A}
\definecolor{poppink}{HTML}{E0457B}
\definecolor{popblue}{HTML}{2F7DE1}
\definecolor{popgold}{HTML}{C98A12}
\definecolor{popmuted}{HTML}{8A7F76}
\definecolor{popline}{HTML}{EADFD2}
\definecolor{popgray}{HTML}{8A7F76}
\colorlet{popgrey}{popgray}
% Alias tolérants (le modèle nomme parfois les couleurs différemment)
\colorlet{popwhite}{white}
\colorlet{popblack}{popink}
\colorlet{popred}{poppink}
\colorlet{popyellow}{popgold}
% Teintes claires (fonds de tableaux/encadrés) — le modèle nomme parfois
% les couleurs avec un suffixe L (ex. popblueL) sans les avoir définies.
\colorlet{poporangeL}{poporange!20}
\colorlet{popdarkL}{popdark!20}
\colorlet{poppurpleL}{poppurple!20}
\colorlet{popgreenL}{popgreen!20}
\colorlet{poppinkL}{poppink!20}
\colorlet{popblueL}{popblue!20}
\colorlet{popgoldL}{popgold!20}
\colorlet{popredL}{popred!20}
\colorlet{popgrayL}{popgray!20}
% Teintes claires pour fonds de tableaux / zones
\colorlet{popblueL}{popblue!10!white}
\colorlet{poporangeL}{poporange!14!white}
\colorlet{popgreenL}{popgreen!12!white}
\colorlet{poppurpleL}{poppurple!12!white}
\colorlet{poppinkL}{poppink!12!white}
\colorlet{popgoldL}{popgold!14!white}

\pagecolor{popcream}

% ============================================================
% Typographie
% ============================================================
\defaultfontfeatures{Ligatures=TeX}
\setmainfont{PlusJakartaSans-Medium.ttf}[Path=../../../fonts/,BoldFont=PlusJakartaSans-Bold.ttf]
% Symboles, API et emojis absents de la police principale : repli sur DejaVu Sans (caractères actifs)
\newfontface\symfont{DejaVuSans.ttf}[Path=../../../fonts/]
\makeatletter
\newcommand\symactive[1]{\catcode#1=13 \begingroup\lccode`\~=#1 \lowercase{\endgroup\def~}{{\symfont\char#1}}}
\newcommand\symblank[1]{\catcode#1=13 \begingroup\lccode`\~=#1 \lowercase{\endgroup\def~}{}}
\newcommand\symhyph[1]{\catcode#1=13 \begingroup\lccode`\~=#1 \lowercase{\endgroup\def~}{\mbox{-}}}
\newcount\symc
\newcommand\symrange[2]{\symc=#1 \@whilenum\symc<#2 \do{\expandafter\symactive\expandafter{\the\symc}\advance\symc by 1 }}
\newcommand\symblankrange[2]{\symc=#1 \@whilenum\symc<#2 \do{\expandafter\symblank\expandafter{\the\symc}\advance\symc by 1 }}
\makeatother
\symrange{"0250}{"02FF}\symactive{"2713}\symactive{"2714}\symactive{"2717}\symactive{"2718}\symactive{"2610}\symactive{"2611}\symactive{"2612}
\symrange{"2600}{"27BF}
\catcode"2705=13 \begingroup\lccode`\~="2705 \lowercase{\endgroup\def~}{{\symfont\char"2713}}\symblank{"26BD}\symblank{"26A1}
\symrange{"2460}{"257F}\symactive{"223C}\symactive{"2248}\symactive{"2260}
\symactive{"01F8}\symactive{"01F9}\symactive{"1E3E}\symactive{"1E3F}
\symhyph{"2011}\symhyph{"2E1A}\symblankrange{"1F300}{"1FAFF}\symblank{"FE0F}
% Caractères chinois : WenQuanYi Zen Hei (sous-ensemble GB2312 + pinyin), gras/italique simulés
\setCJKmainfont{WenQuanYiZenHei-subset.ttf}[Path=../../../fonts/,AutoFakeBold=3,AutoFakeSlant=0.2]
\xeCJKsetup{CJKspace=false}
% Pinyin : voyelles à caron (ǎ ǐ ǒ ǔ ǖ ǘ ǚ ǜ) absentes de la police principale -> caractères actifs
\newfontface\pyfont{WenQuanYiZenHei-subset.ttf}[Path=../../../fonts/]
\newcommand\pyactive[1]{\catcode#1=13 \begingroup\lccode`\~=#1 \lowercase{\endgroup\def~}{{\pyfont\char#1}}}
\pyactive{"01CD}\pyactive{"01CE}\pyactive{"01CF}\pyactive{"01D0}\pyactive{"01D1}\pyactive{"01D2}\pyactive{"01D3}\pyactive{"01D4}\pyactive{"01D5}\pyactive{"01D6}\pyactive{"01D7}\pyactive{"01D8}\pyactive{"01D9}\pyactive{"01DA}\pyactive{"01DB}\pyactive{"01DC}
% Maths en sans-serif (Fira Math) avec chiffres et lettres droites de Plus
% Jakarta, pour que les formules se fondent dans le texte.
\usepackage[math-style=ISO,bold-style=ISO]{unicode-math}
\setmathfont{FiraMath-Regular.otf}[Path=../../../fonts/]
\setmathfont{PlusJakartaSans-Medium.ttf}[Path=../../../fonts/,range={up/num,up/{Latin,latin},\mathup}]
\usepackage{icomma} % virgule décimale sans espace en mode math
% Caractères mathématiques Unicode absents des polices de texte (Plus Jakarta,
% Archivo Black) : rendus par leur équivalent mathématique, en texte comme en
% formule — sinon ils s'affichent en carré vide (ex. « Arithmétique dans ℤ »).
\usepackage{newunicodechar}
\newunicodechar{ℝ}{\ensuremath{\mathbb{R}}}
\newunicodechar{ℤ}{\ensuremath{\mathbb{Z}}}
\newunicodechar{ℂ}{\ensuremath{\mathbb{C}}}
\newunicodechar{ℕ}{\ensuremath{\mathbb{N}}}
\newunicodechar{ℚ}{\ensuremath{\mathbb{Q}}}
\newunicodechar{𝔻}{\ensuremath{\mathbb{D}}}
\newunicodechar{α}{\ensuremath{\alpha}}
\newunicodechar{β}{\ensuremath{\beta}}
\newunicodechar{γ}{\ensuremath{\gamma}}
\newunicodechar{δ}{\ensuremath{\delta}}
\newunicodechar{ε}{\ensuremath{\varepsilon}}
\newunicodechar{θ}{\ensuremath{\theta}}
\newunicodechar{λ}{\ensuremath{\lambda}}
\newunicodechar{μ}{\ensuremath{\mu}}
\newunicodechar{σ}{\ensuremath{\sigma}}
\newunicodechar{τ}{\ensuremath{\tau}}
\newunicodechar{φ}{\ensuremath{\varphi}}
\newunicodechar{ψ}{\ensuremath{\psi}}
\newunicodechar{ω}{\ensuremath{\omega}}
\newunicodechar{Σ}{\ensuremath{\Sigma}}
% majuscules produites par \MakeUppercase dans les titres (α-aminés -> Α)
\newunicodechar{Α}{\ensuremath{\alpha}}
\newunicodechar{Β}{\ensuremath{\beta}}
\newunicodechar{Γ}{\ensuremath{\Gamma}}
\newunicodechar{Δ}{\ensuremath{\Delta}}
\newunicodechar{Ε}{\ensuremath{\varepsilon}}
\newunicodechar{Θ}{\ensuremath{\Theta}}
\newunicodechar{Λ}{\ensuremath{\Lambda}}
\newunicodechar{Μ}{\ensuremath{\mu}}
\newunicodechar{Τ}{\ensuremath{\tau}}
\newunicodechar{Φ}{\ensuremath{\Phi}}
\newunicodechar{Ψ}{\ensuremath{\Psi}}
\newunicodechar{Ω}{\ensuremath{\Omega}}
\newunicodechar{↦}{\ensuremath{\mapsto}}
\newunicodechar{∈}{\ensuremath{\in}}
\newunicodechar{∉}{\ensuremath{\notin}}
\newunicodechar{∧}{\ensuremath{\wedge}}
\newunicodechar{⇔}{\ensuremath{\Leftrightarrow}}
\newunicodechar{⟺}{\ensuremath{\Longleftrightarrow}}
\newunicodechar{⇒}{\ensuremath{\Rightarrow}}
\newunicodechar{⟹}{\ensuremath{\Longrightarrow}}
\newunicodechar{∘}{\ensuremath{\circ}}
\newunicodechar{∪}{\ensuremath{\cup}}
\newunicodechar{∩}{\ensuremath{\cap}}
\newunicodechar{≡}{\ensuremath{\equiv}}
\newunicodechar{⊂}{\ensuremath{\subset}}
\newunicodechar{⊥}{\ensuremath{\perp}}
\newunicodechar{∃}{\ensuremath{\exists}}
\newunicodechar{∀}{\ensuremath{\forall}}
\newunicodechar{∛}{\ensuremath{\sqrt[3]{\,}}}
\newunicodechar{∜}{\ensuremath{\sqrt[4]{\,}}}
\newunicodechar{⃗}{\ensuremath{{}^{\to}}}
\newunicodechar{ⁿ}{\ifmmode^{n}\else\textsuperscript{\ensuremath{n}}\fi}
\newunicodechar{ˣ}{\ifmmode^{x}\else\textsuperscript{\ensuremath{x}}\fi}
\newunicodechar{ᵇ}{\ifmmode^{b}\else\textsuperscript{\ensuremath{b}}\fi}
\newunicodechar{ʳ}{\ifmmode^{r}\else\textsuperscript{\ensuremath{r}}\fi}
\newunicodechar{ᵘ}{\ifmmode^{u}\else\textsuperscript{\ensuremath{u}}\fi}
\newunicodechar{ʸ}{\ifmmode^{y}\else\textsuperscript{\ensuremath{y}}\fi}
\newunicodechar{ᵃ}{\ifmmode^{a}\else\textsuperscript{\ensuremath{a}}\fi}
\newunicodechar{ᵉ}{\ifmmode^{e}\else\textsuperscript{\ensuremath{e}}\fi}
\newunicodechar{ᵏ}{\ifmmode^{k}\else\textsuperscript{\ensuremath{k}}\fi}
\newunicodechar{ᵖ}{\ifmmode^{p}\else\textsuperscript{\ensuremath{p}}\fi}
\newunicodechar{ᴺ}{\ifmmode^{N}\else\textsuperscript{\ensuremath{N}}\fi}
\newunicodechar{ᶜ}{\ifmmode^{c}\else\textsuperscript{\ensuremath{c}}\fi}
\newunicodechar{ʰ}{\ifmmode^{h}\else\textsuperscript{\ensuremath{h}}\fi}
\newunicodechar{ᵠ}{\ifmmode^{q}\else\textsuperscript{\ensuremath{q}}\fi}
\newunicodechar{ₙ}{\ifmmode_{n}\else\textsubscript{\ensuremath{n}}\fi}
\newunicodechar{ₐ}{\ifmmode_{a}\else\textsubscript{\ensuremath{a}}\fi}
\newunicodechar{ᵢ}{\ifmmode_{i}\else\textsubscript{\ensuremath{i}}\fi}
\newunicodechar{ᵣ}{\ifmmode_{r}\else\textsubscript{\ensuremath{r}}\fi}
\newunicodechar{ₖ}{\ifmmode_{k}\else\textsubscript{\ensuremath{k}}\fi}
\newunicodechar{ᵦ}{\ifmmode_{\beta}\else\textsubscript{\ensuremath{\beta}}\fi}
\newunicodechar{ₓ}{\ifmmode_{x}\else\textsubscript{\ensuremath{x}}\fi}
\newfontfamily\popdisplay{ArchivoBlack-Regular.ttf}[Path=../../../fonts/]
\newfontfamily\poplogo{SpaceGrotesk-Bold.ttf}[Path=../../../fonts/]
\newfontfamily\popsemi{PlusJakartaSans-SemiBold.ttf}[Path=../../../fonts/]
\newfontfamily\popxbold{PlusJakartaSans-ExtraBold.ttf}[Path=../../../fonts/]
\newfontfamily\popxboldls{PlusJakartaSans-ExtraBold.ttf}[Path=../../../fonts/,LetterSpace=3.0]

\setlist[enumerate]{leftmargin=1.7em,itemsep=5pt,topsep=4pt}
\setlist[itemize]{leftmargin=1.7em,itemsep=4pt,topsep=4pt,label={\color{poporange}\textbullet}}
\setlength{\parindent}{0pt}
\setlength{\parskip}{5pt}
\renewcommand{\arraystretch}{1.25}

% pgfplots : style Pop par défaut
\pgfplotsset{
  every axis/.append style={axis line style={popink,line width=1pt},tick style={popink},
    grid style={popline,line width=0.5pt},label style={font=\small},tick label style={font=\footnotesize}},
  popcurve/.style={poporange,line width=1.8pt,no markers,smooth},
}
% Même style disponible dans un \draw TikZ simple (hors pgfplots)
\tikzset{popcurve/.style={poporange,line width=1.8pt}}
\tikzset{popgrid/.style={popline,very thin}}
\colorlet{popcurve}{poporange} % le nom du style est parfois utilisé comme couleur

% ============================================================
% Pill / squiggle
% ============================================================
\newcommand{\poppill}[3]{%
  \tikz[baseline=-0.55ex]\node[draw=popink,line width=1.2pt,rounded corners=2.6pt,%
    fill=#2,inner xsep=6.5pt,inner ysep=3pt]{%
    \popxboldls\fontsize{7}{7}\selectfont\color{#3}\MakeUppercase{#1}};%
}
\newcommand{\popsquiggle}[1]{%
  \tikz[baseline=-0.4ex]{
    \draw[#1,line width=2.6pt,line cap=round]
      plot [smooth] coordinates {(0,0.09) (0.34,0.01) (0.68,0.21) (1.02,0.01) (1.36,0.10)};
  }%
}
% Mot-clé surligné (marqueur orange sous le texte)
\newcommand{\cle}[1]{{\popxbold\color{popdark}#1}}

% ============================================================
% En-tête / pied
% ============================================================
\pagestyle{fancy}
\fancyhf{}
\renewcommand{\headrulewidth}{0pt}
\renewcommand{\footrulewidth}{0pt}
\lhead{\raisebox{-0.15ex}{\poplogo\fontsize{13}{13}\selectfont\color{popink}OnBuch\color{poporange}+}}
\rhead{\poppill{\DOCMATIERE\ \textbullet\ \DOCNIVEAU\ \textbullet\ Leçon \DOCLECON}{white}{popink}}
\lfoot{\popxbold\fontsize{7.6}{7.6}\selectfont\color{popmuted}\MakeUppercase{Cours signé OnBuch+}\ \textbullet\ {\popsemi\DOCTITRE}}
\rfoot{\tikz[baseline=-0.6ex]\node[rounded corners=8.5pt,fill=popink,minimum height=17pt,minimum width=17pt,inner xsep=5pt,inner sep=0]{\popxbold\fontsize{8}{8}\selectfont\color{popcream}\thepage\,/\,\pageref*{LastPage}};}

% ============================================================
% Titres
% ============================================================
\newcommand{\chaptitre}[2]{%
  \begin{tcolorbox}[enhanced,colback=poporange,colframe=popink,boxrule=2.6pt,arc=11pt,
      drop shadow=popink,left=15pt,right=15pt,top=12pt,bottom=11pt,before skip=3pt,after skip=18pt]
    \poppill{#2}{popink}{popcream}\par\vspace{9pt}
    {\raggedright\hyphenpenalty=10000\exhyphenpenalty=10000\popdisplay\fontsize{22}{24}\selectfont\color{popcream}\MakeUppercase{#1}\par}
  \end{tcolorbox}
}
% Section numérotée : pastille orange avec le numéro + titre Archivo
\newcounter{popsec}
\newcommand{\coursec}[1]{%
  \stepcounter{popsec}\setcounter{popsub}{0}%
  \par\vspace{16pt}\noindent
  \begin{minipage}{\linewidth}
  \tikz[baseline=-0.7ex]\node[draw=popink,line width=1.6pt,fill=poporange,rounded corners=4pt,minimum size=22pt,inner sep=0]{\popdisplay\fontsize{12}{12}\selectfont\color{popcream}\thepopsec};%
  \hspace{8pt}{\popdisplay\fontsize{15}{17}\selectfont\color{popink}\MakeUppercase{#1}}\par
  \vspace{4pt}\noindent{\color{poporange}\rule{36pt}{3.6pt}}
  \end{minipage}\par\nopagebreak\vspace{8pt}
}
\newcounter{popsub}
\newcommand{\courssub}[1]{%
  \stepcounter{popsub}%
  \par\vspace{9pt}\noindent{\popxbold\fontsize{11.5}{13}\selectfont\color{popink}{\color{poporange}\thepopsec.\thepopsub}\hspace{6pt}#1}\par\nopagebreak\vspace{4pt}
}

% ============================================================
% Boîtes pédagogiques — même recette : fond blanc, bordure épaisse colorée,
% ombre dure encre, badge plein. Titre optionnel [#1].
% ============================================================
\tcbset{popbase/.style={breakable,enhanced,colback=white,boxrule=2.3pt,arc=9pt,
  shadow={3pt}{-3pt}{0pt}{popink},left=14pt,right=14pt,top=11pt,bottom=11pt,before skip=11pt,after skip=15pt,
  fontupper=\popsemi\fontsize{10.6}{14.5}\selectfont\color{popink},
  fontlower=\popsemi\fontsize{10.6}{14.5}\selectfont\color{popink}}}
% Coche dessinée (le glyphe ✓ n'existe pas dans les polices du document)
\newcommand{\popcheck}[1]{\tikz[baseline=-0.5ex]\draw[#1,line width=1.6pt,line cap=round,line join=round](-0.13,0.0)--(-0.04,-0.09)--(0.13,0.1);}
\providecommand{\coche}{\popcheck{popgreen}}
\providecommand{\resultat}[1]{\par\smallskip\noindent\fcolorbox{popgreen}{popgreenL}{\parbox{\dimexpr\linewidth-2\fboxsep-2\fboxrule\relax}{\textbf{Résultat.} #1}}\par\smallskip}
\newcommand{\popboxhead}[3]{\poppill{#1}{#2}{white}\ifx\relax#3\relax\else\hspace{6pt}{\popxbold\fontsize{10.6}{12}\selectfont\color{popink}#3}\fi\par\vspace{7pt}}

\newtcolorbox{aretenir}[1][]{popbase,colframe=popgreen,before upper={\popboxhead{À retenir}{popgreen}{#1}}}
% Texte en chinois (document de lecture, dialogue, extrait) : cadre
% d'étude. Un titre optionnel (source, auteur) s'affiche dans l'en-tête.
\newtcolorbox{textebox}[1][]{popbase,colframe=popblue,before upper={\popboxhead{文本 · Texte}{popblue}{#1}},after upper={}}
% Marques ✓ / ✗ (forme correcte / fautive) souvent utilisées par les modèles
\providecommand{\cross}{\textcolor{poppink}{\ensuremath{\times}}}
\providecommand{\xmark}{\cross}\providecommand{\crossmark}{\cross}\providecommand{\cmark}{\ensuremath{\checkmark}}
% Ligne de réponse / justification à compléter (inventée par les modèles)
\providecommand{\justification}[1]{\par\noindent\textit{Justification~:} \dotfill\par}
% \textipa (package tipa non chargé) : la police ne couvre pas l'API -> simple passage
\providecommand{\textipa}[1]{#1}
% Soulignement (ulem non chargé) : \uline, \ul
\providecommand{\uline}[1]{\underline{#1}}\providecommand{\ul}[1]{\underline{#1}}
% \glossaire{\item ...} inventée par les modèles : liste de glossaire
\providecommand{\glossaire}[1]{\begin{itemize}#1\end{itemize}}
% \alert (beamer) inventée par les modèles : texte mis en évidence
\providecommand{\alert}[1]{\textbf{\textcolor{poppink}{#1}}}
% variantes ASCII d'ordinaux écrites en macro
\providecommand{\iere}{\textsuperscript{re}}\providecommand{\ieme}{\textsuperscript{e}}\providecommand{\ere}{\textsuperscript{re}}\providecommand{\eme}{\textsuperscript{e}}\providecommand{\er}{\textsuperscript{er}}\providecommand{\ie}{\textsuperscript{e}}
% Ordinaux français écrits en macro par les modèles : 1\ière, 3\ième
\newcommand{\ière}{\textsuperscript{re}}\newcommand{\ième}{\textsuperscript{e}}
% \exemple (inventée par les modèles) : étiquette « Exemple : »
\providecommand{\exemple}{\textbf{Exemple~:}\ }
% \square utilisable aussi hors mode math (cases à cocher)
\AtBeginDocument{\let\squareMath\square\renewcommand{\square}{\ensuremath{\squareMath}}}
% Mot ou phrase en chinois à l'intérieur d'un paragraphe français
\newcommand{\zh}[1]{\ifmmode\text{#1}\else{#1}\fi}
\let\eng\zh \let\esp\zh \let\all\zh % alias tolérés
% flèches d'intonation inventées par les modèles
\providecommand{\textsearrow}{}\renewcommand{\textsearrow}{\ensuremath{\searrow}}\providecommand{\textnearrow}{}\renewcommand{\textnearrow}{\ensuremath{\nearrow}}
% \fr{..} : mot français (inventé par les modèles)
\providecommand{\fr}[1]{\textit{#1}}
% zhbox : encadré de texte chinois inventé par les modèles
\newenvironment{zhbox}{\begin{quote}}{\end{quote}}
% \courssubsub : titre de niveau 3 inventé par les modèles
\providecommand{\courssubsub}[1]{\par\medskip\noindent\textbf{#1}\par\smallskip}
% \zhbf : texte chinois en gras inventé par les modèles
\providecommand{\zhbf}[1]{\textbf{#1}}
% \vs : « versus » inventé par les modèles
\providecommand{\vs}{\textit{vs}\ }
% \blank : case à compléter inventée par les modèles
\providecommand{\blank}[1][]{\underline{\hspace{1.6em}}}
\newtcolorbox{exemplebox}[1][]{popbase,colframe=poporange,before upper={\popboxhead{Exemple}{poporange}{#1}}}
\newtcolorbox{definition}[1][]{popbase,colframe=popblue,before upper={\popboxhead{Définition}{popblue}{#1}}}
\newtcolorbox{propriete}[1][]{popbase,colframe=poppurple,before upper={\popboxhead{Propriété}{poppurple}{#1}}}
\newtcolorbox{methode}[1][]{popbase,colframe=popgold,before upper={\popboxhead{Méthode}{popgold}{#1}}}
\newtcolorbox{attention}[1][]{popbase,colframe=poppink,before upper={\popboxhead{Attention}{poppink}{#1}}}
\newtcolorbox{experience}[1][]{popbase,colframe=popblue,colback=popblueL,before upper={\popboxhead{Expérience}{popblue}{#1}}}
% « Le savais-tu ? » : carte violette pleine (contexte, culture, Cameroun)
\newtcolorbox{savaistu}[1][]{popbase,colback=poppurple,colframe=popink,
  fontupper=\popsemi\fontsize{10.4}{14.2}\selectfont\color{white},
  before upper={\poppill{Le savais-tu\,?}{white}{poppurple}\ifx\relax#1\relax\else\hspace{6pt}{\popxbold\color{white}#1}\fi\par\vspace{7pt}}}
% Exercice résolu : énoncé en haut, \tcblower puis la solution sur fond crème
\newcounter{popexo}\newcounter{popexr}
\newtcolorbox{exoresolu}[1][]{popbase,colframe=poporange,colbacklower=popcream,
  segmentation style={popink,line width=1.2pt,dashed},
  before upper={\stepcounter{popexr}\popboxhead{Exercice résolu \thepopexr}{poporange}{#1}},
  before lower={\poppill{Solution}{popink}{popcream}\par\vspace{6pt}}}
% Exercice (non résolu en place) — la correction est dans la partie Corrigés
\newtcolorbox{exercice}[1][]{popbase,colframe=popink,
  before upper={\stepcounter{popexo}\popboxhead{Exercice \thepopexo}{popink}{#1}}}
% Corrigé
\newtcolorbox{corrige}[1][]{popbase,colframe=popgreen,colback=popgreenL,
  before upper={\popboxhead{Corrigé}{popgreen}{#1}}}
% Objectifs / prérequis / bilan
\newtcolorbox{objectifs}{popbase,colframe=popink,colback=poporangeL,
  before upper={\popboxhead{Objectifs de la leçon}{popink}{}}}
\newtcolorbox{prerequis}{popbase,colframe=popink,colback=popblueL,
  before upper={\popboxhead{Prérequis}{popblue}{}}}
\newtcolorbox{bilan}{popbase,colframe=popink,colback=white,boxrule=2.8pt,
  before upper={\popboxhead{Fiche bilan}{poporange}{L'essentiel à connaître pour l'examen}}}
% Figure encadrée avec légende
\newtcolorbox{popfigure}[1][]{enhanced,breakable=false,colback=white,colframe=popline,boxrule=1.4pt,arc=8pt,
  left=10pt,right=10pt,top=10pt,bottom=8pt,before skip=10pt,after skip=12pt,halign=center,
  lower separated=false,halign lower=center,
  fontlower=\popsemi\fontsize{9}{11.5}\selectfont\color{popmuted},
  #1}
\newcommand{\legende}[1]{\par\vspace{6pt}{\popsemi\fontsize{9}{11.5}\selectfont\color{popmuted}#1}}

% ============================================================
% Signature OnBuch+
% ============================================================
\newcommand{\poplogoinline}[1]{{\poplogo\fontsize{#1}{#1}\selectfont\color{popink}OnBuch\color{poporange}+}}
\newcommand{\signatureonbuch}{%
  \begin{tcolorbox}[enhanced,colback=popink,colframe=popink,boxrule=2.2pt,arc=10pt,
      drop shadow=poporange,left=14pt,right=14pt,top=10pt,bottom=10pt,before skip=0pt,after skip=0pt]
    \tikz[baseline=-0.7ex]\node[circle,fill=popgreen,draw=popcream,line width=1.3pt,minimum size=22pt,inner sep=0]{\popcheck{white}};%
    \hspace{9pt}\begin{minipage}[c]{0.62\linewidth}
      {\popxbold\fontsize{12}{13}\selectfont\color{popcream}Signé\ \textbullet\ {\poplogo OnBuch\color{poporange}+}}\par\vspace{2pt}
      {\raggedright\popsemi\fontsize{8}{10.5}\selectfont\color{popline}\MakeUppercase{Équipe pédagogique OnBuch+}\\\MakeUppercase{Conforme au programme officiel MINESEC}\par}
    \end{minipage}\hfill
    {\poplogo\fontsize{22}{22}\selectfont\color{popcream}OnBuch\color{poporange}+}
  \end{tcolorbox}%
}

% ============================================================
% Page de garde
% #1 titre · #2 matière · #3 niveau · #4 module · #5 n° leçon · #6 durée ·
% #7 description / objectifs (texte ou itemize)
% ============================================================
\newcommand{\pagegarde}[7]{%
  \thispagestyle{empty}
  \newgeometry{a4paper,margin=1.7cm,top=1.6cm,bottom=1.4cm}%
  \begin{tikzpicture}[remember picture,overlay]
    \fill[poporange!18] ([xshift=-3.2cm,yshift=-3.6cm]current page.north east) circle (4.6cm);
    \fill[poppurple!14] ([xshift=2.4cm,yshift=7.6cm]current page.south west) circle (3.8cm);
  \end{tikzpicture}%
  {\poplogo\fontsize{30}{30}\selectfont\color{popink}OnBuch\color{poporange}+}\hfill
  \raisebox{0.6ex}{\poppill{Cours signé \textbullet\ Premium}{popink}{popcream}}\par
  \vspace{1pt}\popsquiggle{poppurple}\par
  \vspace{8pt}
  \begin{tcolorbox}[enhanced,colback=poporange,colframe=popink,boxrule=2.8pt,arc=13pt,
      drop shadow=popink,left=18pt,right=18pt,top=12pt,bottom=13pt,after skip=10pt]
    \poppill{#2 \textbullet\ #3}{popink}{popcream}\hspace{5pt}\poppill{#4}{white}{popink}\par\vspace{8pt}
    {\popdisplay\fontsize{14}{14}\selectfont\color{popink}LEÇON #5}\par\vspace{4pt}
    {\raggedright\hyphenpenalty=10000\exhyphenpenalty=10000\popdisplay\fontsize{25}{27}\selectfont\color{popcream}\MakeUppercase{#1}\par}
    \vspace{8pt}
    {\popxbold\fontsize{9.5}{9.5}\selectfont\color{popink}\MakeUppercase{Durée conseillée : #6}}
  \end{tcolorbox}
  \begin{tcolorbox}[enhanced,colback=white,colframe=popink,boxrule=2.2pt,arc=10pt,
      drop shadow=popink,left=16pt,right=16pt,top=10pt,bottom=10pt,after skip=8pt]
    \poppill{Dans cette leçon}{poppurple}{white}\par\vspace{5pt}
    \popsemi\fontsize{9.3}{12.6}\selectfont\color{popink}#7
  \end{tcolorbox}
  \noindent\begin{minipage}[t]{0.487\linewidth}
    \begin{tcolorbox}[enhanced,colback=popgreen,colframe=popink,boxrule=2.2pt,arc=10pt,
        drop shadow=popink,left=11pt,right=11pt,top=10pt,bottom=10pt,height=3.1cm,valign=center]
      \tcbox[colback=white,colframe=popink,boxrule=1.3pt,arc=5pt,boxsep=0pt,nobeforeafter,
        left=3pt,right=3pt,top=3pt,bottom=3pt,valign=center]{\qrcode[height=1.9cm]{https://play.google.com/store/apps/details?id=com.luvvix.onbuch&hl=fr}}%
      \hspace{8pt}\begin{minipage}[c]{0.5\linewidth}
        {\popdisplay\fontsize{11}{12.5}\selectfont\color{white}\MakeUppercase{Télécharge OnBuch}}\par\vspace{4pt}
        {\raggedright\popsemi\fontsize{8.4}{11}\selectfont\color{white}Gratuit \textbullet\ Google Play\\Tuteur IA, annales, concours, résultats\par}
      \end{minipage}
    \end{tcolorbox}
  \end{minipage}\hfill
  \begin{minipage}[t]{0.487\linewidth}
    \begin{tcolorbox}[enhanced,colback=poppurple,colframe=popink,boxrule=2.2pt,arc=10pt,
        drop shadow=popink,left=11pt,right=11pt,top=10pt,bottom=10pt,height=3.1cm,valign=center]
      \tikz[baseline=-0.7ex]\node[circle,fill=white,draw=popink,line width=1.3pt,minimum size=1.6cm,inner sep=0]{\popdisplay\fontsize{24}{24}\selectfont\color{poppurple}+};%
      \hspace{8pt}\begin{minipage}[c]{0.6\linewidth}
        {\popdisplay\fontsize{11}{12.5}\selectfont\color{white}\MakeUppercase{Passe à OnBuch+}}\par\vspace{4pt}
        {\raggedright\popsemi\fontsize{8.4}{11}\selectfont\color{white}Tous les cours signés, Tuteur IA illimité, fiches PDF — onglet OnBuch+ de l'app\par}
      \end{minipage}
    \end{tcolorbox}
  \end{minipage}
  \vspace{8pt}
  \popcontenu
  \vfill
  \signatureonbuch
  \restoregeometry
  \newpage
}

% Rangée « Ce cours contient » (page de garde)
\newcommand{\poptile}[3]{% #1 couleur #2 gros texte #3 légende
  \begin{tcolorbox}[enhanced,colback=#1,colframe=popink,boxrule=2pt,arc=9pt,shadow={2.5pt}{-2.5pt}{0pt}{popink},
      left=6pt,right=6pt,top=9pt,bottom=9pt,halign=center,valign=center,height=1.75cm,width=\linewidth,nobeforeafter]
    {\popdisplay\fontsize{12.5}{13}\selectfont\color{white}\MakeUppercase{#2}}\par\vspace{3pt}
    {\centering\popsemi\fontsize{7.8}{9.5}\selectfont\color{white}#3\par}
  \end{tcolorbox}}
\newcommand{\popcontenu}{%
  \par\noindent{\popxboldls\fontsize{7.5}{7.5}\selectfont\color{popmuted}CE COURS SIGNÉ CONTIENT}\par\vspace{7pt}
  \noindent\begin{minipage}[t]{0.235\linewidth}\poptile{popblue}{Cours}{complet, illustré, pas à pas}\end{minipage}\hfill
  \begin{minipage}[t]{0.235\linewidth}\poptile{poporange}{Méthodes}{et exercices résolus}\end{minipage}\hfill
  \begin{minipage}[t]{0.235\linewidth}\poptile{poppink}{Exercices}{type examen + corrigés}\end{minipage}\hfill
  \begin{minipage}[t]{0.235\linewidth}\poptile{popgreen}{Fiche bilan}{l'essentiel à retenir}\end{minipage}\par}

% Sommaire
\newtcolorbox{sommaire}{enhanced,colback=white,colframe=popink,boxrule=2.2pt,arc=10pt,
  drop shadow=popink,left=18pt,right=18pt,top=13pt,bottom=13pt,before skip=6pt,after skip=20pt,
  before upper={\poppill{Sommaire}{poppurple}{white}\par\vspace{9pt}\popsemi\fontsize{10.6}{14.5}\selectfont\color{popink}}}

% Signature de fin de cours — poussée tout en bas de la dernière page
\newcommand{\findecours}{%
  \par\vfill\nopagebreak
  \begin{center}\popsquiggle{poporange}\end{center}\vspace{4pt}
  \signatureonbuch
}
\AtBeginDocument{\def\llbracket{\mathopen{[\mkern-3.5mu[}}\def\rrbracket{\mathclose{]\mkern-3.5mu]}}} % intervalles d'entiers (glyphe ⟦ absent de la police)
% Glyphes absents de Fira Math : redéfinis à partir de glyphes disponibles
\AtBeginDocument{%
  \def\setminus{\mathbin{\backslash}}\let\smallsetminus\setminus
  \def\vdots{\mathord{\vbox{\baselineskip3.2pt\lineskiplimit0pt\kern4pt\hbox{.}\hbox{.}\hbox{.}}}}%
  \def\ddots{\mathinner{\mkern1mu\raise6.5pt\hbox{.}\mkern2mu\raise3.6pt\hbox{.}\mkern2mu\raise0.7pt\hbox{.}\mkern1mu}}%
  \def\triangle{\mathord{\scalebox{0.9}{$\symup{\Delta}$}}}%
  \def\nabla{\mathord{\raisebox{\depth}{\scalebox{1}[-1]{$\symup{\Delta}$}}}}%
  \def\checkmark{\popcheck{popdark}}%
  \def\star{\mathbin{\ast}}\def\bigstar{\mathord{\ast}}%
  \def\diamond{\mathbin{\scalebox{0.6}{$\Diamond$}}}%
  \def\bigcup{\mathop{\mathchoice{\raisebox{-0.3em}{\scalebox{1.7}{$\cup$}}}{\scalebox{1.25}{$\cup$}}{\cup}{\cup}}\displaylimits}%
  \def\bigcap{\mathop{\mathchoice{\raisebox{-0.3em}{\scalebox{1.7}{$\cap$}}}{\scalebox{1.25}{$\cap$}}{\cap}{\cap}}\displaylimits}%
  \def\lesseqqgtr{\mathrel{\lessgtr}}\def\bigoplus{\mathop{\oplus}\displaylimits}%
}
\providecommand{\ssi}{si et seulement si} % abréviation fréquente
\AtBeginDocument{\providecommand{\wideparen}{\overparen}} % arc de cercle

%% FIGFIX-BEGIN v4 — correctifs globaux des figures (docs/onbuch-generation/tools/figfix)
\usepackage{adjustbox}\usepackage{etoolbox}
% toute figure TikZ/pgfplots plus large que la ligne est réduite à la largeur disponible
\BeforeBeginEnvironment{tikzpicture}{\begin{adjustbox}{max width=\linewidth}}
\AfterEndEnvironment{tikzpicture}{\end{adjustbox}}
% légendes pgfplots lisibles par-dessus les courbes
\pgfplotsset{every axis/.append style={legend style={fill=white,fill opacity=0.92,text opacity=1,draw=black!15,inner sep=3pt}}}
% étiquettes d'arêtes (syntaxe edge["..."]) avec fond blanc
\tikzset{every edge quotes/.append style={fill=white,inner sep=1.2pt}}
%% FIGFIX-END
\begin{document}

\pagegarde{Leçon 34 — Vocabulaire et compréhension de texte : internet et les réseaux sociaux}{Chinois}{Tle A}{Module 5 : Mass médias et télécommunication}{ZH34}{2 h}{Cette leçon de vocabulaire et compréhension de texte porte sur internet et les réseaux sociaux en chinois. À travers un texte support original, un glossaire structuré, l'étude des radicaux et des questions de compréhension, les élèves apprennent à lire, prononcer et employer le lexique du numérique.
\vspace{4pt}
\begin{multicols}{2}\small
\begin{itemize}
\item À la fin de cette leçon, je sais lire à voix haute un court texte chinois sur internet et les réseaux sociaux en respectant le pinyin et les tons.
\item Je sais reconnaître et traduire au moins 20 mots du thème (internet, réseau, téléphone portable, publier, suivre, etc.).
\item Je sais identifier les radicaux (clés) des caractères du thème et les utiliser pour deviner le sens de caractères nouveaux.
\item Je sais écrire correctement quelques caractères clés en respectant l'ordre des traits.
\item Je sais répondre à des questions de compréhension portant sur le texte support.
\item Je sais employer les collocations usuelles du thème (上网, 发朋友圈, 加好友...) dans des phrases complètes.
\end{itemize}\end{multicols}}

\begin{sommaire}
\begin{multicols}{2}
\begin{enumerate}
\item Mise en route et découverte du thème
\item Texte support : lecture et compréhension globale
\item Glossaire du thème
\item Radicaux et ordre des traits
\item Questions de compréhension du texte
\item Production guidée et comparaison Cameroun-Chine
\item Activité d'intégration
\item Méthodes et exercices résolus
\item Exercices
\item Corrigés des exercices
\item Fiche bilan
\end{enumerate}
\end{multicols}
\end{sommaire}

\begin{objectifs}
\begin{itemize}
\item À la fin de cette leçon, je sais lire à voix haute un court texte chinois sur internet et les réseaux sociaux en respectant le pinyin et les tons.
\item Je sais reconnaître et traduire au moins 20 mots du thème (internet, réseau, téléphone portable, publier, suivre, etc.).
\item Je sais identifier les radicaux (clés) des caractères du thème et les utiliser pour deviner le sens de caractères nouveaux.
\item Je sais écrire correctement quelques caractères clés en respectant l'ordre des traits.
\item Je sais répondre à des questions de compréhension portant sur le texte support.
\item Je sais employer les collocations usuelles du thème (上网, 发朋友圈, 加好友...) dans des phrases complètes.
\item Je sais produire un court paragraphe guidé sur mon usage d'internet et des réseaux sociaux.
\item Je sais comparer brièvement les usages numériques des jeunes camerounais et chinois.
\end{itemize}
\end{objectifs}

\begin{prerequis}
\begin{itemize}
\item Vocabulaire des médias et de la communication vu en Première (télévision, radio, téléphone)
\item Structures de base : phrase sujet-verbe-objet, 是, 有, 在
\item Adverbes de fréquence (常常, 每天, 有时候) et compléments de durée
\item Notions de radicaux et d'ordre des traits acquises depuis la Seconde
\item Verbes modaux 会, 可以, 应该
\end{itemize}
\end{prerequis}

\newpage

\coursec{Mise en route et découverte du thème}

\courssub{Situation d'accroche : un matin à Yaoundé, un matin à Pékin}

À Yaoundé comme à Douala, Bamenda ou Buea, presque tous les élèves de Terminale consultent WhatsApp, Facebook ou TikTok avant même d'arriver en classe. Dans le car qui descend vers le marché Mokolo, sur les bancs du lycée, pendant la récréation : les pouces glissent sur les écrans, les messages s'envoient, les vidéos défilent. En Chine aussi, les jeunes passent des heures sur WeChat (\zh{微信}) et Douyin (\zh{抖音}), l'équivalent chinois de TikTok. Les deux jeunesses se ressemblent... mais pas leurs applications, et surtout pas leurs mots !

Mais sauriez-vous parler \emph{en chinois} de votre usage d'internet et des réseaux sociaux ? Comment dire « se connecter », « publier un message » ou « suivre un ami » ? C'est toute la \cle{problématique} de cette leçon : acquérir le vocabulaire chinois du monde numérique et apprendre à comprendre un texte authentique sur ce thème.

\begin{textebox}[Deux questions pour commencer]
你每天都上网吗？你用什么社交软件？\\
\emph{Nǐ měi tiān dōu shàngwǎng ma? Nǐ yòng shénme shèjiāo ruǎnjiàn?}\\
« Est-ce que tu te connectes à internet tous les jours ? Quelles applications sociales utilises-tu ? »
\end{textebox}

Ces deux questions, vous pourrez y répondre en chinois à la fin de la leçon. Gardons-les en tête : elles seront notre fil conducteur.

\courssub{Brainstorming en français : votre vie numérique}

Avant d'ouvrir le chinois, parlons de vous. Répondez en français, à voix haute ou par écrit :

\begin{enumerate}
\item Quels réseaux sociaux ou quelles applications utilisez-vous chaque jour ? (WhatsApp, Facebook, TikTok, Instagram, YouTube...)
\item Combien d'heures par jour passez-vous sur votre téléphone portable ?
\item Pour quoi faire : discuter avec des amis, regarder des vidéos, réviser vos leçons, suivre l'actualité, le football ?
\item Avez-vous des « amis » que vous n'avez jamais rencontrés en vrai ?
\end{enumerate}

L'enseignant note au tableau les mots français qui reviennent : \emph{internet, téléphone portable, se connecter, message, ami en ligne, publier, application}. Chacun de ces mots a un équivalent chinois que nous allons découvrir maintenant. C'est la démarche du bon sinisant : partir de ce que l'on vit pour atteindre ce que l'on doit dire.

\courssub{Découverte des mots-clés du thème}

Voici les sept mots-clés autour desquels toute la leçon est construite. Observez-les d'abord, écoutez l'enseignant les prononcer, puis répétez en faisant attention aux tons.

\begin{center}
\begin{tabularx}{\linewidth}{|X|X|X|X|}
\hline
\rowcolor{popblueL} Caractères & Pinyin & Nature & Traduction \\
\hline
网络 & wǎngluò & nom & internet, le réseau \\
\hline
上网 & shàngwǎng & verbe & se connecter (à internet), aller en ligne \\
\hline
手机 & shǒujī & nom & téléphone portable \\
\hline
微信 & Wēixìn & nom propre & WeChat (application chinoise) \\
\hline
朋友圈 & péngyouquān & nom & « cercle d'amis », fil d'actualité (Moments) \\
\hline
网友 & wǎngyǒu & nom & ami(e) connu(e) sur internet \\
\hline
社交软件 & shèjiāo ruǎnjiàn & nom & application sociale, réseau social \\
\hline
\end{tabularx}
\end{center}

\begin{definition}[Observer les caractères pour mémoriser]
Chaque mot-clé est construit de façon logique ; comprendre sa construction, c'est le mémoriser deux fois plus vite :
\begin{itemize}
\item \zh{网} (wǎng, « filet, réseau ») est la clé de voûte du thème : on le retrouve dans \zh{网络} (réseau), \zh{上网} (se connecter, litt. « monter sur le réseau ») et \zh{网友} (ami du réseau).
\item \zh{手机} = \zh{手} (shǒu, « main ») + \zh{机} (jī, « machine ») : littéralement « la machine de la main ». Image parfaite du téléphone portable !
\item \zh{朋友圈} = \zh{朋友} (péngyou, « ami ») + \zh{圈} (quān, « cercle ») : le « cercle des amis », c'est le fil d'actualité de WeChat où l'on publie photos et messages.
\item \zh{网友} = \zh{网} (réseau) + \zh{友} (yǒu, « ami ») : un ami rencontré en ligne.
\item \zh{社交软件} = \zh{社交} (shèjiāo, « social, relations sociales ») + \zh{软件} (ruǎnjiàn, « logiciel ») : « logiciel social ».
\end{itemize}
\end{definition}

\begin{popfigure}
\begin{tikzpicture}[mindmap, grow cyclic, every node/.style={concept, align=center, font=\small},
root concept/.append style={concept color=popblue, font=\large},
level 1/.append style={level distance=3.4cm, sibling angle=60},
level 1 concept/.append style={concept color=poporangeL}]
\node[root concept, align=center]{网络\\ \emph{wǎngluò}\\ internet}
  child { node[align=center]{手机\\ \emph{shǒujī}\\ téléphone\\ portable} }
  child { node[align=center]{微信\\ \emph{Wēixìn}\\ WeChat} }
  child { node[align=center]{网友\\ \emph{wǎngyǒu}\\ ami en ligne} }
  child { node[align=center]{上网\\ \emph{shàngwǎng}\\ se connecter} }
  child { node[align=center]{朋友圈\\ \emph{péngyouquān}\\ fil d'actualité} }
  child { node[align=center]{社交软件\\ \emph{shèjiāo ruǎnjiàn}\\ appli sociale} };
\end{tikzpicture}
\legende{Carte mentale des mots-clés : tout le vocabulaire du thème tourne autour de \zh{网} (le réseau).}
\end{popfigure}

\begin{attention}[Piège de prononciation : les tons de \zh{网} et de \zh{微信}]
\zh{网} (wǎng) porte le \cle{3e ton}, descendant puis montant. Les francophones le prononcent souvent à plat, comme un 1er ton : on entend alors \zh{汪} (wāng, « ouaf ! », l'aboiement du chien). De même, \zh{微信} (Wēixìn) = 1er ton + 4e ton : « wēi » monte haut et reste haut, « xìn » tombe brusquement. Dire « wěixìn » ou « wéixín » rend le mot méconnaissable pour un Chinois. Entraînez-vous : \emph{Wēi-xìn, Wēi-xìn}.
\end{attention}

\courssub{Radicaux et ordre des traits : les clés du thème}

Connaître le radical (clé) et l'ordre des traits permet de chercher un caractère dans le dictionnaire et de l'écrire correctement. Voici l'analyse pour cinq caractères centraux.

\begin{center}
\begin{tabularx}{\linewidth}{|X|X|X|X|}
\hline
\rowcolor{popblueL} Caractère & Radical (clé) & Nb traits & Ordre des traits (résumé) \\
\hline
网 (wǎng) & \zh{网} (radical 122, « filet ») & 6 & Horizontal, vertical, horizontal, vertical, horizontal, horizontal (le cadre se referme) \\
\hline
手 (shǒu) & \zh{手} (radical 64, « main ») & 4 & Horizontal, crochet vertical, horizontal, trait oblique droit \\
\hline
微 (wēi) & \zh{彳} (radical 60, « pas ») & 13 & Clé \zh{彳} (gauche), puis \zh{徂} (droite : haut → bas) \\
\hline
朋 (péng) & \zh{月} (radical 74, « chair/lune ») & 8 & Deux \zh{月} superposés : haut → bas \\
\hline
友 (yǒu) & \zh{又} (radical 29, « encore/main droite ») & 7 & \zh{上} (haut) puis \zh{又} (bas) \\
\hline
\end{tabularx}
\end{center}

\begin{aretenir}[Règle d'or de l'ordre des traits]
Respectez toujours : \textbf{haut → bas}, \textbf{gauche → droite}, \textbf{horizontal avant vertical}, \textbf{extérieur avant intérieur}, \textbf{fermer en dernier}. Pour \zh{网}, on trace le pourtour avant les traits intérieurs. Pour \zh{友}, on écrit \zh{上} complètement avant \zh{又}.
\end{aretenir}

\courssub{Première écoute du texte support : saisie globale}

Fermez vos cahiers : pas de support écrit. L'enseignant va lire à voix haute, deux fois, un court texte intitulé \zh{我和我的手机} (\emph{Wǒ hé wǒ de shǒujī}, « Mon téléphone portable et moi »). Votre tâche à cette première écoute n'est \emph{pas} de tout comprendre : c'est de saisir l'\cle{idée générale}.

\begin{methode}[Comment aborder un texte chinois en quatre étapes]
\begin{enumerate}
\item \textbf{Lire le titre} : il annonce le thème. Ici, \zh{我和我的手机} nous dit que le narrateur va parler de son rapport à son téléphone.
\item \textbf{Repérer les mots connus} : pendant l'écoute, tendez l'oreille vers les mots-clés de la carte mentale (\zh{网络}, \zh{手机}, \zh{微信}, \zh{朋友圈}, \zh{网友}, \zh{上网}) et les mots déjà appris (\zh{我}, \zh{每天}, \zh{朋友}, \zh{学习}, \zh{时间}).
\item \textbf{Écouter une première fois sans écrire} : notez mentalement qui parle, de quoi, et si le ton est positif ou critique.
\item \textbf{Lire ensuite en détail avec le glossaire} : la compréhension fine viendra à la deuxième étape, texte sous les yeux (section suivante).
\end{enumerate}
\end{methode}

\begin{experience}[Écoute active : les trois questions de saisie globale]
Pendant la lecture par l'enseignant, chaque élève tient mentalement les réponses à trois questions :
\begin{enumerate}
\item De quoi parle le narrateur : de ses études, de sa famille ou de son téléphone ?
\item Quelles applications sont citées ?
\item Le narrateur est-il plutôt content ou plutôt inquiet de son usage d'internet ?
\end{enumerate}
Après la deuxième écoute, comparez vos réponses avec votre voisin, puis mise en commun orale en français. L'enseignant valide sans encore dévoiler le texte écrit.
\end{experience}

\courssub{Hypothèses sur le contenu : prédire avant de lire}

À partir du titre \zh{我和我的手机} et des sept mots-clés, formulez des \cle{hypothèses} sur le contenu du texte. Prédire est une stratégie de lecteur expert : votre cerveau, ayant anticipé, comprend ensuite beaucoup plus vite.

\begin{exemplebox}[Exemples d'hypothèses formulées en classe]
\begin{itemize}
\item « Le narrateur va probablement dire qu'il utilise son téléphone tous les jours. » (mot-clé : \zh{手机}, \zh{每天})
\item « Il va peut-être parler de WeChat et de son fil d'actualité. » (mots-clés : \zh{微信}, \zh{朋友圈})
\item « Il a peut-être des amis en ligne qu'il n'a jamais rencontrés. » (mot-clé : \zh{网友})
\item « Le texte finira peut-être par un conseil : ne pas passer trop de temps sur internet. » (mots-clés : \zh{时间}, \zh{太多})
\end{itemize}
\end{exemplebox}

Notez vos deux meilleures hypothèses dans votre cahier : à la lecture détaillée (section suivante), vous vérifierez si elles étaient justes. Un élève qui vérifie une hypothèse lit activement ; un élève qui attend la traduction lit passivement.

\begin{exoresolu}[Associer le mot chinois à sa fonction]
Énoncé : reliez chaque mot-clé à la situation qui lui correspond.\\
1. \zh{上网} \quad 2. \zh{朋友圈} \quad 3. \zh{网友} \quad 4. \zh{手机} \quad 5. \zh{微信} \quad 6. \zh{社交软件}\\
a. publier une photo que ses amis peuvent commenter \quad b. l'objet qu'on emporte partout \quad c. ouvrir son navigateur ou une application \quad d. une personne rencontrée sur internet \quad e. l'application chinoise de messagerie \quad f. catégorie générale (WhatsApp, TikTok, WeChat...)
\tcblower
Solution détaillée :\\
1 → c : \zh{上网} (\emph{shàngwǎng}) signifie « se connecter à internet », littéralement « monter sur le réseau ».\\
2 → a : \zh{朋友圈} (\emph{péngyouquān}), le « cercle d'amis », est l'endroit où l'on publie photos et messages.\\
3 → d : \zh{网友} (\emph{wǎngyǒu}) = \zh{网} (réseau) + \zh{友} (ami).\\
4 → b : \zh{手机} (\emph{shǒujī}), la « machine de la main ».\\
5 → e : \zh{微信} (\emph{Wēixìn}), WeChat, littéralement « micro-message ».\\
6 → f : \zh{社交软件} (\emph{shèjiāo ruǎnjiàn}) = « logiciel social », hyperonyme de WeChat, WhatsApp, etc.
\end{exoresolu}

\begin{attention}[Faux ami : \zh{网友} n'est pas un « ami » comme les autres]
Un \zh{网友} est une connaissance \emph{en ligne} ; pour un ami réel, le chinois dit \zh{朋友} (péngyou). Ne traduisez donc pas « mon ami chinois » par \zh{我的网友} si vous parlez de votre camarade de classe ! De même, \zh{上网} est un verbe intransitif : on dit \zh{我每天都上网} (\emph{Wǒ měi tiān dōu shàngwǎng}, « je me connecte tous les jours »), jamais \zh{我上网微信} (incorrect). Pour dire « aller sur WeChat », on dit \zh{打开微信} (\emph{dǎkāi Wēixìn}, « ouvrir WeChat ») ou \zh{用微信} (\emph{yòng Wēixìn}, « utiliser WeChat »).
\end{attention}

\courssub{Comparaison Cameroun -- Chine : deux mondes numériques}

\begin{savaistu}[WeChat, l'application qui fait tout]
En Chine, on n'utilise pas WhatsApp ni Facebook (bloqués par le « Grand Pare-feu ») : l'application reine est \zh{微信} (Wēixìn, WeChat), qui sert à discuter, payer ses achats (\zh{微信支付}), commander un taxi (\zh{滴滴} intégré), et même prendre rendez-vous chez le médecin ! Là où un élève camerounais ouvre WhatsApp pour écrire à son correspondant, un lycéen de Shanghai ouvre WeChat. Et là où vous regardez TikTok, lui fait défiler \zh{抖音} (Dǒuyīn, Douyin), la version chinoise de la même application (même éditeur : ByteDance). Même passion pour l'écran, mots différents : voilà pourquoi ce vocabulaire est indispensable pour communiquer avec des jeunes Chinois.
\end{savaistu}

\begin{center}
\begin{tabularx}{\linewidth}{|X|X|X|}
\hline
\rowcolor{popblueL} Situation & Au Cameroun & En Chine \\
\hline
Messagerie quotidienne & WhatsApp & 微信 (Wēixìn, WeChat) \\
\hline
Vidéos courtes & TikTok & 抖音 (Dǒuyīn, Douyin) \\
\hline
Paiement par téléphone & Mobile Money (MTN MoMo, Orange Money) & 微信支付 (Wēixìn zhīfù) / 支付宝 (Zhīfùbǎo) \\
\hline
Fil d'actualité & Statuts WhatsApp / Facebook & 朋友圈 (péngyouquān, Moments) \\
\hline
Appel vidéo & WhatsApp / Messenger & 微信视频 (Wēixìn shìpín) \\
\hline
\end{tabularx}
\end{center}

\courssub{Annonce des objectifs de la leçon}

\begin{aretenir}[Ce que vous saurez faire à la fin de cette leçon]
\begin{itemize}
\item Reconnaître, prononcer (pinyin et tons) et employer le vocabulaire du thème « internet et réseaux sociaux » : \zh{网络}, \zh{上网}, \zh{手机}, \zh{微信}, \zh{朋友圈}, \zh{网友}, \zh{社交软件}...
\item Identifier les radicaux (\zh{网}, \zh{手}, \zh{彳}, \zh{月}, \zh{又}) et respecter l'ordre des traits pour écrire ces caractères.
\item Lire et comprendre un court texte chinois original sur ce thème, en appliquant la méthode en quatre étapes (titre → mots connus → écoute globale → lecture détaillée).
\item Répondre, en français puis en chinois, à des questions de compréhension du texte.
\item Comparer l'usage des réseaux sociaux au Cameroun et en Chine, et répondre à la question d'ouverture : \zh{你每天都上网吗？你用什么社交软件？}
\end{itemize}
\end{aretenir}

Les hypothèses sont posées, les mots-clés sont en place, la carte mentale et les radicaux sont dans vos têtes. Passons maintenant à la lecture détaillée du texte support \zh{我和我的手机}.

\coursec{Mise en route et découverte du thème}

\begin{experience}[Brainstorming : Internet et moi]
\begin{enumerate}
\item \textbf{Questionnaire éclair (oral) :} Le professeur pose les questions suivantes ; les élèves répondent brièvement en français ou en chinois (mots connus).
      \begin{itemize}
      \item Combien d'heures par jour passez-vous sur internet ?
      \item Quelles applications utilisez-vous le plus ? (WeChat, WhatsApp, TikTok, Douyin, Facebook, YouTube, Bilibili…)
      \item Avez-vous déjà acheté quelque chose en ligne ? Payé avec le Mobile Money (Orange Money, MTN MoMo) ou Alipay/WeChat Pay ?
      \item Quel est, selon vous, le plus grand avantage d'internet ? Le plus grand danger ?
      \end{itemize}
\item \textbf{Mise en commun au tableau :} Deux colonnes \textbf{Avantages 好处} / \textbf{Inconvénients 坏处}. Noter les mots-clés en chinois (avec pinyin) au fur et à mesure : \zh{方便 fāngbiàn}, \zh{快 kuài}, \zh{交朋友 jiāo péngyou}, \zh{眼睛坏 yǎnjing huài}, \zh{没时间学习 méi shíjiān xuéxí}…
\item \textbf{Transition :} « Nous allons maintenant lire un texte chinois qui structure exactement ces idées. Votre mission : comprendre le sens global et repérer le vocabulaire du thème. »
\end{enumerate}
\end{experience}

\coursec{Texte support : lecture et compréhension globale}

\courssub{Lecture du texte support}

\begin{textebox}[网络与生活 — Internet et la vie]
现在，互联网在我们的生活中越来越重要。\\
Xiànzài, hùliánwǎng zài wǒmen de shēnghuó zhōng yuè lái yuè zhòngyào.\\
很多年轻人每天都用手机上网。\\
Hěn duō niánqīngrén měitiān dōu yòng shǒujī shàngwǎng.\\
他们通过微信和朋友聊天，在朋友圈发照片，还在网上看视频、听音乐、买东西。\\
Tāmen tōngguò Wēixìn hé péngyou liáotiān, zài péngyouquān fā zhàopiàn, hái zài wǎngshang kàn shìpín, tīng yīnyuè, mǎi dōngxi.\\
网络让我们的生活更方便，也让我们能认识世界各地的朋友。\\
Wǎngluò ràng wǒmen de shēnghuó gèng fāngbiàn, yě ràng wǒmen néng rènshi shìjiè gèdì de péngyou.\\
但是，花太多时间上网对身体不好，也影响学习。\\
Dànshì, huā tài duō shíjiān shàngwǎng duì shēntǐ bù hǎo, yě yǐngxiǎng xuéxí.\\
所以，我们应该合理安排上网时间。\\
Suǒyǐ, wǒmen yīnggāi hélǐ ānpái shàngwǎng shíjiān.
\end{textebox}

\courssub{Traduction française intégrale, phrase par phrase}

\begin{enumerate}
\item \zh{现在，互联网在我们的生活中越来越重要。} — « Aujourd'hui, internet est de plus en plus important dans notre vie. »
\item \zh{很多年轻人每天都用手机上网。} — « Beaucoup de jeunes se connectent à internet avec leur téléphone portable tous les jours. »
\item \zh{他们通过微信和朋友聊天，在朋友圈发照片，还在网上看视频、听音乐、买东西。} — « Ils discutent avec leurs amis grâce à WeChat, publient des photos dans leurs "Moments", et regardent aussi des vidéos, écoutent de la musique et achètent des choses en ligne. »
\item \zh{网络让我们的生活更方便，也让我们能认识世界各地的朋友。} — « Internet rend notre vie plus pratique et nous permet aussi de connaître des amis de partout dans le monde. »
\item \zh{但是，花太多时间上网对身体不好，也影响学习。} — « Mais passer trop de temps sur internet est mauvais pour la santé et nuit aussi aux études. »
\item \zh{所以，我们应该合理安排上网时间。} — « C'est pourquoi nous devons organiser raisonnablement notre temps de connexion. »
\end{enumerate}

\begin{aretenir}[Idée générale du texte]
Le texte présente \cle{internet} dans la vie quotidienne : d'abord les \textbf{usages} (discuter, publier des photos, regarder des vidéos, acheter), puis les \textbf{avantages} (pratique, amis dans le monde entier), ensuite les \textbf{dangers} (santé, études), et enfin un \textbf{conseil} (bien organiser son temps). C'est une structure argumentative typique : constat → avantages → limites → conclusion.
\end{aretenir}

\courssub{Travail de lecture à voix haute}

\begin{methode}[Lire à voix haute : grouper par sens, pas caractère par caractère]
Une bonne lecture chinoise se fait par \textbf{groupes de sens}, jamais syllabe par syllabe. Procédez ainsi :
\begin{enumerate}
\item \textbf{Lecture chorale} : toute la classe lit ensemble après le professeur, en suivant le pinyin.
\item \textbf{Découpage en groupes de sens} : par exemple \zh{现在 / 互联网 / 在我们的生活中 / 越来越重要} (aujourd'hui / internet / dans notre vie / de plus en plus important). Marquez une petite pause entre les groupes.
\item \textbf{Lecture individuelle} : chaque élève lit une phrase ; le professeur corrige les tons.
\item \textbf{Vérification} : relire sans le pinyin, uniquement les caractères.
\end{enumerate}
\end{methode}

\begin{attention}[Pièges de prononciation pour les francophones]
\begin{itemize}
\item \zh{方便} se prononce \emph{fāngbiàn}, avec le \textbf{premier ton} sur 方 (son haut et plat). Les francophones disent souvent \emph{fángbiàn} (2\textsuperscript{e} ton) par habitude de monter la voix : c'est une erreur.
\item \zh{越来越} \emph{yuè lái yuè} : les deux 越 portent le 4\textsuperscript{e} ton, descendant et sec. Ne dites pas \emph{yuě} (3\textsuperscript{e} ton).
\item \zh{上网} \emph{shàngwǎng} : attention au 4\textsuperscript{e} ton de 上 et au 3\textsuperscript{e} ton plein de 网 en fin de groupe (ne pas le réduire en demi-ton 3).
\item \zh{东西} \emph{dōngxi} : le second caractère est atone (ton neutre), ne l'accentuez pas ; le mot signifie « chose, objet », pas « est et ouest » ici.
\item \zh{微信} \emph{Wēixìn} : 信 se prononce \emph{xìn} (4\textsuperscript{e} ton), pas \emph{xīn} (1\textsuperscript{e} ton, « nouveau »).
\end{itemize}
\end{attention}

\courssub{Observation grammaticale : la structure 让}

\begin{exemplebox}[Observer avant d'expliquer]
\begin{itemize}
\item \zh{网络让我们的生活更方便。} \emph{Wǎngluò ràng wǒmen de shēnghuó gèng fāngbiàn.} « Internet rend notre vie plus pratique. »
\item \zh{网络让我们能认识世界各地的朋友。} \emph{Wǎngluò ràng wǒmen néng rènshi shìjiè gèdì de péngyou.} « Internet nous permet de connaître des amis de partout dans le monde. »
\end{itemize}
\end{exemplebox}

\begin{propriete}[La structure causative 让 ràng]
\textbf{Structure :} A + 让 + personne + verbe / adjectif.\\
\textbf{Emploi :} 让 exprime « faire que », « permettre à », « rendre ». A est la cause ; la personne subit ou réalise l'action ou l'état.
\begin{itemize}
\item \textbf{Causatif adjectival (rendre) :} \zh{这个消息让我很高兴。} \emph{Zhège xiāoxi ràng wǒ hěn gāoxìng.} « Cette nouvelle me rend très heureux. »
\item \textbf{Causatif verbal (faire faire / permettre) :} \zh{老师让我们每天读课文。} \emph{Lǎoshī ràng wǒmen měitiān dú kèwén.} « Le professeur nous fait lire le texte tous les jours. » / \zh{请让我试一下。} \emph{Qǐng ràng wǒ shì yīxià.} « Permettez-moi d'essayer. »
\end{itemize}
\textbf{Comparaison avec le français :} le français utilise deux constructions différentes (« rendre + adjectif », « permettre de / faire + infinitif ») ; le chinois utilise la même structure 让 dans les deux cas.
\textbf{Synonymes de registre :} \zh{使} \emph{shǐ} (plus formel, écrit, ex. \zh{这使我想起了童年}) ; \zh{叫} \emph{jiào} (oral, familier, ex. \zh{别叫他等太久}).
\end{propriete}

\begin{exemplebox}[Autres structures clés du texte]
\begin{itemize}
\item \zh{越来越 + adjectif} \emph{yuè lái yuè} — « de plus en plus… ».
      \zh{汉语越来越有意思。} \emph{Hànyǔ yuè lái yuè yǒu yìsi.} « Le chinois devient de plus en plus intéressant. »
\item \zh{花 + temps/argent + (verbe)} \emph{huā} — « passer (dépenser) du temps/argent à… ».
      \zh{他花了两个小时做作业。} \emph{Tā huā le liǎng gè xiǎoshí zuò zuòyè.} « Il a passé deux heures à faire ses devoirs. » (Le \zh{了} marque l'achèvement).
\item \zh{合理安排} \emph{hélǐ ānpái} — « organiser raisonnablement / planifier judicieusement ».
      \zh{请合理安排你的假期。} \emph{Qǐng hélǐ ānpái nǐ de jiàqī.} « Organisez judicieusement vos vacances. »
\item \zh{通过…} \emph{tōngguò} — préposition « par le biais de, grâce à, via » (introduit le moyen/l'instrument).
      \zh{我们通过互联网学习中文。} \emph{Wǒmen tōngguò hùliánwǎng xuéxí Zhōngwén.} « Nous apprenons le chinois via internet. »
\end{itemize}
\end{exemplebox}

\courssub{Compréhension globale du texte}

Répondez d'abord en français, puis essayez en chinois en réutilisant les mots du texte :
\begin{enumerate}
\item De quoi parle ce texte ? (\emph{Le rôle d'internet dans la vie quotidienne des jeunes.})
\item Quels usages d'internet sont cités ? (\emph{Discuter sur WeChat, publier des photos dans les Moments, regarder des vidéos, écouter de la musique, acheter en ligne.})
\item Quels sont les avantages mentionnés ? (\emph{La vie est plus pratique ; on peut connaître des amis du monde entier.})
\item Quels sont les dangers mentionnés ? (\emph{Mauvais pour la santé / le corps ; nuit aux études.})
\item Quel conseil l'auteur donne-t-il ? (\emph{Organiser raisonnablement le temps passé en ligne.})
\end{enumerate}

\begin{popfigure}
\begin{center}
\begin{tikzpicture}[
  box/.style={draw=popline, rounded corners, thick, align=center, inner sep=6pt, font=\small},
  head/.style={font=\bfseries\small, align=center}]
\node[head, text=popgreen] (ha) at (0,0) {Avantages 好处};
\node[head, text=poppink] (hd) at (8,0) {Dangers 坏处};
\node[box, fill=popgreenL, text width=4.2cm] at (0,-1.4) {生活更方便\\ \emph{shēnghuó gèng fāngbiàn}\\ la vie plus pratique};
\node[box, fill=popgreenL, text width=4.2cm] at (0,-3.2) {认识世界各地的朋友\\ \emph{rènshi shìjiè gèdì de péngyou}\\ connaître des amis du monde};
\node[box, fill=poppinkL, text width=4.2cm] at (8,-1.4) {对身体不好\\ \emph{duì shēntǐ bù hǎo}\\ mauvais pour la santé};
\node[box, fill=poppinkL, text width=4.2cm] at (8,-3.2) {影响学习\\ \emph{yǐngxiǎng xuéxí}\\ nuit aux études};
\node[box, fill=popblueL, text width=5cm] at (4,-5) {Conclusion : 合理安排上网时间\\ \emph{hélǐ ānpái shàngwǎng shíjiān}\\ organiser raisonnablement son temps en ligne};
\draw[popline, thick, ->] (ha) -- ++(0,-0.6);
\draw[popline, thick, ->] (hd) -- ++(0,-0.6);
\draw[popline, thick, ->] (0,-4.2) -- (4,-4.8);
\draw[popline, thick, ->] (8,-4.2) -- (4,-4.8);
\end{tikzpicture}
\end{center}
\legende{Avantages et dangers d'internet selon le texte, et conclusion de l'auteur.}
\end{popfigure}

\courssub{Repérage des mots du glossaire dans le texte}

Avant d'étudier le glossaire complet (section suivante), surlignez dans le texte les mots suivants et vérifiez leur sens grâce au contexte :

\begin{center}
\begin{tabularx}{\linewidth}{|X|X|X|X|}
\hline
\rowcolor{popblueL} Caractères & Pinyin & Nature & Traduction \\
\hline
互联网 & hùliánwǎng & nom & internet \\
\hline
上网 & shàngwǎng & verbe & se connecter à internet, aller en ligne \\
\hline
通过 & tōngguò & préposition & par, grâce à, via \\
\hline
聊天 & liáotiān & verbe & bavarder, discuter (en ligne ou oral) \\
\hline
方便 & fāngbiàn & adjectif & pratique, commode \\
\hline
认识 & rènshi & verbe & connaître (une personne), faire connaissance \\
\hline
影响 & yǐngxiǎng & verbe / nom & influencer (souvent négatif), nuire à ; influence \\
\hline
安排 & ānpái & verbe / nom & organiser, planifier ; organisation, emploi du temps \\
\hline
\end{tabularx}
\end{center}

\begin{exoresolu}[Retrouver la phrase du texte]
Énoncé : retrouvez dans le texte la phrase qui signifie « Internet rend notre vie plus pratique » et recopiez-la en caractères, avec le pinyin.
\tcblower
Solution : on cherche le mot \cle{方便} (pratique) et la structure causative \cle{让}. La phrase est :\\
\zh{网络让我们的生活更方便。}\\
\emph{Wǎngluò ràng wǒmen de shēnghuó gèng fāngbiàn.}\\
Décomposition : 网络 (internet) + 让 (rend / fait que) + 我们的生活 (notre vie) + 更 (plus / davantage) + 方便 (pratique). Littéralement : « Internet fait que notre vie [est] plus pratique. »
\end{exoresolu}

\begin{exercice}[Vrai ou faux]
D'après le texte, dites si les phrases suivantes sont vraies (对 duì) ou fausses (不对 bú duì) :
\begin{enumerate}
\item \zh{很多年轻人每天都用手机上网。}
\item \zh{上网对身体很好。}
\item \zh{网络让我们能认识世界各地的朋友。}
\item \zh{我们应该花很多时间上网。}
\end{enumerate}
\end{exercice}

\begin{corrige}[Exercice « Vrai ou faux »]
\begin{enumerate}
\item \textbf{对} (vrai) : c'est la deuxième phrase du texte (avec \zh{都} « tous »).
\item \textbf{不对} (faux) : le texte dit \zh{花太多时间上网对身体不好} — passer \textbf{trop} de temps en ligne est mauvais pour la santé.
\item \textbf{对} (vrai) : phrase 4 du texte (\zh{也让我们能认识\ldots}).
\item \textbf{不对} (faux) : le texte conseille de \zh{合理安排上网时间}, d'organiser \textbf{raisonnablement} son temps, pas d'en passer beaucoup.
\end{enumerate}
\end{corrige}

\coursec{Glossaire du thème : vocabulaire détaillé}

\courssub{Liste complète du vocabulaire (HSK 1--3)}

\begin{center}
\begin{tabularx}{\linewidth}{|X|X|X|X|}
\hline
\rowcolor{popblueL} Caractères & Pinyin & Nature & Traduction / Contexte \\
\hline
互联网 & hùliánwǎng & nom & internet (formel) ; syn. \zh{网络 wǎngluò} (réseau), \zh{网 wǎng} (court) \\
\hline
上网 & shàngwǎng & VO (verbe+objet) & aller sur internet, se connecter (\zh{上} monter/sur + \zh{网} filet/réseau) \\
\hline
手机 & shǒujī & nom & téléphone portable (lit. « machine main ») \\
\hline
微信 & Wēixìn & nom propre & WeChat (application omniprésente en Chine : messagerie, paiement, réseau social) \\
\hline
朋友圈 & péngyouquān & nom & « Moments » (fil d'actualité WeChat) ; \zh{圈 quān} = cercle \\
\hline
发 & fā & verbe & envoyer, publier, émettre (\zh{发照片} publier des photos, \zh{发信息} envoyer un message) \\
\hline
照片 & zhàopiàn & nom & photo ; classif. \zh{张 zhāng} \\
\hline
视频 & shìpín & nom & vidéo ; classif. \zh{个 gè} ou \zh{部 bù} (film/série) \\
\hline
音乐 & yīnyuè & nom & musique \\
\hline
买东西 & mǎi dōngxi & VO & faire des achats, acheter des choses (\zh{东西} atone) \\
\hline
方便 & fāngbiàn & adjectif & pratique, commode ; \textbf{注意 ton 1 sur 方} \\
\hline
认识 & rènshi & verbe & connaître (personne, lieu), faire connaissance ; \zh{认识你很高兴} \\
\hline
世界各地 & shìjiè gèdì & locution & le monde entier, tous les coins du monde \\
\hline
身体 & shēntǐ & nom & corps, santé physique (\zh{身体健康} bonne santé) \\
\hline
影响 & yǐngxiǎng & v./n. & influencer, avoir un impact sur (souvent péjoratif) ; influence, impact \\
\hline
合理 & hélǐ & adjectif & raisonnable, logique, justifié (\zh{不合理} déraisonnable) \\
\hline
安排 & ānpái & v./n. & organiser, planifier, arranger ; arrangement, planning \\
\hline
时间 & shíjiān & nom & temps, heure ; \zh{安排时间} gérer son emploi du temps \\
\hline
通过 & tōngguò & prép. / verbe & via, par le biais de ; traverser, passer (examen) \\
\hline
聊天 & liáotiān & VO & bavarder ; \zh{聊} discuter + \zh{天} ciel/jour (originellement « parler du temps ») \\
\hline
越来越 & yuè lái yuè & adv. & de plus en plus (structure : \zh{越来越} + adj.) \\
\hline
花 & huā & verbe & dépenser (temps, argent) ; structure \zh{花 + 时间/金钱 + V} \\
\hline
太 & tài & adv. & trop, excessivement (suivi souvent de \zh{了} en fin de phrase pour l'excès) \\
\hline
应该 & yīnggāi & verbe modal & devoir, devrait (obligation morale / conseil) \\
\hline
\end{tabularx}
\end{center}

\courssub{Focus sur les classificateurs (quantificateurs) du thème}

\begin{center}
\begin{tabularx}{\linewidth}{|X|X|X|}
\hline
\rowcolor{poporangeL} Classificateur & Pinyin & Emploi typique (exemples du thème) \\
\hline
个 & gè & 一个视频、一个朋友、一个网站 \\
\hline
张 & zhāng & 一张照片、一张图片 \\
\hline
部 & bù & 一部电影、一部手机 (appareil complexe) \\
\hline
条 & tiáo & 一条信息、一条视频 (objet long/filiforme) \\
\hline
次 & cì & 上网一次、看两次视频 (fois, occurrence) \\
\hline
小时 & xiǎoshí & 两个小时 (durée, pas de \zh{个}) \\
\hline
\end{tabularx}
\end{center}

\coursec{Radicaux et ordre des traits}

\courssub{Analyse de caractères clés du texte}

Maîtriser les radicaux (clés) permet de deviner le sens et de retrouver le caractère dans un dictionnaire. L'ordre des traits (haut→bas, gauche→droite, horizontal→vertical, extérieur→intérieur→fermeture) assure une écriture lisible, rapide et esthétique.

\begin{center}
\begin{tabularx}{\linewidth}{|X|X|X|X|X|}
\hline
\rowcolor{poppurpleL} Caractère & Pinyin & Radical (Clé) & Nb traits (reste) & Ordre des traits (description) \\
\hline
网 & wǎng & \zh{纟} (sī, soie) & 6 & 1. Écrire 纟 (3 traits : point, crochet, horizontal) 2. 亡 (4 traits : horizontaux, verticale, point) \\
\hline
线 & xiàn & \zh{纟} (sī, soie) & 10 & 1. 纟 2. 泉 (quán, source : blanc + eau) — « fil de soie qui relie » \\
\hline
聊 & liáo & \zh{讠} (yán, parole) & 10 & 1. 讠 (2 traits : point, courbe) 2. 尞 (liáo, phonétique : petit + grand + feu) \\
\hline
影 & yǐng & \zh{彡} (shān, ornements/cheveux) & 12 & 1. 景 (jǐng, paysage : jour + soleil + petit) 2. 彡 (3 traits) — « ombre portée » \\
\hline
安 & ān & \zh{宀} (mián, toit) & 3 & 1. 宀 (3 traits : point, pli, horizontal) 2. 女 (nǚ, femme : courbe, trait, horizontal) — « femme sous le toit = paix » \\
\hline
\end{tabularx}
\end{center}

\begin{methode}[Entraînement à l'écriture sur cahier]
\begin{enumerate}
\item \textbf{Grille 10×10} : écrire chaque caractère 5 fois en respectant la grille (centrer le radical).
\item \textbf{Dictée de caractères} : le professeur donne le pinyin + définition ; l'élève écrit le caractère.
\item \textbf{Recomposition} : mélanger les composants (ex. 纟 + 亡 → 网 ; 讠 + 尞 → 聊) et les reassembler.
\end{enumerate}
\end{methode}

\coursec{Questions de compréhension du texte (approfondissement)}

Répondez en phrases complètes en chinois (caractères + pinyin) :

\begin{exercice}[Questions ouvertes]
\begin{enumerate}
\item \zh{现在互联网在我们生活中怎么样？}
\item \zh{年轻人通过什么和朋友聊天？}
\item \zh{网络让我们的生活变得怎样？}
\item \zh{花太多时间上网有什么坏处？}
\item \zh{作者建议我们怎么做？}
\end{enumerate}
\end{exercice}

\begin{corrige}[Questions ouvertes — propositions de réponses]
\begin{enumerate}
\item \zh{现在互联网在我们生活中越来越重要。} (Xiànzài hùliánwǎng zài wǒmen shēnghuó zhōng yuè lái yuè zhòngyào.)
\item \zh{他们通过微信和朋友聊天。} (Tāmen tōngguò Wēixìn hé péngyou liáotiān.)
\item \zh{网络让我们的生活更方便。} (Wǎngluò ràng wǒmen de shēnghuó gèng fāngbiàn.)
\item \zh{对身体不好，也影响学习。} (Duì shēntǐ bù hǎo, yě yǐngxiǎng xuéxí.)
\item \zh{我们应该合理安排上网时间。} (Wǒmen yīnggāi hélǐ ānpái shàngwǎng shíjiān.)
\end{enumerate}
\end{corrige}

\coursec{Production guidée et comparaison Cameroun – Chine}

\courssub{Expression écrite : Mon utilisation d'internet}

\begin{methode}[Rédiger un court paragraphe (60--80 caractères)]
\textbf{Consigne :} Écrivez 5--6 phrases pour décrire vos habitudes et votre opinion. Utilisez le vocabulaire et les structures de la leçon (\zh{让}, \zh{越来越}, \zh{花时间}, \zh{合理安排}, \zh{通过}).
\textbf{Plan suggéré :}
\begin{enumerate}
\item Fréquence / support : \zh{我每天用手机/电脑上网。}
\item Activités (3 minimum) : \zh{我通过微信/WhastApp和朋友聊天，看视频，买东西。}
\item Avantage ressenti : \zh{网络让我的生活更方便，也让我认识了新朋友。}
\item Danger / limite : \zh{但是花太多时间上网对眼睛不好，影响睡眠/学习。}
\item Conclusion / résolution : \zh{所以我要合理安排上网时间，多运动，多读书。}
\end{enumerate}
\end{methode}

\begin{exemplebox}[Modèle de production (niveau Terminale A)]
\zh{我每天都用手机上网。通过微信和朋友聊天，发照片，还看短视频。网络让我的生活更方便，让我认识世界各地的朋友。但是花太多时间上网对眼睛不好，也影响学习。所以我应该合理安排上网时间，少看手机，多读书、多运动。}\\
\emph{Wǒ měitiān dōu yòng shǒujī shàngwǎng. Tōngguò Wēixìn hé péngyou liáotiān, fā zhàopiàn, hái kàn duǎn shìpín. Wǎngluò ràng wǒ de shēnghuó gèng fāngbiàn, ràng wǒ rènshi shìjiè gèdì de péngyou. Dànshì huā tài duō shíjiān shàngwǎng duì yǎnjing bù hǎo, yě yǐngxiǎng xuéxí. Suǒyǐ wǒ yīnggāi hélǐ ānpái shàngwǎng shíjiān, shǎo kàn shǒujī, duō dúshū, duō yùndòng.}
\end{exemplebox}

\courssub{Débat oral : Cameroun vs Chine — Usages du numérique}

\begin{experience}[Table ronde : Deux mondes connectés]
\textbf{Modalités :} Groupes de 4 (2 « Cameroun », 2 « Chine » — rôles tirés au sort). 10 min préparation, 5 min débat, 3 min restitution par groupe.
\textbf{Thèmes imposés (vocabulaire fourni) :}
\begin{center}
\begin{tabularx}{\linewidth}{|X|X|X|}
\hline
\rowcolor{popgoldL} Sujet & Cameroun (Réalités 2024) & Chine (Réalités 2024) \\
\hline
Réseau social principal & WhatsApp, Facebook, TikTok & \zh{微信 Wēixìn} (super-app), \zh{抖音 Dǒuyīn}, \zh{小红书 Xiǎohóngshū} \\
\hline
Paiement mobile & \zh{手机支付} : Orange Money, MTN MoMo (USSD/App) & \zh{微信支付 Wēixìn Zhīfù}, \zh{支付宝 Zhīfùbǎo} (QR code omniprésent) \\
\hline
Achats en ligne & Jumia, Kaymu, boutiques Instagram/WhastApp & \zh{淘宝 Táobǎo}, \zh{京东 Jīngdōng}, \zh{拼多多 Pīnduōduō}, livraison 24h \\
\hline
Accès internet & 4G villes, fibre en déploiement, cybercafés encore utiles & 5G généralisée, wifi gratuit partout, très haut débit \\
\hline
Régulation & Libre, lois cybercriminalité en cours & \zh{防火墙} (Great Firewall), intranet national, identité réelle obligatoire \\
\hline
\end{tabularx}
\end{center}

\textbf{Phrases-outils pour le débat :}
\begin{itemize}
\item \zh{在喀麦隆，我们习惯用\ldots} (Au Cameroun, nous avons l'habitude d'utiliser…)
\item \zh{在中国，\ldots 更普遍。} (En Chine, … est plus répandu.)
\item \zh{我觉得\ldots 很方便，因为\ldots} (Je trouve … très pratique, car…)
\item \zh{但是\ldots 也有问题，比如\ldots} (Mais … a aussi des problèmes, par ex.…)
\item \zh{我们可以向中国学习\ldots} (Nous pouvons apprendre de la Chine …)
\end{itemize}

\textbf{Consigne :} Comparez \textbf{un} point précis (ex. paiement, réseau social, livraison). Dites ce qui vous semble meilleur chez l'un ou l'autre, et pourquoi. Utilisez \zh{比 bǐ} (comparatif) et \zh{更 gèng} / \zh{最 zuì}.
\end{experience}

\begin{savaistu}[Le saviez-vous ? — Le « Grand Firewall » et l'intranet chinois]
En Chine, l'internet international est filtré par le \textbf{Grand Pare-feu (GFW / 防火墙 Fánghuǒqiáng)}. Google, YouTube, Facebook, WhatsApp, Instagram, Wikipédia sont inaccessibles sans VPN (légal seulement pour entreprises/universités agréées). Les Chinois utilisent un écosystème parallèle complet : \zh{百度 Bǎidù} (moteur de recherche), \zh{哔哩哔哩 Bīlībīlī} (vidéos), \zh{知乎 Zhīhū} (Q\&R), \zh{微博 Wēibó} (microblog). Ce n'est pas une « absence » d'internet, mais un \textbf{intranet national souverain} de 1 milliard d'utilisateurs, le plus grand au monde. Au Cameroun, l'accès est ouvert mais la fracture numérique (coût, infrastructure, électricité) reste le vrai défi.
\end{savaistu}

\begin{attention}[Erreurs fréquentes à éviter en production]
\begin{itemize}
\item \textbf{Ordre des mots (Complément de moyen) :} \zh{我用手机上网} (Suj + 用 + Outil + V) \textbf{PAS} \zh{我上网用手机}.
\item \textbf{让 vs 使 vs 叫 :} À l'écrit scolaire, privilégiez \zh{让} (standard) ou \zh{使} (formel). Évitez \zh{叫}.
\item \textbf{花 + temps :} \zh{我花两个小时上网} (correct) ; \zh{我在上网花两个小时} (incorrect, place du complément de durée).
\item \textbf{方便 vs 方便 :} N'écrivez pas \zh{fangbian} sans tons. \zh{方} = 1\textsuperscript{er} ton (haut), \zh{便} = 4\textsuperscript{e} ton (descendant).
\item \textbf{认识 vs 知道 :} \zh{认识} = connaître une personne/lieu (expérience) ; \zh{知道} = savoir une info/fait. \zh{我知道他} (je sais qui il est) vs \zh{我认识他} (on s'est rencontrés).
\end{itemize}
\end{attention}

\begin{aretenir}[Bilan de la leçon 34 — Check-list compétences]
\begin{itemize}
\item [\checkmark] \textbf{Vocabulaire :} 25+ mots/thème (noms, verbes, adj., prép., adv.) — reconnaître, prononcer, classer par radical.
\item [\checkmark] \textbf{Lecture :} Texte 6 phrases, HSK 2-3 — lecture fluide par groupes de sens, tons corrects.
\item [\checkmark] \textbf{Grammaire :} \zh{让} (causatif/permissif), \zh{越来越+Adj}, \zh{花+Temps+V}, \zh{通过+Moyen}, \zh{合理安排}.
\item [\checkmark] \textbf{Compréhension :} V/F, Q/R ouvertes, résumé structuré (Avantages/Inconvénients/Conseil).
\item [\checkmark] \textbf{Écriture :} 5 radicaux (纟, 讠, 彡, 宀, 女) + ordre traits ; rédaction paragraphe guidé.
\item [\checkmark] \textbf{Oral :} Lecture chorale/individuelle, débat comparatif Cameroun/Chine avec phrases-outils.
\item [\checkmark] \textbf{Culture :} Écosystème numérique chinois (WeChat, Alipay, GFW) vs camerounais (Mobile Money, WhatsApp, 4G).
\end{itemize}
\end{aretenir}

\coursec{Glossaire du thème : Internet et les réseaux sociaux}

Après la lecture globale du texte support, nous allons maintenant explorer en détail le vocabulaire clé du thème « Internet et les réseaux sociaux ». Cette section vous permettra de maîtriser la prononciation, le sens, la nature grammaticale et les emplois courants de chaque mot, afin de pouvoir les réutiliser à l'oral comme à l'écrit.

\courssub{Tableau récapitulatif du vocabulaire}

Le tableau ci-dessous regroupe les 17 mots et expressions du thème. Apprenez-les par cœur en associant caractères, pinyin (avec tons), nature et traduction. La colonne « Nature » utilise les abréviations standard : n. = nom, v. = verbe, adj. = adjectif, n. propre = nom propre, loc. v. = locution verbale.

\begin{center}
\begin{tabularx}{\linewidth}{|X|X|X|X|}
\hline
\rowcolor{popblue} \textbf{Caractères} & \textbf{Pinyin} & \textbf{Nature} & \textbf{Traduction} \\
\hline
\zh{网络} & wǎngluò & n. & internet, réseau (général) \\
\hline
\rowcolor{popblueL}
\zh{互联网} & hùliánwǎng & n. & internet (le réseau mondial) \\
\hline
\zh{上网} & shàngwǎng & v. & se connecter à internet, surfer \\
\hline
\rowcolor{popblueL}
\zh{手机} & shǒujī & n. & téléphone portable \\
\hline
\zh{微信} & Wēixìn & n. propre & WeChat \\
\hline
\rowcolor{popblueL}
\zh{朋友圈} & péngyouquān & n. & fil d'actualité (Moments) \\
\hline
\zh{网友} & wǎngyǒu & n. & ami(e) en ligne \\
\hline
\rowcolor{popblueL}
\zh{聊天} & liáotiān & v. & bavarder, discuter (en ligne) \\
\hline
\zh{发} & fā & v. & envoyer, publier, poster \\
\hline
\rowcolor{popblueL}
\zh{照片} & zhàopiàn & n. & photo \\
\hline
\zh{视频} & shìpín & n. & vidéo \\
\hline
\rowcolor{popblueL}
\zh{软件} & ruǎnjiàn & n. & logiciel, application \\
\hline
\zh{方便} & fāngbiàn & adj. & pratique, commode \\
\hline
\rowcolor{popblueL}
\zh{影响} & yǐngxiǎng & v./n. & influencer / influence \\
\hline
\zh{安排} & ānpái & v./n. & organiser, planifier / organisation \\
\hline
\rowcolor{popblueL}
\zh{时间} & shíjiān & n. & temps \\
\hline
\zh{加好友} & jiā hǎoyǒu & loc. v. & ajouter un ami (sur un réseau social) \\
\hline
\end{tabularx}
\end{center}

\courssub{Prononciation guidée et travail des tons}

Avant d'utiliser ces mots, assurons-nous d'une prononciation exacte. En classe, nous procédons en deux temps : \textbf{répétition chorale} (toute la classe ensemble) puis \textbf{répétition individuelle} (chaque élève passe au tableau ou s'enregistre).

\begin{definition}[Points d'attention phonétique]
\begin{itemize}
    \item \textbf{Le 3e ton « plein » vs « demi-ton » (变调 biàndiào) :} \zh{网} \emph{wǎng} (dans \zh{网络}, \zh{上网}, \zh{网友}, \zh{网上}) est un 3e ton isolé (descendant-remontant : 214). Enchaîné \textbf{devant un 4e ton}, il devient un « demi-3e ton » (bas seulement : 21). Exemple : \zh{网络} \emph{wǎngluò} $\rightarrow$ \emph{wǎngluò} (3e demi-ton + 4e ton). Devant un 1er, 2e ou 3e ton, il reste bas. Devant une pause ou en fin d'énoncé (\zh{上网}), il est souvent plein (214).
    \item \textbf{Tons neutres (轻声 qīngshēng) :} Aucun mot de cette liste n'a de ton neutre standard dans sa citation dictionnaire. Attention cependant à \zh{朋友圈} \emph{péngyouquān} : \zh{圈} garde son 1er ton (quān), ne le réduisez pas.
    \item \textbf{Ü (yu) :} \zh{软} \emph{ruǎn} contient le son \emph{ü} (écrit u après j, q, x, y ; écrit ü après n, l). Ici \emph{ruǎn} s'entend [ʐuan] (comme « rwan » français mais avec la pointe de la langue retroflexe).
    \item \textbf{Rétroflexes (zh, ch, sh, r) :} \zh{上} \emph{shàng}, \zh{手} \emph{shǒu}, \zh{视} \emph{shì}, \zh{上} \emph{shàng} (dans \zh{上网}) exigent la pointe de la langue relevée vers le palais dur. \zh{网} \emph{wǎng} n'est \textbf{pas} une rétroflexe (c'est une vélaire).
    \item \textbf{Noms propres :} \zh{微信} \emph{Wēixìn} : majuscule au W, 1er ton + 4e ton. C'est une marque déposée (Tencent).
\end{itemize}
\end{definition}

\begin{exemplebox}[Modèles de prononciation à répéter]
\begin{itemize}
    \item \zh{我每天上网两个小时。} \emph{Wǒ měitiān shàngwǎng liǎng gè xiǎoshí.} (Je surfe deux heures par jour.)
    \item \zh{请加我的微信。} \emph{Qǐng jiā wǒ de Wēixìn.} (Ajoute mon WeChat s'il te plaît.)
    \item \zh{这个软件很方便。} \emph{Zhège ruǎnjiàn hěn fāngbiàn.} (Cette appli est très pratique.)
\end{itemize}
\end{exemplebox}

\courssub{Distinction essentielle : \zh{网络} vs \zh{互联网}}

C'est une confusion classique chez les francophones qui traduisent les deux par « internet ».

\begin{propriete}[网络 (wǎngluò) vs 互联网 (hùliánwǎng)]
\begin{center}
\begin{tabularx}{\linewidth}{|X|X|X|}
\hline
\rowcolor{popblue} \textbf{Mot} & \textbf{Sens précis} & \textbf{Exemple d'emploi} \\
\hline
\zh{网络} \emph{wǎngluò} & « Réseau » au sens large (réseau local, réseau social, réseau de contacts, réseau informatique). & \zh{校园网络} \emph{xiàoyuán wǎngluò} (réseau du campus / intranet scolaire), \zh{社交网络} \emph{shèjiāo wǎngluò} (réseau social). \\
\hline
\zh{互联网} \emph{hùliánwǎng} & « Internet » spécifiquement : le réseau mondial interconnecté (Inter-connected Network). & \zh{互联网时代} \emph{hùliánwǎng shídài} (l'ère d'Internet), \zh{接入互联网} \emph{jiērù hùliánwǎng} (accéder à Internet). \\
\hline
\end{tabularx}
\end{center}
\textbf{Règle mnémotechnique :} \zh{互} (mutuel) + \zh{联} (connecter) + \zh{网} (réseau) = « réseau d'interconnexion mutuelle » $\rightarrow$ \textbf{Internet mondial}. \zh{网络} est l'hyperonyme (terme plus général).
\end{propriete}

\begin{attention}[Erreur fréquente : « Je n'ai pas de 网络 »]
Si vous voulez dire « Je n'ai pas de connexion Internet / pas de data », dites :
\zh{我没网。} \emph{Wǒ méi wǎng.} (familier, très courant) ou \zh{我连不上互联网。} \emph{Wǒ lián bù shàng hùliánwǎng.}
Ne dites pas \zh{我没有互联网} (sonne comme « je ne possède pas l'internet mondial »).
\end{attention}

\courssub{Collocations indispensables (搭配 dāpèi)}

En chinois, les mots ne vivent pas seuls. Mémorisez ces \textbf{blocs lexicaux} (collocations) pour parler naturellement.

\begin{propriete}[Collocations clés du thème]
\begin{center}
\begin{tabularx}{\linewidth}{|>{\bfseries}X|X|X|}
\hline
\rowcolor{popblue} \textbf{Collocation (chinois + pinyin)} & \textbf{Structure} & \textbf{Traduction} \\
\hline
\zh{上网} \emph{shàng wǎng} & Verbe séparable (V + O) & Se connecter, surfer \\
\hline
\zh{发朋友圈} \emph{fā péngyouquān} & Verbe + Objet composé & Publier sur son fil (Moments) \\
\hline
\zh{加好友} \emph{jiā hǎoyǒu} & Verbe + Objet & Ajouter un ami (demande d'ami) \\
\hline
\zh{看视频} \emph{kàn shìpín} & Verbe + Objet & Regarder des vidéos \\
\hline
\zh{网上购物} \emph{wǎngshàng gòuwù} & Locution adverbiale + VO & Faire des achats en ligne \\
\hline
\zh{聊天} \emph{liáo tiān} & Verbe séparable (V + O) & Bavarder / Discuter en ligne \\
\hline
\zh{发照片 / 发视频 / 发信息} \emph{fā zhàopiàn / shìpín / xìnxī} & Verbe + Objet & Envoyer/Poster photo / vidéo / message \\
\hline
\end{tabularx}
\end{center}
\end{propriete}

\begin{exemplebox}[Exemples contextuels traduits]
\begin{enumerate}
    \item \zh{我每天用手机上网。} \emph{Wǒ měitiān yòng shǒujī shàngwǎng.} — Je me connecte chaque jour avec mon téléphone.
    \item \zh{她在朋友圈发了很多照片。} \emph{Tā zài péngyouquān fā le hěn duō zhàopiàn.} — Elle a publié beaucoup de photos sur son fil (Moments).
    \item \zh{我想加你为好友，可以吗？} \emph{Wǒ xiǎng jiā nǐ wéi hǎoyǒu, kěyǐ ma?} — Je veux t'ajouter comme ami, est-ce possible ?
    \item \zh{我们在微信上聊了半天天。} \emph{Wǒmen zài Wēixìn shàng liáo le bàntiān.} — Nous avons bavardé sur WeChat pendant une demi-journée.
    \item \zh{网上购物很方便，不用出门。} \emph{Wǎngshàng gòuwù hěn fāngbiàn, bùyòng chūmén.} — Faire des achats en ligne est pratique, pas besoin de sortir.
\end{enumerate}
\end{exemplebox}

\courssub{Focus grammatical : Le verbe \zh{发} (fā) et le verbe séparable \zh{上网}}

Deux points de grammaire lexical sont cruciaux ici pour éviter les fautes de structure de phrase.

\begin{propriete}[Le verbe polyvalent \zh{发} fā : « envoyer / publier / émettre »]
En français, on change de verbe : \emph{poster} (sur un mur), \emph{envoyer} (un message privé), \emph{publier} (une vidéo), \emph{diffuser} (en direct). En chinois, \zh{发} couvre tous ces sens avec un objet direct.
\begin{center}
\begin{tabularx}{\linewidth}{|X|X|}
\hline
\rowcolor{popblue} \textbf{Objet} & \textbf{Traduction contextuelle} \\
\hline
\zh{发朋友圈} \emph{fā péngyouquān} & Publier sur son fil public (Moments) \\
\zh{发照片} \emph{fā zhàopiàn} & Envoyer/Poster une photo \\
\zh{发视频} \emph{fā shìpín} & Poster/Envoyer une vidéo \\
\zh{发信息 / 发消息} \emph{fā xìnxī / fā xiāoxi} & Envoyer un message (texte, vocal) \\
\zh{发语音} \emph{fā yǔyīn} & Envoyer un message vocal \\
\zh{发红包} \emph{fā hóngbāo} & Envoyer une enveloppe rouge (argent) \\
\hline
\end{tabularx}
\end{center}
\textbf{Structure :} Sujet + \zh{发} + Objet (+ \zh{了} pour l'accompli).
\end{propriete}

\begin{attention}[Piège majeur : \zh{上网} est un VERBE SÉPARABLE (离合词 líhécí)]
\zh{上网} = \zh{上} (verbe : monter, aller sur) + \zh{网} (nom/objet : le réseau/Internet).
\textbf{Conséquence :} On ne peut rien insérer entre le verbe et son objet \textbf{sauf} les marqueurs d'aspect (\zh{了}, \zh{过}, \zh{着}) ou les compléments de durée/fréquence \textbf{après le verbe \zh{上}}.
\begin{itemize}
    \item ✓ \zh{我上了两个小时网。} \emph{Wǒ shàng le liǎng gè xiǎoshí wǎng.} (J'ai surfé 2h.) $\leftarrow$ Durée \textbf{après} \zh{上} + \zh{了}.
    \item ✓ \zh{你上网了吗？} \emph{Nǐ shàngwǎng le ma?} (T'as surfé ?) $\leftarrow$ \zh{了} aspectuel après le verbe \zh{上} (forme non séparée acceptée pour l'aspect parfaitif simple).
    \item ✗ \zh{我上网了两个小时。} \emph{Wǒ shàngwǎng le liǎng gè xiǎoshí.} (FAUX : \zh{了} collé à \zh{网} suggère « l'internet est fini »).
    \item ✗ \zh{我上了网两个小时。} (Ambigu, évitez).
\end{itemize}
\textbf{Astuce :} Traitez \zh{上网} comme « aller sur le net ». On dit « Je suis \textbf{allé} deux heures \textbf{sur le net} », pas « J'ai \textbf{allé-sur-le-net} deux heures ».
\end{attention}

\begin{exoresolu}[Application : Transformer les phrases avec la durée]
\textbf{Énoncé :} Réécrivez les phrases en insérant la durée indiquée entre parenthèses.
\begin{enumerate}
    \item 我上网。(两个小时) $\rightarrow$ \textbf{Correction :} \zh{我上了两个小时网。} \emph{Wǒ shàng le liǎng gè xiǎoshí wǎng.}
    \item 他每天聊天。(半小时) $\rightarrow$ \textbf{Correction :} \zh{他每天聊半小时天。} \emph{Tā měitiān liáo bàn xiǎoshí tiān.} (\zh{聊天} aussi séparable : \zh{聊} + \zh{天}).
    \item 我们发朋友圈。(三次) $\rightarrow$ \textbf{Correction :} \zh{我们发了三次朋友圈。} \emph{Wǒmen fā le sān cì péngyouquān.} (\zh{发朋友圈} : \zh{发} n'est pas séparable de \zh{朋友圈} ici, \zh{朋友圈} est un nom composé unique. On met \zh{了} après \zh{发}. Le classifieur \zh{次} (fois) quantifie l'action).
\end{enumerate}
\tcblower
\textbf{Solution détaillée :}
1. \zh{上网} séparable $\rightarrow$ Durée après \zh{上} + \zh{了}.
2. \zh{聊天} séparable $\rightarrow$ Durée après \zh{聊} (pas de \zh{了} car habitude \zh{每天}).
3. \zh{发朋友圈} : \zh{发} est verbe transitif direct, \zh{朋友圈} est l'objet unique. \zh{了} marque l'accompli après \zh{发}. Le classifieur \zh{次} (fois) quantifie l'action.
\end{exoresolu}

\courssub{Le saviez-vous ? : Étymologie et société}

\begin{savaistu}[Le mot \zh{网友} (wǎngyǒu) : un néologisme des années 90]
Le terme \zh{网友} est la contraction de \zh{网络} (réseau) + \zh{朋友} (ami) $\rightarrow$ \zh{网} + \zh{友}.
\begin{itemize}
    \item Il apparaît en Chine continentale au milieu des années 1990 avec la démocratisation de l'accès à Internet (modems 56k, cybercafés \zh{网吧} \emph{wǎngbā}).
    \item Au début, il désignait spécifiquement les personnes rencontrées sur les \textbf{BBS} (Bulletin Board Systems) et les premiers forums (如 \zh{天涯社区} \emph{Tiānyá Shèqū}, \zh{猫扑} \emph{Māopū}).
    \item Aujourd'hui, il couvre tout contact connu uniquement en ligne (WeChat, QQ, Weibo, Douyin, jeux vidéo). On distingue souvent \zh{现实朋友} \emph{xiànshí péngyou} (amis dans la vraie vie / IRL) et \zh{网友}.
    \item \textbf{Comparaison Cameroun-Chine :} Au Cameroun, on parle souvent de « contacts WhatsApp » ou « amis Facebook ». En Chine, \zh{微信好友} \emph{Wēixìn hǎoyǒu} (amis WeChat) est la norme professionnelle et personnelle ; ajouter un \zh{网友} sur WeChat se dit \zh{加微信} \emph{jiā Wēixìn} (souvent en scannant un \zh{二维码} \emph{èrwéimǎ} — QR code).
\end{itemize}
\end{savaistu}

\courssub{Entraînement rapide : Association et Traduction}

\courssub{Exercice 1 : Association Mot — Définition}
Associez chaque terme à sa définition.
\begin{exercice}[Association]
\begin{enumerate}
    \item \zh{软件} \hfill A. L'endroit où l'on voit les photos de ses amis sur WeChat.
    \item \zh{朋友圈} \hfill B. L'action de scanner un code pour se connecter / ajouter un contact.
    \item \zh{网友} \hfill C. WeChat, TikTok (Douyin), QQ sont des exemples.
    \item \zh{二维码} (supplément) \hfill D. Une personne qu'on ne connaît que par écran interposé.
\end{enumerate}
\end{exercice}
\begin{corrige}[Exercice 1]
1-C, 2-A, 3-D, 4-B.
\end{corrige}

\courssub{Exercice 2 : Quiz de traduction bidirectionnel}
Traduisez sans regarder le tableau ci-dessus.

\begin{exercice}[Chinois $\rightarrow$ Français]
\begin{enumerate}
    \item 这个软件对学习中文有影响。
    \item 请安排一下时间，我们视频聊天。
    \item 我在网上买了一件衣服，很方便。
\end{enumerate}
\end{exercice}
\begin{corrige}[Exercice 2 (Chinois $\rightarrow$ Français)]
1. Ce logiciel a une influence sur l'apprentissage du chinois. (\zh{影响} = n. influence / v. influencer). \\
2. Organise un peu le temps, on fera un appel vidéo / on bavardera en vidéo. (\zh{视频} peut être nom « vidéo » ou verbe/adverbe « par vidéo » ici \zh{视频聊天}). \\
3. J'ai acheté un vêtement en ligne, c'est très pratique. (\zh{网上} = en ligne / sur le net ; \zh{买} \emph{mǎi} acheter).
\end{corrige}

\begin{exercice}[Français $\rightarrow$ Chinois (caractères + pinyin)]
\begin{enumerate}
    \item Je veux t'ajouter sur WeChat. (Utiliser \zh{加} + \zh{微信}).
    \item Elle a posté une vidéo hier.
    \item Nous n'avons pas le temps de surfer aujourd'hui.
\end{enumerate}
\end{exercice}
\begin{corrige}[Exercice 2 (Français $\rightarrow$ Chinois)]
1. \zh{我想加你微信。} \emph{Wǒ xiǎng jiā nǐ Wēixìn.} (Ou \zh{加你为好友}). \\
2. \zh{她昨天发了一个视频。} \emph{Tā zuótiān fā le yí gè shìpín.} (Classifieur \zh{个} ou \zh{条} \emph{tiáo} pour vidéo). \\
3. \zh{我们今天没时间上网。} \emph{Wǒmen jīntiān méi shíjiān shàngwǎng.}
\end{corrige}

\courssub{Texte d'entraînement complémentaire (Approfondissement)}

Ce texte original réemploie \textbf{au moins 12 mots du glossaire}. Lisez-le à haute voix, soulignez les mots connus, puis répondez aux questions.

\begin{textebox}[La vie numérique de Léa à Yaoundé]
莱雅住在雅温得。她是高三学生，每天很忙。但是，她每天晚上都要上网。她用手机登录微信，看朋友圈。她的网友很多，有喀麦隆的同学，也有中国的朋友。周末的时候，她喜欢看视频，学习做照片编辑的软件。她觉得网上购物很方便，不用去市场。可是，妈妈说：“上网要安排好时间，不要影响学习。” 莱雅回答：“我知道，我会安排好的。”
\end{textebox}

\courssub{Exercice 3 : Questions de compréhension sur le texte de Léa}
\begin{exercice}[Compréhension écrite]
\begin{enumerate}
    \item Où habite Léa et quel est son statut ?
    \item Quand se connecte-t-elle ? Avec quel appareil ?
    \item Que fait-elle sur WeChat ? (Deux actions précises).
    \item Qui sont ses \zh{网友} ?
    \item Quelle activité fait-elle le week-end liée aux logiciels ?
    \item Pourquoi aime-t-elle \zh{网上购物} ?
    \item Quel conseil sa mère lui donne-t-elle ? (Utilisez \zh{安排} et \zh{影响}).
    \item Comment Léa répond-elle ? (Traduisez sa réponse).
\end{enumerate}
\end{exercice}
\begin{corrige}[Exercice 3]
1. Elle habite à Yaoundé (\zh{雅温得} Yǎwēndé) et est en Terminale (\zh{高三} gāosān). \\
2. Chaque soir (\zh{每天晚上} \emph{měitiān wǎnshang}) avec son téléphone portable (\zh{手机}). \\
3. Elle regarde le fil d'actualité (\zh{看朋友圈}) et discute/bavarde (\zh{聊天} implicite). Le texte dit explicitement : \zh{看朋友圈}. \\
4. Des camarades camerounais et des amis chinois (\zh{喀麦隆的同学，也有中国的朋友}). \\
5. Elle regarde des vidéos (\zh{看视频}) pour apprendre à utiliser des logiciels de retouche photo (\zh{学习做照片编辑的软件}). \\
6. Parce que c'est pratique (\zh{方便}), pas besoin d'aller au marché (\zh{不用去市场} \emph{bùyòng qù shìchǎng}). \\
7. « Il faut bien organiser ton temps pour surfer, ne laisse pas internet influencer tes études. » (\zh{上网要安排好时间，不要影响学习。}) \\
8. « Je sais, je m'en occuperai / je l'organiserai bien. » (\zh{我知道，我会安排好的。} \emph{Wǒ zhīdào, wǒ huì ānpái hǎo de.}).
\end{corrige}

\courssub{Synthèse orthographique : Clés (radicaux) et ordre des traits}

Pour écrire correctement les caractères nouveaux de cette leçon, identifions leur clé (radical \zh{部首} \emph{bùshǒu}) et l'ordre des traits principal.

\begin{definition}[Analyse graphémique de 6 caractères clés]
\begin{center}
\begin{tabularx}{\linewidth}{|c|X|X|X|}
\hline
\rowcolor{popblue} \textbf{Caractère} & \textbf{Clé (Radical)} & \textbf{Nb traits (reste / total)} & \textbf{Ordre des traits (règle principale)} \\
\hline
\zh{网} \emph{wǎng} & \zh{纟} (soie, fil) — 3 traits & 3 / 6 & 1. \zh{纟} : 撇 (piě), 横折钩 (héng zhé gōu), 提 (tí). 2. \zh{亡} (partie basse) : 横 (héng), 横 (héng), 竖 (shù). *Note : le point final de \zh{亡} est omis dans le comptage officiel (6 traits).* \\
\hline
\zh{互} \emph{hù} & \zh{二} (deux) — 2 traits & 2 / 4 & Haut $\rightarrow$ bas : deux horizontales \zh{二} (la 2e plus longue), puis \zh{丿} (piě) et \zh{乀} (zhé/ná) qui s'entrecroisent. \\
\hline
\zh{联} \emph{lián} & \zh{耳} (oreille) — 6 traits & 1 / 7 & Gauche \zh{耳} (haut $\rightarrow$ bas, cadre puis traits intérieurs), puis droite \zh{(trait brisé)} (crochet unique, forme simplifiée de \zh{关}). \\
\hline
\zh{发} \emph{fā} & \zh{又} (encore/main droite) — 2 traits & 3 / 5 & 1. \zh{丿} (piě). 2. \zh{㇇} (héng zhé tí / wān gōu - trait unique formant le haut). 3. \zh{丿} (piě, gauche de \zh{又}). 4. \zh{㇏} (nà, droite de \zh{又}). 5. \zh{一} (héng, base de \zh{又}). *Attention : \zh{发} (simplifié, cheveux/émettre) $\neq$ \zh{髮} (traditionnel, cheveux, clé \zh{髟}).* \\
\hline
\zh{视} \emph{shì} & \zh{礻} (\zh{示} simplifié, autel/voir) — 4 traits & 4 / 8 & Gauche \zh{礻} (点, 横, 点, 横) puis droite \zh{见} (见 : 竖, 横折, 横, 横). \textbf{Clé = \zh{礻} (shì), pas \zh{见}.} \\
\hline
\zh{频} \emph{pín} & \zh{页} (page/tête) — 6 traits & 7 / 13 & Haut \zh{品} (trois \zh{口} empilés : chaque \zh{口} = 3 traits : 竖, 横折, 横 ; ordre : gauche $\rightarrow$ droite, haut $\rightarrow$ bas) puis bas \zh{页} (横, 竖, 3 横, 竖 final). \\
\hline
\end{tabularx}
\end{center}
\end{definition}

\begin{methode}[Comment mémoriser l'écriture ?]
\begin{enumerate}
    \item \textbf{Décomposer :} Identifiez la clé (sens) et le phonétique (son). Ex: \zh{视} = \zh{礻} (sacré/voir) + \zh{见} (voir) $\rightarrow$ « regarder attentivement ».
    \item \textbf{Tracer dans l'air :} Suivez l'ordre strict (horizontal avant vertical, gauche avant droite, extérieur avant intérieur, fermer en dernier).
    \item \textbf{Écrire 5 fois :} Sur papier quadrillé (田字格 \emph{tiánzìgé}), en disant le pinyin à chaque trait.
    \item \textbf{Composer :} Écrivez le mot complet : \zh{视频}, \zh{发朋友圈}, \zh{上网}.
\end{enumerate}
\end{methode}

\begin{aretenir}[Checklist de fin de section « Glossaire »]
\begin{itemize}
    \item Je sais lire, écrire (pinyin + tons) et traduire les 17 mots du tableau.
    \item Je distingue \zh{网络} (réseau général) et \zh{互联网} (Internet mondial).
    \item Je sais utiliser \zh{发} + objet (照片, 视频, 朋友圈, 信息).
    \item Je maîtrise la structure du verbe séparable \zh{上网} : \zh{上了} + durée + \zh{网}.
    \item Je connais l'origine du mot \zh{网友} et la place centrale de \zh{微信} en Chine vs WhatsApp/Facebook au Cameroun.
    \item J'ai pratiqué la lecture à voix haute du texte de Léa et répondu aux questions.
\end{itemize}
\end{aretenir}

\coursec{Radicaux et ordre des traits}

Dans la section précédente, nous avons découvert le glossaire du thème « internet et les réseaux sociaux » : \zh{网络} (wǎngluò, réseau), \zh{聊天} (liáotiān, bavarder), \zh{影响} (yǐngxiǎng, influence), \zh{方便} (fāngbiàn, pratique), \zh{排队} (páiduì, faire la queue)... Ces caractères peuvent sembler complexes au premier regard. Bonne nouvelle : ils ne sont pas des dessins arbitraires ! Chaque caractère chinois est construit selon une logique précise, avec une \cle{clé} (ou \cle{radical}, \zh{部首} bùshǒu) qui donne une indication de sens, et souvent un \cle{composant phonétique} qui donne une indication de prononciation. Cette section vous apprend à « lire dans » les caractères et à les écrire dans le bon ordre.

\courssub{Pourquoi étudier les radicaux ?}

\begin{definition}[Le radical (部首 bùshǒu)]
Le \cle{radical} est la partie d'un caractère qui sert à le classer dans les dictionnaires et qui indique généralement sa \textbf{famille de sens}. Par exemple, les caractères contenant \zh{氵} (l'eau, forme compressée de \zh{水}) parlent souvent de liquides : \zh{河} (hé, rivière), \zh{洗} (xǐ, laver), \zh{海} (hǎi, mer).
\end{definition}

Observer le radical, c'est comme deviner le sens d'un mot français grâce à sa racine : quand vous voyez « aquatique », vous pensez à l'eau. En chinois, cette stratégie est encore plus puissante, car le radical est \emph{visible} dans le caractère.

\courssub{Étude des caractères du thème}

\textbf{1. 网 wǎng — réseau, filet (6 traits)}\par

\zh{网} est un caractère \emph{pictographique} : à l'origine, c'était le dessin d'un \textbf{filet} tendu pour pêcher ou attraper des oiseaux. Sa structure associe \zh{冂} (jiōng, le cadre extérieur) et \zh{爻} (yáo, les mailles croisées à l'intérieur). \textbf{Attention} : dans les dictionnaires modernes (chinois simplifié), sa clé de classement (radical dictionnaire) est \zh{纟} (sī, soie), car un filet est tissé ; mais pour l'analyse structurelle et la mémorisation, on retient le cadre \zh{冂}. Aujourd'hui, 网 sert à dire « réseau » : \zh{网络} (wǎngluò, réseau), \zh{上网} (shàngwǎng, aller sur internet), \zh{网友} (wǎngyǒu, ami rencontré en ligne). Internet est donc, en chinois, un immense filet qui relie le monde !

Ordre des traits (6 traits) : 
\begin{enumerate}
\item \zh{丨} Vertical gauche (bord gauche du cadre)
\item \zh{フ} Horizontal-crochet supérieur (bord haut + angle droit descendant)
\item \zh{乂} Premier croisement (maille gauche : trait oblique gauche \zh{丿} puis trait oblique droit \zh{丶} — ici écrit comme un petit trait croisé)
\item \zh{乂} Deuxième croisement (maille droite, même ordre)
\end{enumerate}
Règle générale respectée : \emph{de haut en bas, de gauche à droite, l'extérieur avant l'intérieur}.

\textbf{2. 聊 liáo — bavarder (11 traits)}\par

\zh{聊} est un caractère \emph{phono-sémantique} (形声字). Il se décompose en deux parties : à gauche \zh{耳} (ěr, l'oreille), qui est le \cle{radical de sens} (clé 耳), et à droite \zh{卯} (mǎo), le \cle{composant phonétique} qui rapproche la prononciation (liáo ≈ mǎo). Le sens est très logique : \textbf{bavarder, c'est prêter l'oreille} à quelqu'un ! D'où \zh{聊天} (liáotiān, bavarder, littéralement « discuter du ciel / passer le temps ») et \zh{聊天室} (liáotiānshì, salon de discussion).

Ordre des traits (11 traits) : on écrit d'abord \zh{耳} à gauche (6 traits : \zh{一} horizontal, \zh{丨} vertical, \zh{丨} vertical, \zh{一} horizontal, \zh{一} horizontal, \zh{一} horizontal), puis \zh{卯} à droite (5 traits : \zh{一} horizontal, \zh{丨} vertical, \zh{一} horizontal, \zh{一} horizontal, \zh{乚} crochet final). Règle : \emph{la partie gauche avant la partie droite}.

\begin{popfigure}
\begin{center}
\begin{tikzpicture}[
  charbox/.style={draw=popblue, thick, rounded corners=3pt, fill=popblueL, minimum width=1.4cm, minimum height=1.4cm, font=\Huge},
  partbox/.style={draw=poppurple, thick, rounded corners=3pt, fill=poppurpleL, minimum width=1.1cm, minimum height=1.1cm, font=\huge},
  lbl/.style={font=\small, align=center, text=popdark}
]
\node[charbox] (liao) at (0,0) {聊};
\node[partbox] (er) at (-2.6,-2.4) {耳};
\node[partbox] (mao) at (2.6,-2.4) {卯};
\node[lbl, above=0.05cm of er.north, yshift=0.15cm, align=center]{radical de sens : \zh{ěr} (oreille)\\(écouter)};
\node[lbl, above=0.05cm of mao.north, yshift=0.15cm, align=center]{composant phonétique : \zh{mǎo}\\(son : liáo)};
\draw[-{Stealth}, thick, poppurple] (liao.south west) -- (er.north);
\draw[-{Stealth}, thick, poppurple] (liao.south east) -- (mao.north);
\node[lbl, below=0.2cm of er.south] {\zh{ěr} : oreille};
\node[lbl, below=0.2cm of mao.south] {\zh{mǎo} : indique le son};
\end{tikzpicture}
\legende{Décomposition du caractère \zh{聊} (liáo, bavarder) : radical de sens + composant phonétique.}
\end{center}
\end{popfigure}

\textbf{3. 影 yǐng — ombre, image, reflet (15 traits)}\par

\zh{影} est un caractère \emph{phono-sémantique}. Il se décompose en \zh{景} (jǐng, paysage, lumière du jour) à gauche — qui donne aussi une indication phonétique (jǐng ≈ yǐng) — et \zh{彡} (shān) à droite, le \cle{radical}, qui représente des \emph{traits de lumière, des poils ou des ornements} — ici, il évoque l'ombre portée ou le reflet. 影 apparaît dans \zh{电影} (diànyǐng, cinéma, litt. « ombre électrique »), \zh{照片} (zhàopiàn, photo, via \zh{照} qui partage le même radical \zh{彡}) et surtout \zh{影响} (yǐngxiǎng, influence).

Ordre des traits (15 traits) : on écrit \textbf{d'abord \zh{景}} (12 traits, de haut en bas : \zh{日} soleil 4 traits, puis \zh{京} capitale 8 traits — \zh{亠} \zh{口} \zh{小}), \textbf{puis \zh{彡} à droite} (3 traits : trois courbes \zh{丿} de haut en bas). Règle respectée : gauche avant droite, haut avant bas.

\textbf{4. 响 xiǎng — résonner, bruit, son (9 traits)}\par

\zh{响} est un caractère \emph{phono-sémantique}. Son radical est \zh{口} (kǒu, la bouche) à gauche : logique, car un son est produit par la bouche/la voix ! La partie droite \zh{向} (xiàng, vers, direction) donne le son (xiǎng ≈ xiàng). \zh{响} signifie « résonner, faire du bruit » : \zh{音响} (yīnxiǎng, sono, chaîne stéréo), \zh{影响} (yǐngxiǎng, influence — littéralement « l'ombre et l'écho », belle image : l'influence est l'ombre et l'écho que l'on laisse !).

Ordre des traits (9 traits) : \zh{口} à gauche (3 traits : \zh{丨} vertical, \zh{フ} horizontal-crochet, \zh{一} horizontal de fermeture), puis \zh{向} à droite (6 traits : \zh{丶} point, \zh{丿} oblique, \zh{一} horizontal, \zh{丨} vertical, \zh{フ} horizontal-crochet, \zh{一} horizontal).

\textbf{5. 便 biàn — pratique, commode (9 traits)}\par

\zh{便} est un caractère \emph{phono-sémantique}. Il porte le radical \zh{亻} (rén, forme compressée de \zh{人}, personne), toujours placée à gauche. La partie droite \zh{更} (gèng, changer, encore) est le composant phonétique (biàn ≈ gèng) et suggère l'idée de « changement » vers quelque chose de plus aisé. On le trouve dans \zh{方便} (fāngbiàn, pratique, commode) : « Internet est très pratique » se dit \zh{网络很方便。} (Wǎngluò hěn fāngbiàn.)

Ordre des traits (9 traits) : d'abord \zh{亻} (2 traits : \zh{丿} courbe descendante, \zh{丨} vertical), puis \zh{更} (7 traits, de haut en bas : \zh{一} \zh{丨} \zh{一} \zh{一} \zh{丨} \zh{一} \zh{乚}).

\textbf{6. 排 pái — ranger, aligner, éliminer (11 traits)}\par

\zh{排} est un caractère \emph{phono-sémantique}. Il porte le radical \zh{扌} (shǒu, forme compressée de \zh{手}, la main). C'est un indice précieux : les caractères avec \zh{扌} sont presque toujours des \textbf{verbes d'action manuelle} : \zh{打} (dǎ, frapper, taper), \zh{找} (zhǎo, chercher), \zh{拍} (pāi, frapper, photographier), \zh{排} (pái, ranger, aligner). On dit \zh{排队} (páiduì, faire la queue) et \zh{安排} (ānpái, organiser, arranger). La partie droite \zh{非} (fēi, non) est le composant phonétique (pái ≈ fēi).

Ordre des traits (11 traits) : \zh{扌} à gauche (3 traits : \zh{一} horizontal, \zh{亅} vertical-crochet, \zh{㇀} courbe montante / trait relevé), puis \zh{非} à droite (8 traits : colonne de gauche \zh{一}\zh{丨}\zh{一}\zh{一}, colonne de droite \zh{一}\zh{丨}\zh{一}\zh{一}).

\courssub{Tableau récapitulatif}

\begin{center}
\begin{tabularx}{\linewidth}{|X|X|X|X|}
\hline
\rowcolor{popblueL} Caractère & Pinyin et sens & Radical (clé) & Nombre de traits \\
\hline
网 & wǎng : réseau, filet & 冂 (cadre) / 纟 (dictionnaire) & 6 \\
\hline
聊 & liáo : bavarder & 耳 (oreille) & 11 \\
\hline
影 & yǐng : ombre, image & 彡 (traits de lumière) & 15 \\
\hline
响 & xiǎng : résonner, son & 口 (bouche) & 9 \\
\hline
便 & biàn : pratique & 亻 (personne) & 9 \\
\hline
排 & pái : ranger, aligner & 扌 (main) & 11 \\
\hline
\end{tabularx}
\end{center}

\begin{propriete}[Le radical 亻 (rén), la clé de la personne]
Le radical \zh{亻} (forme gauche de \zh{人}, personne) apparaît dans de très nombreux caractères courants. Il signale très souvent un lien avec l'\textbf{être humain}, ses actions, ses qualités ou ses relations.
\begin{itemize}
\item \zh{信} (xìn) : personne + parole → confiance, lettre, message.
\item \zh{休} (xiū) : personne + arbre → se reposer (à l'ombre de l'arbre).
\item \zh{们} (men) : personne + porte → groupe, pluriel (\zh{我们} wǒmen, \zh{同学们} tóngxuémen).
\item \zh{便} (biàn) : personne + changer → pratique (ce qui arrange la personne).
\item \zh{做} (zuò) : personne + faire → faire, agir.
\end{itemize}
Quand vous rencontrez un caractère inconnu avec \zh{亻}, pensez d'abord : « Cela concerne probablement un être humain ».
\end{propriete}

\begin{exemplebox}[Deux mots du thème éclairés par leurs radicaux]
\zh{聊天} (liáotiān, bavarder) : \zh{耳} = l'oreille → bavarder, c'est \emph{écouter} et parler tour à tour.\\
\zh{影响} (yǐngxiǎng, influence) : \zh{彡} évoque l'\emph{ombre/reflet}, \zh{口} la \emph{voix/écho} → l'influence est l'ombre et l'écho que l'on laisse sur les autres. Retenir ces images mentales rend les mots inoubliables !
\end{exemplebox}

\begin{methode}[Mémoriser un caractère en trois étapes]
\begin{enumerate}
\item \textbf{Identifier le radical (clé de sens)} : il donne la famille de sens (ex. \zh{扌} → action de la main ; \zh{亻} → personne).
\item \textbf{Identifier le composant phonétique} : il rapproche la prononciation (ex. \zh{卯} dans \zh{聊} liáo ; \zh{非} dans \zh{排} pái).
\item \textbf{Écrire le caractère 5 à 10 fois} en comptant les traits à voix haute, dans le bon ordre (haut → bas, gauche → droite, extérieur → intérieur, fermer en dernier).
\end{enumerate}
\end{methode}

\begin{attention}[Caractères proches à ne pas confondre]
Ne confondez pas :
\begin{itemize}
\item \zh{网} (wǎng, réseau / filet) et \zh{同} (tóng, pareil, même) : \zh{网} contient deux croix \zh{爻} (les mailles) dans un cadre ; \zh{同} contient \zh{一} et \zh{口} sous un toit.
\item \zh{便} (biàn, pratique) et \zh{使} (shǐ, utiliser, faire, envoyer) : \zh{便} contient \zh{更} (gèng) à droite ; \zh{使} contient \zh{吏} (lì, fonctionnaire) à droite. Une seule différence de trait change tout le sens : vérifiez toujours le composant de droite !
\item \zh{排} (pái, ranger) et \zh{非} (fēi, non / pas) : \zh{排} a le radical main \zh{扌} à gauche, \zh{非} n'en a pas.
\end{itemize}
\end{attention}

\begin{exoresolu}[Deviner le sens grâce au radical]
Énoncé — Sans dictionnaire, devinez la famille de sens des caractères suivants en observant leur radical : \zh{信} (xìn), \zh{休} (xiū), \zh{们} (men).
\tcblower
Solution détaillée — Les trois caractères portent le radical \zh{亻} (personne) à gauche. Ils concernent donc tous l'être humain.
\begin{itemize}
\item \zh{信} (xìn) : \zh{亻} (personne) + \zh{言} (yán, parole) → la parole de l'homme en qui l'on a confiance → \emph{croire, confiance, lettre, message}.
\item \zh{休} (xiū) : \zh{亻} (personne) + \zh{木} (mù, arbre) → l'homme s'appuie contre un arbre pour souffler → \emph{se reposer, pause}.
\item \zh{们} (men) : \zh{亻} (personne) + \zh{门} (mén, porte) → phonétique \zh{mén} + sens « groupe de gens » (autrefois à la porte du village) → \emph{marque du pluriel} (\zh{我们} wǒmen, nous ; \zh{老师们} lǎoshīmen, les professeurs).
\end{itemize}
Le radical a permis de deviner la logique commune des trois caractères !
\end{exoresolu}

\begin{exercice}[Entraînement à l'écriture et à la reconnaissance]
\begin{enumerate}
\item Dans votre cahier, écrivez chacun des caractères suivants \textbf{cinq fois} en comptant les traits à voix haute et en respectant l'ordre appris : \zh{网} (6 traits), \zh{聊} (11 traits), \zh{便} (9 traits), \zh{排} (11 traits).
\item Recopiez la phrase suivante en caractères, pinyin (avec tons) et traduction française :\\
\zh{我们在网上聊天很方便。} (Wǒmen zài wǎngshang liáotiān hěn fāngbiàn. — « Nous bavardons sur internet, c'est très pratique. »)
\item Complétez avec le bon caractère (\zh{响} ou \zh{影} ; \zh{便} ou \zh{使}) :\\
a) 电\_\_ (diàn\_\_, cinéma) \quad b) 音\_\_ (yīn\_\_, sono) \quad c) 方\_\_ (fāng\_\_, pratique) \quad d) \_\_用 (\_\_yòng, utiliser)
\end{enumerate}
\end{exercice}

\begin{corrige}[Exercice : Entraînement à l'écriture et reconnaissance]
\begin{enumerate}
\item \textbf{Ordre des traits à vérifier par le professeur} :\\
\zh{网} : 1) \zh{丨} vertical gauche, 2) \zh{フ} horizontal-crochet sup., 3) \zh{乂} maille g., 4) \zh{乂} maille d. (ext. avant int.).\\
\zh{聊} : \zh{耳} complet à gauche (6 traits), puis \zh{卯} à droite (5 traits).\\
\zh{便} : \zh{亻} (2 traits), puis \zh{更} (7 traits, haut → bas).\\
\zh{排} : \zh{扌} (3 traits), puis \zh{非} (8 traits, colonne g. puis colonne d.).
\item \zh{我们在网上聊天很方便。} (Wǒmen zài wǎngshang liáotiān hěn fāngbiàn. — « Nous bavardons sur internet, c'est très pratique. »)
\item a) \zh{电影} (diànyǐng) — radical \zh{彡} (ombre/image).\\
b) \zh{音响} (yīnxiǎng) — radical \zh{口} (son/bouche).\\
c) \zh{方便} (fāngbiàn) — radical \zh{亻} (personne).\\
d) \zh{使用} (shǐyòng) — radical \zh{亻} (personne qui agit/emploie).
\end{enumerate}
\end{corrige}

\begin{aretenir}[À retenir de cette section]
Chaque caractère chinois se lit comme une petite histoire : le \cle{radical} (clé de sens) donne la famille de sens (\zh{耳} l'oreille pour \zh{聊}, \zh{扌} la main pour \zh{排}, \zh{亻} la personne pour \zh{便}, \zh{口} la bouche pour \zh{响}, \zh{彡} l'ombre pour \zh{影}), le \cle{composant phonétique} rapproche le son. L'écriture suit toujours les mêmes règles universelles : de haut en bas, de gauche à droite, de l'extérieur vers l'intérieur (fermer en dernier). Observer les radicaux, c'est apprendre deux fois plus vite et écrire correctement !
\end{aretenir}

\coursec{Leçon 34 — Vocabulaire et compréhension de texte : internet et les réseaux sociaux}

\courssub{Mise en route et découverte du thème}

\begin{experience}[Discussion d'amorce : nos habitudes numériques]
En classe entière, le professeur lance la discussion en français, puis en chinois simple : \zh{你上网吗？} \emph{Nǐ shàngwǎng ma?} (Tu te connectes ?), \zh{你用手机还是电脑？} \emph{Nǐ yòng shǒujī háishì diànnǎo?} (Téléphone ou ordinateur ?), \zh{你用微信吗？} \emph{Nǐ yòng Wēixìn ma?} (Tu utilises WeChat ?). Les élèves répondent par \zh{是} / \zh{不是} ou \zh{我用……} \emph{Wǒ yòng…}. On note au tableau les mots-clés qui émergent (\zh{手机}, \zh{微信}, \zh{聊天}, \zh{看视频}) pour introduire le vocabulaire de la leçon.
\end{experience}

\courssub{Texte support : lecture et compréhension globale}

\begin{textebox}[网络与我们的生活 — Internet et notre vie]
现在，网络对我们的生活非常重要。年轻人一般用手机上网。\\
\emph{Xiànzài, wǎngluò duì wǒmen de shēnghuó fēicháng zhòngyào. Niánqīngrén yìbān yòng shǒujī shàngwǎng.}\\
« Aujourd'hui, internet est très important dans notre vie. Les jeunes se connectent généralement avec leur téléphone portable. »\\[4pt]
他们在微信上和朋友们聊天，也在网上看视频、买东西。\\
\emph{Tāmen zài Wēixìn shàng hé péngyoumen liáotiān, yě zài wǎngshàng kàn shìpín, mǎi dōngxi.}\\
« Ils discutent avec leurs amis sur WeChat, et regardent aussi des vidéos et achètent des choses en ligne. »\\[4pt]
网络让我们的生活更方便，还让我们认识世界各地的朋友。\\
\emph{Wǎngluò ràng wǒmen de shēnghuó gèng fāngbiàn, hái ràng wǒmen rènshi shìjiè gèdì de péngyou.}\\
« Internet rend notre vie plus pratique et nous permet aussi de connaître des amis du monde entier. »\\[4pt]
但是，如果花太多时间上网，对身体不好，也会影响学习。\\
\emph{Dànshì, rúguǒ huā tài duō shíjiān shàngwǎng, duì shēntǐ bù hǎo, yě huì yǐngxiǎng xuéxí.}\\
« Mais si l'on passe trop de temps en ligne, c'est mauvais pour la santé et cela nuit aussi aux études. »\\[4pt]
所以，我们应该合理安排上网时间。\\
\emph{Suǒyǐ, wǒmen yīnggāi hélǐ ānpái shàngwǎng shíjiān.}\\
« Nous devrions donc organiser raisonnablement notre temps de connexion. »
\end{textebox}

\begin{methode}[Lecture active en trois temps]
\begin{enumerate}
\item \textbf{Lecture à voix haute} par le professeur (modèle de tons, de liaison et d'intonation), puis par les élèves (chorale, puis individuelle).
\item \textbf{Repérage visuel} : surligner les mots connus, encadrer les nouveaux caractères, cercler les mots-outils (\zh{的}, \zh{了}, \zh{也}, \zh{如果}, \zh{所以}, \zh{应该}).
\item \textbf{Traduction segmentée} : phrase par phrase, du chinois vers le français, sans traduction mot à mot, pour saisir le sens global.
\end{enumerate}
\end{methode}

\courssub{Glossaire du thème (Vocabulaire)}

\begin{center}
\begin{tabularx}{\linewidth}{|X|X|X|X|}
\hline
\rowcolor{popblueL} Caractères & Pinyin & Nature & Traduction \\
\hline
网络 & wǎngluò & nom & internet, réseau, toile \\
\hline
重要 & zhòngyào & adj. & important \\
\hline
年轻人 & niánqīngrén & nom & jeune, jeune personne \\
\hline
手机 & shǒujī & nom & téléphone portable, smartphone \\
\hline
上网 & shàngwǎng & verbe & se connecter, aller sur internet \\
\hline
微信 & Wēixìn & nom propre & WeChat (appli messagerie/paiement) \\
\hline
聊天 & liáotiān & verbe & discuter, bavarder (en ligne ou oral) \\
\hline
视频 & shìpín & nom & vidéo \\
\hline
买东西 & mǎi dōngxi & loc. verbale & faire des achats, acheter des choses \\
\hline
方便 & fāngbiàn & adj. & pratique, commode \\
\hline
认识 & rènshi & verbe & connaître (personne), reconnaître \\
\hline
世界各地 & shìjiè gèdì & loc. nominale & le monde entier, les quatre coins du monde \\
\hline
花 (时间) & huā (shíjiān) & verbe & passer (du temps), dépenser (argent) \\
\hline
身体 & shēntǐ & nom & corps, santé physique \\
\hline
影响 & yǐngxiǎng & verbe/nom & influencer, nuire à / influence, impact \\
\hline
学习 & xuéxí & verbe/nom & étudier, apprendre / études \\
\hline
所以 & suǒyǐ & conjonction & donc, par conséquent \\
\hline
合理 & hélǐ & adj. & raisonnable, rationnel \\
\hline
安排 & ānpái & verbe & organiser, arranger, planifier \\
\hline
如果 & rúguǒ & conjonction & si (condition) \\
\hline
因为 / 所以 & yīnwèi / suǒyǐ & conjonctions & parce que / donc (corrélatives) \\
\hline
\end{tabularx}
\end{center}

\begin{aretenir}[Mots-outils clés du texte]
\begin{itemize}
\item \zh{也} (yě) « aussi, également » : se place \textbf{avant} le verbe (\zh{也在网上看}).
\item \zh{还} (hái) « en plus, encore » : insiste sur l'addition d'un deuxième avantage (\zh{更方便，还让我们认识…}).
\item \zh{如果}… (rúguǒ) « si… » : introduit la condition ; la conséquence suit souvent avec \zh{就} (jiù) (omissible à l'oral) ou directement le verbe (\zh{如果花太多时间…, 对身体不好}).
\item \zh{应该} (yīnggāi) « devoir, devrait » : modalité de conseil/obligation morale, placé avant le verbe (\zh{应该合理安排}).
\item \zh{对}…\zh{不好} (duì… bù hǎo) « mauvais pour… » : structure fixe, \zh{对} introduit l'objet affecté.
\end{itemize}
\end{aretenir}

\courssub{Radicaux et ordre des traits}

\begin{definition}[Clés (radicaux) et composition des caractères nouveaux]
Le tableau ci-dessous présente l'analyse graphique des caractères principaux du thème. La connaissance des clés facilite la mémorisation, la recherche dans le dictionnaire et la compréhension du sens.
\end{definition}

\begin{center}
\begin{tabularx}{\linewidth}{|X|X|X|X|X|}
\hline
\rowcolor{popblueL} Caractère & Pinyin & Clé (Radical) & Nombre de traits & Structure / Ordre des traits principal \\
\hline
网 & wǎng & \zh{纟} (sī, soie) & 6 & Structure haut-bas : écrire \zh{纟} (3 traits : frémissement, frémissement, verticale) puis \zh{亡} (wáng, 3 traits : horizontale, horizontale, crochet vertical). Sens : filet, réseau. \\
\hline
络 & luò & \zh{纟} (sī, soie) & 11 & Structure gauche-droite : \zh{纟} puis \zh{各} (gè). \zh{各} : \zh{口} (bouche) + \zh{夂} (pied qui traîne). Sens : filet, maillage, réseau. \\
\hline
信 & xìn & \zh{亻} (rén, homme) & 9 & Structure gauche-droite : \zh{亻} (2 traits) puis \zh{言} (yán, parole, 7 traits : point, horizontale, horizontale, horizontale, verticale, horizontale, point). Sens : confiance, lettre, message,信息. \\
\hline
聊 & liáo & \zh{讠} (yán, parole) & 11 & Structure gauche-droite : \zh{讠} (2 traits : point, courbe) puis \zh{尞} (liáo, phonétique). Sens : bavarder, discuter. \\
\hline
微 & wēi & \zh{彳} (chì, pas) & 13 & Structure gauche-droite : \zh{彳} (3 traits) puis \zh{徽} (huī). Sens : minuscule,微弱, 微信. \\
\hline
视 & shì & \zh{示} (shì, autel/montrer) & 11 & Structure haut-bas : \zh{示} (5 traits) puis \zh{见} (jiàn, voir, 4 traits). Sens : regarder, vision, 视频. \\
\hline
频 & pín & \zh{页} (yè, page/tête) & 15 & Structure haut-bas : \zh{页} (6 traits) puis \zh{步} (bù, pas, 7 traits). Sens : fréquence, répété, 视频. \\
\hline
花 & huā & \zh{艹} (cǎo, herbe) & 7 & Structure haut-bas : \zh{艹} (3 traits) puis \zh{化} (huà, 4 traits). Sens : fleur, passer (temps/argent). \\
\hline
身 & shēn & \zh{身} (shēn, corps) & 7 & Caractère-clé. Traits : horizontale, verticale, horizontale, horizontale, verticale, horizontale, horizontale. Sens : corps, personne. \\
\hline
体 & tǐ & \zh{骨} (gǔ, os) & 7 & Structure gauche-droite : \zh{骨} (simplifié en partie haute) + \zh{亻} ? Non, \zh{体} simplifié : \zh{亻} + \zh{本} (běn, racine). 7 traits. Sens : corps, corps physique. \\
\hline
\end{tabularx}
\end{center}

\begin{attention}[Pièges graphiques fréquents]
\begin{itemize}
\item \zh{网} vs \zh{罗} (luó, filet / rassembler) : même clé \zh{纟}, phonétique différent (\zh{亡} vs \zh{罗}).
\item \zh{信} : ne pas confondre \zh{言} (parole, 7 traits) et \zh{讠} (clé parole à gauche, 2 traits). \zh{信} s'écrit avec \zh{言} à droite.
\item \zh{视} : le bas est \zh{见} (voir), pas \zh{贝} (coquillage/argent) ni \zh{见} mal formé (le trait vertical central ne traverse pas la première horizontale).
\item \zh{体} (simplifié) vs \zh{體} (traditionnel) : en simplifié, c'est \zh{亻} + \zh{本}. Ne pas écrire \zh{骨} à gauche.
\end{itemize}
\end{attention}

\courssub{Questions de compréhension du texte}

Après avoir étudié les radicaux et l'ordre des traits des caractères du thème (\zh{网}, \zh{络}, \zh{信}, \zh{聊}, \zh{微}, \zh{视}, \zh{频}, \zh{花}, \zh{身}, \zh{体}…), nous revenons maintenant au texte support pour vérifier notre compréhension. Savoir répondre à une question de compréhension en chinois est une compétence clé du programme de Terminale : il ne suffit pas de comprendre, il faut aussi \cle{formuler une réponse complète et correcte}.

\courssub{Rappel du texte support}

Relisons le texte étudié dans la section 2 avant de répondre aux questions.

\begin{textebox}[网络与我们的生活 — Internet et notre vie]
现在，网络对我们的生活非常重要。年轻人一般用手机上网。\\
\emph{Xiànzài, wǎngluò duì wǒmen de shēnghuó fēicháng zhòngyào. Niánqīngrén yìbān yòng shǒujī shàngwǎng.}\\
« Aujourd'hui, internet est très important dans notre vie. Les jeunes se connectent généralement avec leur téléphone portable. »\\[4pt]
他们在微信上和朋友们聊天，也在网上看视频、买东西。\\
\emph{Tāmen zài Wēixìn shàng hé péngyoumen liáotiān, yě zài wǎngshàng kàn shìpín, mǎi dōngxi.}\\
« Ils discutent avec leurs amis sur WeChat, et regardent aussi des vidéos et achètent des choses en ligne. »\\[4pt]
网络让我们的生活更方便，还让我们认识世界各地的朋友。\\
\emph{Wǎngluò ràng wǒmen de shēnghuó gèng fāngbiàn, hái ràng wǒmen rènshi shìjiè gèdì de péngyou.}\\
« Internet rend notre vie plus pratique et nous permet aussi de connaître des amis du monde entier. »\\[4pt]
但是，如果花太多时间上网，对身体不好，也会影响学习。\\
\emph{Dànshì, rúguǒ huā tài duō shíjiān shàngwǎng, duì shēntǐ bù hǎo, yě huì yǐngxiǎng xuéxí.}\\
« Mais si l'on passe trop de temps en ligne, c'est mauvais pour la santé et cela nuit aussi aux études. »\\[4pt]
所以，我们应该合理安排上网时间。\\
\emph{Suǒyǐ, wǒmen yīnggāi hélǐ ānpái shàngwǎng shíjiān.}\\
« Nous devrions donc organiser raisonnablement notre temps de connexion. »
\end{textebox}

\courssub{Méthode : répondre à une question de compréhension}

\begin{methode}[Répondre en trois étapes]
\begin{enumerate}
\item \textbf{Repérer le mot interrogatif.} Chaque mot interrogatif appelle un type de réponse précis : \zh{什么} (quoi) → une chose ; \zh{谁} (qui) → une personne ; \zh{为什么} (pourquoi) → une cause introduite par \zh{因为} ; \zh{怎么} (comment) → une manière de faire ; \zh{哪儿} (où) → un lieu.
\item \textbf{Localiser la phrase du texte.} Souligner dans le texte la phrase qui contient l'information demandée. Les mots de la question (par exemple \zh{上网}, \zh{微信}) servent d'indices.
\item \textbf{Recopier et adapter en phrase complète.} Reprendre la structure de la phrase du texte en l'adaptant à la question. Une réponse en un seul mot est une réponse incomplète : elle sera pénalisée à l'examen.
\end{enumerate}
\end{methode}

\begin{center}
\begin{tabularx}{\linewidth}{|X|X|X|X|}
\hline
\rowcolor{popblueL} Mot interrogatif & Pinyin & Type de réponse & Début de réponse typique \\
\hline
什么 & shénme & une chose, une action & Sujet + verbe + complément \\
\hline
谁 & shéi & une personne & 是 + personne \\
\hline
为什么 & wèishénme & une cause & 因为… \\
\hline
怎么 & zěnme & une manière & 应该 / 可以 + verbe \\
\hline
哪儿 & nǎr & un lieu & 在 + lieu \\
\hline
\end{tabularx}
\end{center}

\courssub{Les cinq questions de compréhension}

\textbf{Question 1 :} \zh{年轻人用什么上网？} \emph{Niánqīngrén yòng shénme shàngwǎng?} « Que les jeunes utilisent-ils pour se connecter ? »

Le mot interrogatif est \zh{什么} : on cherche une chose. La phrase du texte est « \zh{年轻人一般用手机上网} ». Réponse complète attendue :

\zh{年轻人一般用手机上网。} \emph{Niánqīngrén yìbān yòng shǒujī shàngwǎng.} « Les jeunes se connectent généralement avec leur téléphone portable. »

Réponse insuffisante (à éviter) : \zh{手机。} — un seul mot ne montre pas que l'élève a compris la structure de la phrase.

\textbf{Question 2 :} \zh{他们在微信上做什么？} \emph{Tāmen zài Wēixìn shàng zuò shénme?} « Que font-ils sur WeChat ? »

On cherche une action. Le texte dit « \zh{他们在微信上和朋友们聊天} ». Réponse complète :

\zh{他们在微信上和朋友们聊天。} \emph{Tāmen zài Wēixìn shàng hé péngyoumen liáotiān.} « Ils discutent avec leurs amis sur WeChat. »

\textbf{Question 3 :} \zh{网络给我们的生活带来什么好处？} \emph{Wǎngluò gěi wǒmen de shēnghuó dài lái shénme hǎo chù?} « Quels avantages internet apporte-t-il à notre vie ? »

Le texte donne deux avantages : \zh{更方便} et \zh{认识世界各地的朋友}. Une bonne réponse cite les deux :

\zh{网络让我们的生活更方便，还让我们认识世界各地的朋友。} \emph{Wǎngluò ràng wǒmen de shēnghuó gèng fāngbiàn, hái ràng wǒmen rènshi shìjiè gèdì de péngyou.} « Internet rend notre vie plus pratique et nous permet de connaître des amis du monde entier. »

\textbf{Question 4 :} \zh{花太多时间上网有什么问题？} \emph{Huā tài duō shíjiān shàngwǎng yǒu shénme wèntí?} « Quels problèmes pose le fait de passer trop de temps en ligne ? »

Réponse complète :

\zh{花太多时间上网对身体不好，也会影响学习。} \emph{Huā tài duō shíjiān shàngwǎng duì shēntǐ bù hǎo, yě huì yǐngxiǎng xuéxí.} « Passer trop de temps en ligne est mauvais pour la santé et nuit aussi aux études. »

\textbf{Question 5 :} \zh{作者建议我们怎么做？} \emph{Zuòzhě jiànyì wǒmen zěnme zuò?} « Que l'auteur nous conseille-t-il de faire ? »

Le mot interrogatif est \zh{怎么} : on cherche une manière d'agir, introduite par \zh{应该}. Réponse complète :

\zh{作者建议我们应该合理安排上网时间。} \emph{Zuòzhě jiànyì wǒmen yīnggāi hélǐ ānpái shàngwǎng shíjiān.} « L'auteur nous conseille d'organiser raisonnablement notre temps de connexion. »

\courssub{La question « pourquoi » : la paire 因为…所以…}

\begin{exemplebox}[Répondre à une question en 为什么]
Question : \zh{为什么网络让我们的生活更方便？} \emph{Wèishénme wǎngluò ràng wǒmen de shēnghuó gèng fāngbiàn?} « Pourquoi internet rend-il notre vie plus pratique ? »\\[4pt]
Réponse : \zh{因为我们能在网上看视频、买东西。} \emph{Yīnwèi wǒmen néng zài wǎngshàng kàn shìpín, mǎi dōngxi.} « Parce que nous pouvons regarder des vidéos et acheter des choses en ligne. »\\[4pt]
Réponse complète avec conclusion : \zh{因为我们能在网上看视频、买东西，所以网络让我们的生活更方便。} \emph{Yīnwèi wǒmen néng zài wǎngshàng kàn shìpín, mǎi dōngxi, suǒyǐ wǎngluò ràng wǒmen de shēnghuó gèng fāngbiàn.} « Parce que nous pouvons regarder des vidéos et acheter des choses en ligne, internet rend donc notre vie plus pratique. »
\end{exemplebox}

\begin{attention}[Les pièges des francophones avec 因为 et 所以]
\begin{itemize}
\item En chinois, la réponse à \zh{为什么} commence obligatoirement par \zh{因为} (parce que). Ne jamais répondre directement par le verbe comme en français oral (« Pourquoi ? — On peut acheter en ligne »).
\item Contrairement au français, le chinois \emph{autorise} et même \emph{apprécie} la combinaison \zh{因为…所以…} dans la même phrase : \zh{因为…，所以…} « parce que…, donc… ». En français, « parce que… donc… » serait une faute de syntaxe ; en chinois, c'est la structure correcte et naturelle.
\item Attention à la graphie de \zh{所以} : ne pas oublier le caractère \zh{所} (clé \zh{戶} hù, la porte). Écrire seulement \zh{以} (yǐ) est une erreur fréquente qui change le sens (\zh{以} = par, selon, pour).
\item Ponctuation : dans la structure \zh{因为…，所以…}, une virgule sépare généralement les deux propositions.
\end{itemize}
\end{attention}

\courssub{La structure du texte en organigramme}

Comprendre un texte, c'est aussi comprendre son organisation. Le texte suit un plan en cinq étapes, très fréquent dans les textes argumentatifs chinois (structure \cle{总分总} ou \cle{现象-分析-建议}) :

\begin{popfigure}
\begin{tikzpicture}[node distance=0.55cm,
  boite/.style={rectangle, rounded corners=3pt, draw=popblue, fill=popblueL, align=center, text width=4.2cm, minimum height=0.9cm, font=\small},
  boitefin/.style={rectangle, rounded corners=3pt, draw=popgreen, fill=popgreenL, align=center, text width=4.2cm, minimum height=0.9cm, font=\small},
  fleche/.style={-{Stealth[length=2.5mm]}, thick, popdark}]
\node[boite, align=center] (a){网络非常重要\\ \emph{Importance d'internet}};
\node[boite, below=of a, align=center] (b){上网和微信的用法\\ \emph{Les usages concrets}};
\node[boite, below=of b, align=center] (c){好处：方便、交朋友\\ \emph{Les avantages}};
\node[boite, below=of c, align=center] (d){Problèmes : santé, études\\ \emph{Les dangers / inconvénients}};
\node[boitefin, below=of d, align=center] (e){建议：合理安排时间\\ \emph{Le conseil / Solution}};
\draw[fleche] (a) -- (b);
\draw[fleche] (b) -- (c);
\draw[fleche] (c) -- (d);
\draw[fleche] (d) -- (e);
\end{tikzpicture}
\legende{Structure argumentative du texte : de l'importance d'internet au conseil final.}
\end{popfigure}

\courssub{Exercice résolu et reformulation}

\begin{exoresolu}[Répondre par une phrase complète]
Énoncé : répondez par une phrase complète en chinois à la question suivante : \zh{年轻人为什么喜欢微信？} \emph{Niánqīngrén wèishénme xǐhuan Wēixìn?} « Pourquoi les jeunes aiment-ils WeChat ? »
\tcblower
Solution détaillée :
\begin{enumerate}
\item Mot interrogatif : \zh{为什么} → la réponse doit commencer par \zh{因为}.
\item Information dans le texte : \zh{他们在微信上和朋友们聊天} (ils discutent avec leurs amis sur WeChat).
\item Réponse complète : \zh{因为他们能在微信上和朋友们聊天，所以年轻人喜欢微信。} \emph{Yīnwèi tāmen néng zài Wēixìn shàng hé péngyoumen liáotiān, suǒyǐ niánqīngrén xǐhuan Wēixìn.} « Parce qu'ils peuvent discuter avec leurs amis sur WeChat, les jeunes aiment donc WeChat. »
\end{enumerate}
La réponse \zh{因为聊天。} est trop courte et grammaticalement pauvre : elle serait sanctionnée.
\end{exoresolu}

\begin{methode}[Résumer un texte en 2-3 phrases avec ses propres mots]
\begin{enumerate}
\item Identifier l'idée de chaque partie du plan (voir l'organigramme ci-dessus).
\item Remplacer les expressions du texte par des synonymes simples : \zh{非常重要} → \zh{很有用} (très utile) ; \zh{花太多时间} → \zh{上网时间太长} (temps de connexion trop long) ; \zh{合理安排} → \zh{好好安排} (bien organiser).
\item Assembler en 2 ou 3 phrases avec des connecteurs logiques : \zh{但是} (mais), \zh{所以} (donc), \zh{因为} (parce que).
\end{enumerate}
Modèle de résumé : \zh{网络很有用，年轻人用它聊天、买东西。但是上网时间太长对身体和学习不好，所以我们应该好好安排上网时间。} \emph{Wǎngluò hěn yǒuyòng, niánqīngrén yòng tā liáotiān, mǎi dōngxi. Dànshì shàngwǎng shíjiān tài cháng duì shēntǐ hé xuéxí bù hǎo, suǒyǐ wǒmen yīnggāi hǎo hǎo ānpái shàngwǎng shíjiān.} « Internet est très utile, les jeunes s'en servent pour discuter et acheter des choses. Mais passer trop de temps en ligne est mauvais pour la santé et les études, nous devrions donc bien organiser notre temps de connexion. »
\end{methode}

\courssub{Production guidée et comparaison Cameroun — Chine}

\begin{experience}[Débat oral : Internet à Yaoundé/Douala vs Pékin/Shanghai]
En binômes, préparez un dialogue de 2 minutes. L'élève A est un lycéen camerounais, l'élève B un lycéen chinois. Utilisez le vocabulaire de la leçon.
\textbf{Pistes :}
\begin{itemize}
\item Applications : \zh{微信} (WeChat) vs \zh{WhatsApp} / \zh{Facebook} / \zh{TikTok} (\zh{抖音} Dǒuyīn en Chine).
\item Paiement : \zh{微信支付} / \zh{支付宝} (Alipay) — omniprésents en Chine (marchés, taxis) vs \zh{Mobile Money} (MTN MoMo, Orange Money) — dominant au Cameroun.
\item Accès : \zh{流量} (données 4G/5G) pas cher en Chine vs coût relatif au Cameroun ; \zh{防火长城} (Grand Firewall) en Chine (pas Google, pas YouTube, pas Facebook) vs internet ouvert au Cameroun.
\item Usage scolaire : \zh{网课} (cours en ligne) — généralisés en Chine (post-Covid) ; développement au Cameroun.
\end{itemize}
\textbf{Modèle de début :} \zh{A: 你上网用什么软件？ B: 我用微信，你呢？ A: 我用WhatsApp。微信能付钱吗？ B: 当然，买东西、坐地铁都用微信支付。}
\end{experience}

\begin{exercice}[À vous de répondre — Écrit]
Répondez par des phrases complètes en chinois (caractères + pinyin) :
\begin{enumerate}
\item \zh{年轻人用什么上网？}
\item \zh{为什么网络让我们的生活更方便？}
\item \zh{花太多时间上网有什么问题？}
\item Résumez le texte en deux phrases avec vos propres mots (utilisez \zh{但是} et \zh{所以}).
\end{enumerate}
\end{exercice}

\begin{corrige}[Exercice « À vous de répondre »]
\begin{enumerate}
\item \zh{年轻人一般用手机上网。} \emph{Niánqīngrén yìbān yòng shǒujī shàngwǎng.} « Les jeunes se connectent généralement avec leur téléphone portable. »
\item \zh{因为我们能在网上看视频、买东西，所以网络让我们的生活更方便。} \emph{Yīnwèi wǒmen néng zài wǎngshàng kàn shìpín, mǎi dōngxi, suǒyǐ wǎngluò ràng wǒmen de shēnghuó gèng fāngbiàn.} « Parce que nous pouvons regarder des vidéos et acheter des choses en ligne, internet rend donc notre vie plus pratique. »
\item \zh{花太多时间上网对身体不好，也会影响学习。} \emph{Huā tài duō shíjiān shàngwǎng duì shēntǐ bù hǎo, yě huì yǐngxiǎng xuéxí.} « Passer trop de temps en ligne est mauvais pour la santé et nuit aux études. »
\item Exemple : \zh{网络很有用，但是上网时间太长不好，所以我们要合理安排时间。} \emph{Wǎngluò hěn yǒuyòng, dànshì shàngwǎng shíjiān tài cháng bù hǎo, suǒyǐ wǒmen yào hélǐ ānpái shíjiān.} « Internet est utile, mais trop de temps en ligne est mauvais, nous devons donc organiser raisonnablement notre temps. »
\end{enumerate}
\end{corrige}

\begin{aretenir}[L'essentiel de la leçon 34]
\begin{itemize}
\item \textbf{Vocabulaire} : maîtriser les 20 mots du glossaire (caractères, pinyin, sens, nature). Savoir écrire \zh{网络}, \zh{微信}, \zh{视频}, \zh{聊天}, \zh{身体}, \zh{影响}, \zh{安排}.
\item \textbf{Radicaux} : identifier \zh{纟} (réseau), \zh{讠/言} (parole/info), \zh{彳} (action), \zh{示} (voir), \zh{页} (tête/page), \zh{艹} (plante), \zh{亻} (humain), \zh{骨/身} (corps).
\item \textbf{Compréhension} : méthode en 3 étapes (mot interrogatif $\rightarrow$ localisation $\rightarrow$ phrase complète).
\item \textbf{Grammaire clé} : structure \zh{因为…所以…} pour répondre à \zh{为什么} ; \zh{如果…} (condition) ; \zh{应该} (conseil) ; \zh{对…不好} (mauvais pour…).
\item \textbf{Expression} : savoir résumer un texte argumentatif (plan : situation $\rightarrow$ usages $\rightarrow$ avantages $\rightarrow$ inconvénients $\rightarrow$ conseil) avec connecteurs \zh{但是}, \zh{所以}.
\item \textbf{Culture} : WeChat (\zh{微信}) est un « super-app » (messagerie + réseau social + paiement + services publics) sans équivalent unique au Cameroun (WhatsApp + Mobile Money + Facebook séparés). La régulation d'internet (\zh{防火长城}) diffère radicalement.
\end{itemize}
\end{aretenir}

\coursec{Production guidée et comparaison Cameroun-Chine}

Après avoir vérifié votre compréhension du texte support grâce aux questions de la section précédente, il est temps de passer à l'étape la plus importante de l'approche par les compétences : \textbf{produire vous-même}. Vous allez maintenant réutiliser le vocabulaire et les structures de la leçon pour écrire un court paragraphe sur vos propres habitudes sur internet, comparer les réseaux sociaux du Cameroun et de la Chine, puis débattre en binômes. C'est exactement le type de tâche demandé à l'examen : rédiger quelques phrases correctes, connectées et personnelles sur un thème du programme.

\courssub{Observer d'abord : un modèle de production}

Avant d'écrire, observons comment un élève de Terminale peut parler de sa vie numérique en quelques phrases simples. Lisez le modèle suivant, repérez les connecteurs et la trame utilisée.

\begin{textebox}[Modèle de production : ma vie sur internet]
我每天用手机上网。\\
\emph{Wǒ měitiān yòng shǒujī shàngwǎng.}\\
« Je me connecte à internet tous les jours avec mon téléphone. »\\[4pt]
我用微信和朋友聊天，也在朋友圈发照片。\\
\emph{Wǒ yòng Wēixìn hé péngyou liáotiān, yě zài péngyouquān fā zhàopiàn.}\\
« J'utilise WeChat pour discuter avec mes amis, et je publie aussi des photos dans les moments. »\\[4pt]
网络让我的生活更方便。\\
\emph{Wǎngluò ràng wǒ de shēnghuó gèng fāngbiàn.}\\
« Internet rend ma vie plus pratique. »\\[4pt]
但是花太多时间上网影响学习，所以我应该合理安排时间。\\
\emph{Dànshì huā tài duō shíjiān shàngwǎng yǐngxiǎng xuéxí, suǒyǐ wǒ yīnggāi hélǐ ānpái shíjiān.}\\
« Mais passer trop de temps en ligne nuit aux études, donc je dois organiser mon temps de façon raisonnable. »
\end{textebox}

Observez la logique du texte : \textbf{fréquence} (\zh{每天}) → \textbf{outil} (\zh{用手机}) → \textbf{activités} (\zh{聊天、发照片}) → \textbf{avantage} (\zh{更方便}) → \textbf{limite} (\zh{但是}…) → \textbf{conclusion personnelle} (\zh{所以我应该}…). C'est exactement la trame que vous allez suivre. Notez aussi deux points de grammaire déjà vus et réemployés ici : le complément de moyen avec \zh{用} + outil placé \textbf{avant le verbe} (\zh{用手机上网}), et la structure \zh{让} + personne + adjectif (\zh{让我的生活更方便} « rendre ma vie plus pratique »).

\courssub{La trame de production guidée}

Voici la trame à trous à compléter avec vos propres informations. Chaque blanc correspond à un élément du vocabulaire de la leçon.

\begin{methode}[Rédiger 5 à 6 phrases sur internet : la trame]
\begin{enumerate}
\item \zh{我} \_\_\_ (fréquence : \zh{每天、常常、有时候、很少}) \zh{上网。} — J'indique \textbf{à quelle fréquence} je me connecte. L'adverbe de fréquence se place entre le sujet et le verbe.
\item \zh{我用} \_\_\_ (appareil : \zh{手机、电脑、平板}) \zh{。} — Je précise \textbf{l'outil} utilisé.
\item \zh{我在网上} \_\_\_ (activités : \zh{聊天、看视频、学习、网购、发照片}) \zh{。} — Je décris \textbf{mes activités} en ligne ; \zh{在网上} « sur internet » se place avant le verbe.
\item \zh{网络让我的生活} \_\_\_ (\zh{更方便、更丰富}) \zh{。} — Je donne \textbf{un effet positif}.
\item \zh{但是} \_\_\_ (un inconvénient : \zh{花太多时间、影响学习、对身体不好}) \zh{。} — Je nuance avec \textbf{un inconvénient}.
\item \zh{所以我应该} \_\_\_ (une résolution : \zh{合理安排时间、少玩游戏、多学习}) \zh{。} — Je conclus par \textbf{une résolution}.
\end{enumerate}
\end{methode}

\begin{exemplebox}[Banque de mots et connecteurs]
\begin{itemize}
\item \zh{因为网络很方便，所以大家都喜欢上网。} \emph{Yīnwèi wǎngluò hěn fāngbiàn, suǒyǐ dàjiā dōu xǐhuan shàngwǎng.} « Comme internet est pratique, tout le monde aime se connecter. »
\item \zh{上网有好处，也有坏处。} \emph{Shàngwǎng yǒu hǎochù, yě yǒu huàichù.} « Internet a des avantages et des inconvénients. »
\item \zh{我不仅用微信聊天，还用它付钱。} \emph{Wǒ bùjǐn yòng Wēixìn liáotiān, hái yòng tā fùqián.} « Je n'utilise pas seulement WeChat pour discuter, je m'en sers aussi pour payer. »
\item \zh{越来越多的喀麦隆年轻人用社交媒体。} \emph{Yuè lái yuè duō de Kāmàilóng niánqīngrén yòng shèjiāo méitǐ.} « De plus en plus de jeunes Camerounais utilisent les réseaux sociaux. »
\item \zh{我也喜欢在网上看足球比赛。} \emph{Wǒ yě xǐhuan zài wǎngshang kàn zúqiú bǐsài.} « Moi aussi, j'aime regarder les matchs de football en ligne. »
\end{itemize}
\end{exemplebox}

\begin{propriete}[Les connecteurs de la trame : forme et emploi]
\begin{itemize}
\item \textbf{因为…所以…} « parce que… donc… » : en chinois, on peut utiliser les deux ensemble dans la même phrase (contrairement au français où « parce que… donc » est incorrect). Structure : \cle{因为} + cause，\cle{所以} + conséquence。On peut aussi n'employer que l'un des deux.
\item \textbf{但是} « mais » : se place \textbf{en tête de la seconde proposition}, devant le sujet s'il y en a un.
\item \textbf{也} « aussi » et \textbf{还} « en plus, encore » : se placent \textbf{après le sujet, avant le verbe}. \zh{我也上网。} \emph{Wǒ yě shàngwǎng.} « Moi aussi je me connecte. »
\item \textbf{越来越} + adjectif : « de plus en plus… ». \zh{越来越方便} \emph{yuè lái yuè fāngbiàn} « de plus en plus pratique ». Attention : l'adjectif après \zh{越来越} ne prend jamais \zh{很}.
\end{itemize}
\end{propriete}

\begin{attention}[La place de 但是 : l'erreur classique du francophone]
En français, « mais » peut sembler mobile ; en chinois, \zh{但是} est \textbf{toujours en début de seconde proposition}, jamais entre le sujet et le verbe.\\[4pt]
\zh{但是花太多时间上网对身体不好。} ✓ \emph{Dànshì huā tài duō shíjiān shàngwǎng duì shēntǐ bù hǎo.} « Mais passer trop de temps en ligne est mauvais pour la santé. »\\[4pt]
Ne dites jamais : \zh{花太多时间但是上网不好} ✗ — le connecteur doit ouvrir la proposition. Même vigilance pour \zh{所以} : il se place en tête de la proposition de conséquence, pas à l'intérieur. Autre piège : ne dites pas \zh{很方便越来越} ✗ ; la bonne forme est \zh{越来越方便} ✓.
\end{attention}

\courssub{Comparaison guidée : les réseaux sociaux au Cameroun et en Chine}

Un point essentiel de culture numérique : les applications ne sont pas les mêmes partout. Au Cameroun, les jeunes utilisent surtout WhatsApp, Facebook et TikTok ; en Chine, ce sont \zh{微信} (Wēixìn, WeChat), \zh{微博} (Wēibó, Weibo) et \zh{抖音} (Dǒuyīn, Douyin). Observons ce tableau comparatif.

\begin{popfigure}
\begin{center}
\begin{tikzpicture}
\node[fill=popgreenL, rounded corners=3pt, minimum width=6.2cm, minimum height=0.9cm, align=center, font=\bfseries] (cam) at (0,0) {Cameroun};
\node[fill=popblueL, rounded corners=3pt, minimum width=6.2cm, minimum height=0.9cm, align=center, font=\bfseries] (chn) at (7,0) {Chine};
\node[draw=popline, rounded corners=3pt, minimum width=6.2cm, minimum height=0.8cm, align=center] at (0,-1) {WhatsApp — messagerie};
\node[draw=popline, rounded corners=3pt, minimum width=6.2cm, minimum height=0.8cm, align=center] at (7,-1) {微信 Wēixìn — messagerie};
\node[draw=popline, rounded corners=3pt, minimum width=6.2cm, minimum height=0.8cm, align=center] at (0,-2) {Facebook — réseau social};
\node[draw=popline, rounded corners=3pt, minimum width=6.2cm, minimum height=0.8cm, align=center] at (7,-2) {微博 Wēibó — microblog};
\node[draw=popline, rounded corners=3pt, minimum width=6.2cm, minimum height=0.8cm, align=center] at (0,-3) {TikTok — vidéos courtes};
\node[draw=popline, rounded corners=3pt, minimum width=6.2cm, minimum height=0.8cm, align=center] at (7,-3) {抖音 Dǒuyīn — vidéos courtes};
\draw[<->, thick, poporange] (3.2,-1) -- (3.8,-1);
\draw[<->, thick, poporange] (3.2,-2) -- (3.8,-2);
\draw[<->, thick, poporange] (3.2,-3) -- (3.8,-3);
\end{tikzpicture}
\end{center}
\legende{Équivalences fonctionnelles entre les applications utilisées au Cameroun et en Chine}
\end{popfigure}

Pour comparer à l'oral ou à l'écrit, utilisez les structures de comparaison déjà vues :

\begin{center}
\begin{tabularx}{\linewidth}{|X|X|X|}
\hline
\rowcolor{popblueL} Structure & Exemple en chinois + pinyin & Traduction \\
\hline
A 和 B 不一样 & \zh{微信和WhatsApp不一样。} \emph{Wēixìn hé WhatsApp bù yíyàng.} & WeChat et WhatsApp ne sont pas pareils. \\
\hline
A 比 B + adj. & \zh{在中国，微信比微博更常用。} \emph{Zài Zhōngguó, Wēixìn bǐ Wēibó gèng chángyòng.} & En Chine, WeChat est plus courant que Weibo. \\
\hline
越来越 + adj. & \zh{网购越来越流行。} \emph{Wǎnggòu yuè lái yuè liúxíng.} & Les achats en ligne sont de plus en plus populaires. \\
\hline
\end{tabularx}
\end{center}

\begin{savaistu}[Douyin et TikTok : deux sœurs qui ne se rencontrent jamais]
\zh{抖音} (Dǒuyīn) et TikTok appartiennent à la même entreprise chinoise, ByteDance, mais ce sont \textbf{deux applications totalement séparées} : TikTok n'existe pas en Chine continentale, et Douyin n'est pas disponible dans les boutiques d'applications étrangères. Les contenus, les serveurs et même les comptes sont différents. Autre fait utile : \zh{微博} (Wēibó), littéralement « micro-blog », fonctionne comme un équivalent chinois de X (Twitter) : on y suit l'actualité et les célébrités.
\end{savaistu}

\textbf{Complément utile à l'examen} (à retenir) :

\begin{center}
\begin{tabularx}{\linewidth}{|X|X|X|X|}
\hline
\rowcolor{popblueL} Caractères & Pinyin & Nature & Traduction \\
\hline
微博 & Wēibó & n. propre & Weibo (microblog chinois) \\
\hline
抖音 & Dǒuyīn & n. propre & Douyin (TikTok chinois) \\
\hline
网购 & wǎnggòu & n. / v. & achats en ligne ; acheter en ligne \\
\hline
朋友圈 & péngyouquān & n. & « moments » (fil d'actualité WeChat) \\
\hline
合理安排 & hélǐ ānpái & loc. v. & organiser de façon raisonnable \\
\hline
\end{tabularx}
\end{center}

\courssub{Débat en binômes : 上网好还是不好？}

\begin{experience}[Débat guidé : Internet, bien ou mal ?]
Par binômes, l'un défend la thèse « internet est bon » (\zh{上网好}), l'autre la thèse « internet est mauvais » (\zh{上网不好}). Chacun donne au moins deux arguments tirés du texte support, avec la structure \zh{因为…所以…}.\\[4pt]
Arguments possibles « pour » : \zh{上网可以学习、和朋友联系、了解新闻、网购很方便。} « On peut étudier, rester en contact avec ses amis, suivre l'actualité, et les achats en ligne sont pratiques. »\\[4pt]
Arguments possibles « contre » : \zh{花太多时间上网影响学习、对眼睛不好、有时候有假消息。} « Passer trop de temps en ligne nuit aux études, fatigue les yeux, et il y a parfois de fausses informations. »\\[4pt]
Conclusion commune obligatoire : \zh{上网有好处，也有坏处，所以我们应该合理安排时间。} « Internet a des avantages et des inconvénients, donc nous devons organiser notre temps raisonnablement. »
\end{experience}

\courssub{Exercice résolu et correction collective d'une production modèle}

\begin{exoresolu}[Complétez la trame avec les mots de la banque]
Énoncé : complétez le paragraphe suivant avec les mots de la banque (\zh{每天、手机、聊天、方便、但是、所以}) :\\
\zh{我\_\_\_上网。我用\_\_\_和朋友\_\_\_。网络让我的生活更\_\_\_。\_\_\_花太多时间上网影响学习，\_\_\_我应该合理安排时间。}
\tcblower
Solution détaillée :\\
\zh{我}\textbf{每天}\zh{上网。} — la fréquence se place avant le verbe \zh{上网}.\\
\zh{我用}\textbf{手机}\zh{和朋友}\textbf{聊天}\zh{。} — \zh{用} + outil ; \zh{和朋友聊天} « discuter avec des amis ».\\
\zh{网络让我的生活更}\textbf{方便}\zh{。} — \zh{让} + personne + adjectif : « rendre quelqu'un/quelque chose + adj. ».\\
\textbf{但是}\zh{花太多时间上网影响学习，}\textbf{所以}\zh{我应该合理安排时间。} — \zh{但是} ouvre la proposition d'opposition, \zh{所以} ouvre la conclusion.\\
Traduction complète : « Je me connecte tous les jours. J'utilise mon téléphone pour discuter avec mes amis. Internet rend ma vie plus pratique. Mais passer trop de temps en ligne nuit aux études, donc je dois organiser mon temps raisonnablement. »
\end{exoresolu}

\begin{methode}[Correction collective d'une production au tableau]
\begin{enumerate}
\item Un élève écrit sa production de 5 à 6 phrases au tableau, sans la lire.
\item La classe vérifie d'abord les \textbf{connecteurs} : \zh{但是} et \zh{所以} sont-ils en tête de proposition ? \zh{也} et \zh{还} sont-ils avant le verbe ?
\item On vérifie ensuite l'\textbf{ordre des mots} : sujet + fréquence + verbe ; \zh{用} + outil avant l'action ; \zh{在网上} avant le verbe.
\item On corrige enfin les \textbf{caractères} : attention à \zh{网} (clé \zh{罒}, le filet), à \zh{聊} (clé \zh{耳}, l'oreille) et à \zh{影} de \zh{影响} (ne pas oublier la partie \zh{景} à gauche).
\item Le professeur réécrit la version corrigée en couleur ; chaque élève compare avec sa propre copie et note ses erreurs.
\end{enumerate}
\end{methode}

\begin{exercice}[À vous de produire]
Rédigez votre propre paragraphe de 5 à 6 phrases en suivant la trame : fréquence, outil, activités, avantage, inconvénient, résolution. Utilisez au moins une fois \zh{因为…所以…} et une fois \zh{越来越}. Terminez par une phrase de comparaison Cameroun-Chine, par exemple : \zh{在喀麦隆我们用WhatsApp，在中国人们用微信。} « Au Cameroun nous utilisons WhatsApp, en Chine les gens utilisent WeChat. »
\end{exercice}

\begin{corrige}[Exercice : production libre]
Exemple de production attendue (une correction possible parmi d'autres) :\\
\zh{我常常用手机上网。我在网上聊天、看视频，也学习汉语。网络让我的生活越来越方便。但是花太多时间上网对眼睛不好，因为上网有好处也有坏处，所以我应该合理安排时间。在喀麦隆我们用WhatsApp，在中国人们用微信。}\\
\emph{Wǒ chángcháng yòng shǒujī shàngwǎng. Wǒ zài wǎngshang liáotiān, kàn shìpín, yě xuéxí Hànyǔ. Wǎngluò ràng wǒ de shēnghuó yuè lái yuè fāngbiàn. Dànshì huā tài duō shíjiān shàngwǎng duì yǎnjing bù hǎo, yīnwèi shàngwǎng yǒu hǎochù yě yǒu huàichù, suǒyǐ wǒ yīnggāi hélǐ ānpái shíjiān. Zài Kāmàilóng wǒmen yòng WhatsApp, zài Zhōngguó rénmen yòng Wēixìn.}\\
Critères de réussite : 5 à 6 phrases complètes, connecteurs bien placés, au moins un mot du complément d'examen (\zh{网购}, \zh{微博}, \zh{抖音}) bienvenu.
\end{corrige}

\begin{aretenir}[L'essentiel de la section]
\begin{itemize}
\item Une bonne production suit un plan : fréquence → outil → activités → avantage → inconvénient → résolution.
\item \zh{因为…所以…} peut s'employer en entier dans une même phrase chinoise ; \zh{但是} et \zh{所以} ouvrent toujours leur proposition ; \zh{也} et \zh{还} se placent avant le verbe ; après \zh{越来越}, l'adjectif ne prend pas \zh{很}.
\item Cameroun : WhatsApp, Facebook, TikTok ; Chine : \zh{微信}, \zh{微博}, \zh{抖音}. À l'examen, connaissez aussi \zh{网购} « achats en ligne » et \zh{朋友圈} « moments ».
\item Pour débattre : deux arguments avec \zh{因为…所以…}, puis une conclusion équilibrée : \zh{上网有好处，也有坏处。}
\end{itemize}
\end{aretenir}

\newpage

\coursec{Activité d'intégration}

\coursec{Leçon 34 — Vocabulaire et compréhension de texte : internet et les réseaux sociaux}

\courssub{Objectifs de la leçon}
À l'issue de cette leçon, l'élève de Terminale A sera capable de :
\begin{itemize}
\item lire et comprendre un texte authentique (niveau HSK 3) traitant de l'usage d'Internet et des réseaux sociaux au Cameroun et en Chine ;
\item reconnaître, prononcer correctement (pinyin, tons) et écrire les caractères clés du thème ;
\item mobiliser un lexique d'au moins 15 mots (verbes, noms, adjectifs, connecteurs) pour décrire ses habitudes numériques ;
\item employer les structures argumentatives \zh{因为……所以……} (cause-conséquence), \zh{但是} (opposition) et \zh{以后，我要……} (résolution) dans un ordre des mots correct ;
\item rédiger un court texte personnel structuré (6--8 phrases) : \zh{我和网络} ;
\item adopter une attitude critique et responsable face aux écrans, en comparant les contextes camerounais et chinois.
\end{itemize}

\courssub{Mise en route et découverte du thème}
\begin{experience}[Brainstorming : « Mon écran et moi »]
L'enseignant projette ou affiche au tableau deux images : un jeune Camerounais à Douala tenant un smartphone (logo WhatsApp, Facebook visible) et un jeune Chinois à Pékin avec WeChat, Douyin. En binômes, les élèves listent en 2 minutes 5 mots chinois qu'ils associent à ces images (ex. : \zh{手机}, \zh{上网}, \zh{聊天}, \zh{视频}, \zh{新闻}). Mise en commun au tableau : l'enseignant écrit les caractères, corrige les tons, introduit le radical \zh{网} (filet/réseau) et annonce le thème du module 5 : \zh{大众传媒与电信} (Médias de masse et télécommunications).
\end{experience}

\courssub{Texte support : Lecture et compréhension globale}
\begin{textebox}[Texte original : Les jeunes et Internet — Yaoundé / Pékin]
\zh{保罗是喀麦隆雅温得的高中生。他每天放学后用手机上网，通常上两三个小时。他最常用WhatsApp和朋友聊天，也看足球视频和新闻。李明是中国北京的高中生。他也每天用手机上网，但他常用微信，还在网上学习英语和法语。保罗觉得网络很方便，因为可以很快找到信息。李明也觉得网络有好处，但是提醒：上网太久对眼睛不好，容易浪费时间。两人都说：以后，我要合理安排时间，少玩手机，多读书、多运动。}\\
Pǔluó shì Kāmàilóng Yǎwēndé de gāozhōngshēng. Tā měi tiān fàngxué hòu yòng shǒujī shàngwǎng, tōngcháng shàng liǎng sān ge xiǎoshí. Tā zuì cháng yòng WhatsApp hé péngyou liáotiān, yě kàn zúqiú shìpín hé xīnwén. Lǐ Míng shì Zhōngguó Běijīng de gāozhōngshēng. Tā yě měi tiān yòng shǒujī shàngwǎng, dàn tā cháng yòng Wēixìn, hái zài wǎngshang xuéxí Yīngyǔ hé Fǎyǔ. Pǔluó juéde wǎngluò hěn fāngbiàn, yīnwèi kěyǐ hěn kuài zhǎodào xìnxī. Lǐ Míng yě juéde wǎngluò yǒu hǎochù, dànshì tíxǐng: shàngwǎng tài jiǔ duì yǎnjing bù hǎo, róngyì làngfèi shíjiān. Liǎng rén dōu shuō: yǐhòu, wǒ yào hélǐ ānpái shíjiān, shǎo wán shǒujī, duō dúshū, duō yùndòng.\\
« Paul est lycéen à Yaoundé, Cameroun. Après les cours, il se connecte chaque jour avec son portable, généralement 2 à 3 heures. Il utilise surtout WhatsApp pour discuter avec ses amis, regarde aussi des vidéos de foot et les infos. Li Ming est lycéen à Pékin, Chine. Il se connecte aussi tous les jours sur portable, mais utilise surtout WeChat, et étudie l'anglais et le français en ligne. Paul trouve Internet très pratique car on peut y trouver vite de l'information. Li Ming trouve aussi Internet avantageux, mais prévient : rester trop longtemps en ligne est mauvais pour les yeux et fait perdre du temps. Tous deux disent : à l'avenir, je dois organiser mon temps raisonnablement, moins jouer au téléphone, plus lire et plus faire du sport. »
\end{textebox}

\begin{exercice}[Compréhension du texte support]
Répondez en chinois (phrases complètes) aux questions suivantes :
\begin{enumerate}
\item \zh{保罗和李明都是谁？他们住在哪儿？}
\item \zh{保罗常用什么软件聊天？李明呢？}
\item \zh{保罗为什么觉得网络方便？}
\item \zh{李明提醒大家上网太久有什么坏处？}
\item \zh{他们两人共同的决定是什么？（用“以后，我要……”）}
\end{enumerate}
\end{exercice}

\begin{corrige}[Exercice Compréhension]
\begin{enumerate}
\item \zh{他们都是高中生。保罗住在喀麦隆雅温得，李明住在中国北京。} (Tāmen dōu shì gāozhōngshēng. Pǔluó zhù zài Kāmàilóng Yǎwēndé, Lǐ Míng zhù zài Zhōngguó Běijīng.)
\item \zh{保罗常用WhatsApp，李明常用微信。} (Pǔluó cháng yòng WhatsApp, Lǐ Míng cháng yòng Wēixìn.)
\item \zh{因为可以在网上很快找到信息。} (Yīnwèi kěyǐ zài wǎngshang hěn kuài zhǎodào xìnxī.)
\item \zh{对眼睛不好，也容易浪费时间。} (Duì yǎnjing bù hǎo, yě róngyì làngfèi shíjiān.)
\item \zh{以后，我要合理安排时间，少玩手机，多读书、多运动。} (Yǐhòu, wǒ yào hélǐ ānpái shíjiān, shǎo wán shǒujī, duō dúshū, duō yùndòng.)
\end{enumerate}
\end{corrige}

\courssub{Glossaire du thème : Vocabulaire essentiel}
\begin{aretenir}[Lexique du thème « Internet et réseaux sociaux » (HSK 2--4)]
\begin{center}
\begin{tabularx}{\linewidth}{|X|X|X|X|}
\hline
\rowcolor{popblueL} Caractères & Pinyin & Nature & Traduction \\
\hline
上网 & shàngwǎng & verbe & se connecter (à Internet), aller en ligne \\
\hline
网络 & wǎngluò & nom & Internet, le réseau, la toile \\
\hline
手机 & shǒujī & nom & téléphone portable, smartphone \\
\hline
电脑 & diànnǎo & nom & ordinateur \\
\hline
微信 & Wēixìn & nom propre & WeChat (appli chinoise omniprésente) \\
\hline
WhatsApp & WhatsApp & nom propre & WhatsApp (très utilisé au Cameroun) \\
\hline
聊天 & liáotiān & verbe & bavarder, discuter (en ligne ou oral) \\
\hline
新闻 & xīnwén & nom & informations, actualités, nouvelles \\
\hline
视频 & shìpín & nom & vidéo, clip vidéo \\
\hline
信息 & xìnxī & nom & information, message, donnée \\
\hline
方便 & fāngbiàn & adjectif & pratique, commode \\
\hline
好处 & hǎochù & nom & avantage, côté positif, bienfait \\
\hline
危险 & wēixiǎn & adjectif / nom & dangereux ; danger, risque \\
\hline
浪费 & làngfèi & verbe & gaspiller (temps, argent, eau\ldots) \\
\hline
眼睛 & yǎnjing & nom & œil(s) \\
\hline
合理 & hélǐ & adjectif & raisonnable, sensé \\
\hline
安排 & ānpái & verbe & arranger, organiser, planifier \\
\hline
运动 & yùndòng & verbe / nom & faire du sport ; sport \\
\hline
\hline
\multicolumn{4}{|l|}{\textbf{Connecteurs et mots-outils du thème}} \\
\hline
因为……所以…… & yīnwèi\ldots suǒyǐ\ldots & conj. & parce que\ldots donc\ldots (cause-conséquence) \\
\hline
但是 & dànshì & conj. & mais, cependant (opposition) \\
\hline
以后 & yǐhòu & nom temps & à l'avenir, plus tard, après \\
\hline
通常 & tōngcháng & adv. & généralement, d'habitude \\
\hline
最常 & zuì cháng & adv. & le plus souvent \\
\hline
\end{tabularx}
\end{center}
\end{aretenir}

\courssub{Radicaux et ordre des traits : Clés d'écriture}
\begin{propriete}[Analyse de 6 caractères clés du thème]
\begin{center}
\begin{tabularx}{\linewidth}{|X|X|X|X|}
\hline
\rowcolor{popblueL} Caractère & Radical (clé) & Nombre de traits & Ordre des traits (règles principales) \\
\hline
\zh{网} (wǎng) & \zh{纟} (sī, soie/filer) — 3 traits & 6 & 1. \zh{纟} (haut→bas : point, pli, point) ; 2. \zh{亡} (wáng : horizontales, verticale, point final) \\
\hline
\zh{手} (shǒu) & \zh{手} (shǒu, main) — 4 traits & 4 & Verticale centrale → horizontale courte → crochet vertical → horizontale longue (base) \\
\hline
\zh{机} (jī) & \zh{木} (mù, bois/arbre) — 4 traits & 16 & 1. \zh{木} (horiz, vert, g. gauche, g. droite) ; 2. \zh{幂} (jī, phonétique : enclos + traits intérieurs) \\
\hline
\zh{信} (xìn) & \zh{亻} (rén, homme) — 2 traits & 9 & 1. \zh{亻} (traits gauche→droite) ; 2. \zh{言} (yán, parole : point, horiz, horiz, vert, 3 horizontales, point final) \\
\hline
\zh{视} (shì) & \zh{见} (jiàn, voir) — 7 traits & 11 & 1. \zh{礻} (shì, autel/rite : point, horiz, g. gauche, g. droite, point) ; 2. \zh{见} (horiz, vert, horiz, 3 horizontales, point) \\
\hline
\zh{危} (wēi) & \zh{危} (wēi, danger) — 6 traits (clé elle-même) & 6 & Haut \zh{厂} (hán, falaise : horiz, verticale, horiz) → Bas \zh{㐆} (qí : verticale, 2 horizontales, crochet) \\
\hline
\end{tabularx}
\end{center}
\emph{Rappel des règles universelles :} haut → bas ; gauche → droite ; horizontale avant verticale ; trait central avant ailes ; enclos avant contenu ; trait de fermeture final.
\end{propriete}

\begin{popfigure}
\begin{tikzpicture}[
  mindmap,
  grow cyclic,
  every node/.style={concept, execute at begin node=\hskip0pt},
  root concept/.append style={concept color=popblue, minimum size=2.5cm, text=white, font=\bfseries\large},
  level 1 concept/.append style={concept color=popblue!70, sibling angle=90, minimum size=1.8cm, text=white, font=\bfseries},
  level 2 concept/.append style={concept color=poporange!80, sibling angle=45, minimum size=1.4cm, text=white, font=\small},
  level 3 concept/.append style={concept color=popgreen!70, sibling angle=30, minimum size=1.2cm, text=white, font=\tiny}
]
\node {Internet\ \& Réseaux}
  child [clockwise from=60] {
    node {Appareils}
    child { node {\zh{手机} shǒujī} }
    child { node {\zh{电脑} diànnǎo} }
  }
  child [clockwise from=150] {
    node {Actions}
    child { node {\zh{上网} shàngwǎng} }
    child { node {\zh{聊天} liáotiān} }
    child { node {\zh{看视频} kàn shìpín} }
    child { node {\zh{学习} xuéxí} }
  }
  child [clockwise from=240] {
    node {Applis}
    child { node {\zh{微信} Wēixìn} }
    child { node {WhatsApp} }
  }
  child [clockwise from=330] {
    node {Évaluation}
    child { node {\zh{方便} fāngbiàn} }
    child { node {\zh{好处} hǎochù} }
    child { node {\zh{危险} wēixiǎn} }
    child { node {\zh{浪费时间} làngfèi shíjiān} }
  };
\end{tikzpicture}
\legende{Carte mentale du vocabulaire : organisez vos révisions par catégories.}
\end{popfigure}

\courssub{Points de grammaire : Connecteurs argumentatifs}
\begin{propriete}[Structure \cle{因为……所以……} (Cause $\rightarrow$ Conséquence)]
\textbf{Schéma :} \textbf{Sujet + 因为 + Cause , 所以 + Conséquence.}
\begin{itemize}
\item Les deux marqueurs \zh{因为} (yīnwèi) et \zh{所以} (suǒyǐ) sont généralement tous deux présents en chinois écrit/formel, contrairement au français où « parce que » suffit souvent seul.
\item L'ordre est immuable : Cause d'abord, Conséquence ensuite.
\item \zh{所以} peut être précédé du sujet s'il change, ou omis si le sujet est identique (mais on le garde pour la clarté en Terminale).
\end{itemize}
\end{propriete}

\begin{exemplebox}[Exemples \zh{因为……所以……}]
\zh{因为网络很方便，所以我每天都上网。} (Yīnwèi wǎngluò hěn fāngbiàn, suǒyǐ wǒ měi tiān dōu shàngwǎng.) — « Parce qu'Internet est pratique, je me connecte tous les jours. » \\
\zh{我喜欢中文，因为我觉得汉字很有意思，所以我想去中国留学。} (Wǒ xǐhuan Zhōngwén, yīnwèi wǒ juéde Hànzì hěn yǒu yìsi, suǒyǐ wǒ xiǎng qù Zhōngguó liúxué.) — « J'aime le chinois, car je trouve les caractères intéressants, donc je veux aller étudier en Chine. »
\end{exemplebox}


\begin{propriete}[Structure \cle{但是} (Opposition / Concession)]
\textbf{Schéma :} \textbf{Proposition 1 , 但是 + Proposition 2.}
\begin{itemize}
\item \zh{但是} (dànshì) splace en début de seconde proposition (après la virgule). Il marque un contraste fort.
\item Synonyme plus formel : \zh{可是} (kěshì) ; plus littéraire : \zh{然而} (rán'ér).
\item Ne pas confondre avec \zh{不过} (bùguò) « toutefois, simplement » (plus léger, souvent oral).
\end{itemize}
\end{propriete}

\begin{exemplebox}[Exemples \zh{但是}]
\zh{网络很有用，但是也有危险。} (Wǎngluò hěn yǒuyòng, dànshì yě yǒu wēixiǎn.) — « Internet est très utile, mais présente aussi des dangers. » \\
\zh{我想玩游戏，但是我要复习功课。} (Wǒ xiǎng wán yóuxì, dànshì wǒ yào fùxí gōngkè.) — « Je veux jouer, mais je dois réviser. »
\end{exemplebox}


\begin{propriete}[Structure \cle{以后，我要……} (Résolution future)]
\textbf{Schéma :} \textbf{以后 , Sujet + 要 / 打算 / 计划 + Verbe (+ Complément).}
\begin{itemize}
\item \zh{以后} (yǐhòu) = « à l'avenir, dorénavant » (marqueur temporel en position initiale de phrase).
\item \zh{要} (yào) exprime l'intention forte, l'obligation morale ou le projet décidé. \zh{打算} (dǎsuàn) et \zh{计划} (jìhuà) sont plus souples (« avoir l'intention de », « prévoir de »).
\item Structure antithétique fréquente : \zh{少……多……} (moins\ldots plus\ldots).
\end{itemize}
\end{propriete}

\begin{exemplebox}[Exemples Résolution]
\zh{以后，我要少上网，多读书。} (Yǐhòu, wǒ yào shǎo shàngwǎng, duō dúshū.) — « À l'avenir, je me connecterai moins, lirai plus. » \\
\zh{以后，我打算每天只上网一小时。} (Yǐhòu, wǒ dǎsuàn měi tiān zhǐ shàngwǎng yī ge xiǎoshí.) — « Dorénavant, je prévois de me connecter seulement 1h par jour. »
\end{exemplebox}


\courssub{Attention : Erreurs fréquentes des francophones}
\begin{attention}[Pièges à éviter absolument]
\begin{itemize}
\item \textbf{\zh{上网} vs \zh{上网络}} : \zh{上网} (shàngwǎng) est un verbe lexicalisé indissociable (verbe + objet incorporé). On ne dit \textbf{JAMAIS} \zh{上网络}. \checkmark \zh{我上网} / \zh{我用手机上网}. \cross \zh{我上网络} / \zh{我用手机上网络}.
\item \textbf{Place du complément de moyen \zh{用\ldots}} : En chinois, l'instrument précède le verbe. \checkmark \zh{我用手机上网} (Sujet + 用 + Instrument + Verbe). \cross \zh{我上网用手机} (ordre français SVO incorrect ici).
\item \textbf{\zh{很} + Adjectif prédicat} : Un adjectif en fonction de prédicat (sans verbe \zh{是}) requiert presque toujours un modificateur de degré (\zh{很}, \zh{真}, \zh{特别}, \zh{比较}). \checkmark \zh{网络很方便}. \cross \zh{网络方便} (phrase incomplète, sonne « télégraphique »).
\item \textbf{Tons critiques} :
  \begin{itemize}
  \item \zh{上网} : \textbf{4\textsuperscript{e} + 3\textsuperscript{e}} (shàngwǎng) — attention au 3\textsuperscript{e} ton bas sur \zh{网}.
  \item \zh{手机} : \textbf{3\textsuperscript{e} + 1\textsuperscript{er}} (shǒujī) — 3\textsuperscript{e} ton bas sur \zh{手}, 1\textsuperscript{er} ton haut plat sur \zh{机}.
  \item \zh{微信} : \textbf{1\textsuperscript{er} + 4\textsuperscript{e}} (Wēixìn) — ne pas dire Wéixìn (2+4) ni Wèixìn (4+4).
  \item \zh{危险} : \textbf{1\textsuperscript{er} + 3\textsuperscript{e}} (wēixiǎn).
  \end{itemize}
\item \textbf{Classificateur \zh{个} (ge) pour les heures} : \zh{两个小时} (liǎng ge xiǎoshí) — \zh{两} (liǎng) obligatoire pour « deux » devant classifieur, pas \zh{二} (èr).
\end{itemize}
\end{attention}

\courssub{Activité d'intégration : Concours « Mon journal numérique »}
\courssub{Situation de communication}
Dans le cadre de la Semaine des Langues Vivantes au Lycée de Biyem-Assi (Yaoundé), le club de chinois lance le concours \zh{我和网络} (« Internet et moi »). Chaque élève de Terminale A rédige son journal personnel en chinois simplifié. Les 3 meilleurs seront lus en public et affichés aux côtés d'un article sur la jeunesse chinoise connectée.

\begin{textebox}[Modèle fourni : Journal de Li Ming (Pékin)]
\zh{我叫李明，我是北京的一名高中生。我每天用手机上网，一般晚上上两个小时。我用微信和朋友聊天，也在网上看新闻、听音乐。网络很方便，因为我可以在网上学习汉语和英语，所以我的汉语越来越好了。但是，上网太久对眼睛不好，也容易浪费时间。以后，我要少玩游戏，多学习。}\\
Wǒ jiào Lǐ Míng, wǒ shì Běijīng de yī míng gāozhōngshēng. Wǒ měi tiān yòng shǒujī shàngwǎng, yībān wǎnshang shàng liǎng ge xiǎoshí. Wǒ yòng Wēixìn hé péngyou liáotiān, yě zài wǎngshang kàn xīnwén, tīng yīnyuè. Wǎngluò hěn fāngbiàn, yīnwèi wǒ kěyǐ zài wǎngshang xuéxí Hànyǔ hé Yīngyǔ, suǒyǐ wǒ de Hànyǔ yuè lái yuè hǎo le. Dànshì, shàngwǎng tài jiǔ duì yǎnjing bù hǎo, yě róngyì làngfèi shíjiān. Yǐhòu, wǒ yào shǎo wán yóuxì, duō xuéxí.\\
« Je m'appelle Li Ming, lycéen à Pékin. Je me connecte chaque jour au portable, en général 2h le soir. J'utilise WeChat pour discuter, je lis les infos et j'écoute de la musique en ligne. Internet est pratique : comme j'étudie le chinois et l'anglais en ligne, mon chinois s'améliore de plus en plus. Mais trop longtemps en ligne nuit aux yeux et fait perdre du temps. À l'avenir, je jouerai moins, j'étudierai plus. »
\end{textebox}

\courssub{Consignes de l'activité (5 tâches)}

\begin{methode}[Tâche 1 — Compréhension du modèle (individuel, 5 min)]
Lisez le journal de Li Ming deux fois. Répondez par écrit en chinois :
\begin{enumerate}
\item \zh{李明用什么上网？}
\item \zh{他在网上做什么？(至少3 activités)}
\item \zh{上网太久有什么坏处？(deux points)}
\end{enumerate}
\textit{Vérification orale collective : l'enseignant attend les réponses en phrases complètes.}
\end{methode}

\begin{methode}[Tâche 2 — Préparation lexicale personnelle (cahier, 5 min)]
Complétez cette trame avec VOS informations, en piochant dans le glossaire :
\begin{enumerate}
\item Fréquence : \zh{我每天 / 常常 / 有时候 / 很少上网。}
\item Appareil \& Applis : \zh{我用手机 / 电脑上网。我用微信 / WhatsApp / Facebook……}
\item Activités : \zh{我在网上看新闻 / 看视频 / 听音乐 / 聊天 / 学习汉语 / 找信息。}
\item Avantage : \zh{网络很方便 / 有好处，因为……所以……}
\item Danger : \zh{但是，上网太久对眼睛不好 / 浪费时间 / 有危险。}
\item Résolution : \zh{以后，我要少……多……}
\end{enumerate}
\end{methode}

\begin{methode}[Tâche 3 — Rédaction guidée : \zh{我和网络} (15 min)]
Rédigez votre journal (6--8 phrases) en respectant \textbf{obligatoirement} cet ordre :
\begin{enumerate}
\item Présentation : Nom, ville (Yaoundé, Douala, Buea\ldots), classe (\zh{高中生}).
\item Fréquence : \zh{我每天\ldots} / \zh{我通常\ldots}
\item Appareil \& Réseaux : \zh{我用\ldots 上网。我用\ldots}
\item Activités : \zh{我在网上\ldots{}、\ldots{}、\ldots{}。}
\item Avantage (cause-conséquence) : \zh{网络很方便/有好处，因为\ldots{}，所以\ldots{}。}
\item Danger (opposition) : \zh{但是，\ldots}
\item Résolution : \zh{以后，我要\ldots}
\end{enumerate}
\textbf{Contraintes :} Minimum 6 mots du glossaire. Au moins \zh{因为……所以……} OU \zh{但是}. Caractères soignés, pinyin noté au-dessus si besoin.
\end{methode}

\begin{experience}[Tâche 4 — Échange en binômes / Pair-work (10 min)]
\begin{enumerate}
\item Échangez votre feuille avec votre voisin.
\item Lisez son texte à voix basse. Cochez les mots du glossaire utilisés.
\item Posez-vous \textbf{deux questions} de compréhension (ex. : \zh{你用什么上网？} / \zh{你每天上几个小时？} / \zh{你觉得网络有什么好处？}).
\item Corrigez ensemble : caractères mal formés, tons hésitants, ordre des mots (\zh{用手机上网}).
\item Signez la feuille de votre camarade : \zh{阅读人：[Votre nom]}.
\end{enumerate}
\end{experience}

\begin{experience}[Tâche 5 — Lecture publique et vote (10 min)]
\begin{enumerate}
\item 3 volontaires lisent leur journal devant la classe (debout, voix claire, tons soignés).
\item La classe note sur un papier : les 3 mots du thème les plus entendus + une question pour l'auteur.
\item Vote à main levée pour le « Prix de la meilleure plume », « Prix de la meilleure prononciation », « Prix de la réflexion la plus responsable ».
\item Affichage des 3 lauréats au tableau du club.
\end{enumerate}
\end{experience}

\courssub{Critères d'évaluation (Grille enseignant)}
\begin{aretenir}[Grille d'évaluation — Sur 10 points]
\begin{center}
\begin{tabularx}{\linewidth}{|X|X|X|}
\hline
\rowcolor{popblueL} Critère & Indicateurs & Points \\
\hline
Vocabulaire thème & $\ge$ 6 mots du glossaire employés correctement (caractères + sens) & 3 \\
\hline
Grammaire / Connecteurs & \zh{因为……所以……} et/ou \zh{但是} bien construits ; ordre mots (用\ldots 上网 ; 在网上\ldots) & 3 \\
\hline
Prononciation / Tons & Lecture fluide, tons justes sur mots-clés (\zh{上网}, \zh{手机}, \zh{微信}, \zh{危险}, \zh{因为}) & 2 \\
\hline
Contenu / Longueur & 6--8 phrases ; structure respectée ; résolution personnelle cohérente & 2 \\
\hline
\end{tabularx}
\end{center}
\end{aretenir}

\courssub{Production modèle et corrigé commenté}
\begin{exoresolu}[Journal numérique modèle : \zh{我和网络} (Élève Terminale A, Yaoundé)]
Rédigez votre journal en 6--8 phrases suivant la trame guidée.
\tcblower
\textbf{Production de l'élève (Aimée, Lycée de Biyem-Assi) :}\\[0.3em]
\zh{我叫艾美丽，我是雅温得的一名高中生。我每天用手机上网，一般晚上上三个小时。我用微信和WhatsApp跟朋友聊天，也在网上看新闻、看视频。网络很方便，因为我可以在网上学习汉语，所以我的汉语越来越好了。但是，上网太久对眼睛不好，也浪费很多时间。以后，我要少看手机，多看书，多运动。}\\[0.3em]
Wǒ jiào Àiměilì, wǒ shì Yǎwēndé de yī míng gāozhōngshēng. Wǒ měi tiān yòng shǒujī shàngwǎng, yībān wǎnshang shàng sān ge xiǎoshí. Wǒ yòng Wēixìn hé WhatsApp gēn péngyou liáotiān, yě zài wǎngshang kàn xīnwén, kàn shìpín. Wǎngluò hěn fāngbiàn, yīnwèi wǒ kěyǐ zài wǎngshang xuéxí Hànyǔ, suǒyǐ wǒ de Hànyǔ yuè lái yuè hǎo le. Dànshì, shàngwǎng tài jiǔ duì yǎnjing bù hǎo, yě làngfèi hěn duō shíjiān. Yǐhòu, wǒ yào shǎo kàn shǒujī, duō kàn shū, duō yùndòng.\\[0.3em]
« Je m'appelle Aimée, je suis lycéenne à Yaoundé. Je me connecte chaque jour au portable, en général 3h le soir. J'utilise WeChat et WhatsApp pour discuter avec mes amis, et je regarde les infos et des vidéos en ligne. Internet est très pratique : comme je peux étudier le chinois en ligne, mon chinois s'améliore de plus en plus. Mais rester trop longtemps en ligne est mauvais pour les yeux et fait perdre beaucoup de temps. À l'avenir, je dois moins regarder mon téléphone, plus lire, plus faire du sport. »\\[0.5em]
\textbf{Commentaire linguistique détaillé :}
\begin{itemize}
\item \textbf{Lexique (9/15 mots glossaire) :} \zh{上网, 手机, 微信, 聊天, 新闻, 视频, 网络, 方便, 浪费, 眼睛, 以后, 少, 多} — critère largement rempli. Note : \zh{WhatsApp} conservé en alphabet latin (usage admis pour noms propres étrangers non sinisés).
\item \textbf{Grammaire / Connecteurs :} Phrase 5 : structure complète \zh{因为……所以……} (cause : étude en ligne $\rightarrow$ conséquence : progrès). Phrase 6 : opposition nette avec \zh{但是} introduisant deux inconvénients (\zh{对眼睛不好}, \zh{浪费时间}). Phrase 7 : résolution \zh{以后，我要……} + antithèse \zh{少看手机，多看书，多运动} (structure \zh{少V, 多V} très idiomatique).
\item \textbf{Ordre des mots (Syntax) :}
  \begin{itemize}
  \item Moyen : \zh{用手机上网} (Sujet + 用 + Instrument + Verbe) — \checkmark correct.
  \item Lieu : \zh{在网上看新闻} (Sujet + 在 + Lieu + Verbe) — \checkmark correct.
  \item Durée : \zh{上三个小时} (Verbe + Durée) — \checkmark correct (post-verbal).
  \end{itemize}
\item \textbf{Structure progressive :} \zh{越来越好} (yuè lái yuè hǎo) « de mieux en mieux » + \zh{了} (particule d'accomplissement / changement d'état) — niveau HSK 3 maîtrisé.
\item \textbf{Adaptation contextuelle :} Citation de Yaoundé (\zh{雅温得}), double usage WeChat/WhatsApp (réalité camerounaise), résolution santé (\zh{多运动}) pertinente.
\end{itemize}
\end{exoresolu}

\courssub{Bilan de la leçon et ouverture socioculturelle}
\begin{savaistu}[Cameroun — Chine : Deux jeunesses connectées]
\begin{itemize}
\item \textbf{Pénétration mobile :} Au Cameroun (2024), $\approx$ 95\% des connexions Internet se font via mobile (Données ARTC/GSMA). En Chine, 99,9\% des internautes utilisent le smartphone (CNNIC 2024). Le \zh{手机} est l'outil unique dominant dans les deux pays.
\item \textbf{Écosystème d'applications :}
  \begin{itemize}
  \item \textbf{Cameroun :} WhatsApp (messagerie n°1), Facebook / Facebook Lite, YouTube, TikTok (Douyin version intl), Instagram. Paiement mobile : Mobile Money (MTN MoMo, Orange Money) — \zh{手机支付} (shǒujī zhīfù) très développé.
  \item \textbf{Chine :} \zh{微信} (WeChat) = « super-app » tout-en-un (messagerie, paiement \zh{微信支付}, réseaux sociaux \zh{朋友圈}, mini-programmes, admin). \zh{抖音} (Douyin) = TikTok version chinoise. \zh{支付宝} (Alipay) = concurrent paiement. \textbf{Grand Firewall} : Google, WhatsApp, Facebook, YouTube, Instagram sont bloqués ; les Chinois utilisent des équivalents locaux.
  \end{itemize}
\item \textbf{Usage éducatif :} Les deux pays ont massivement basculé vers l'\zh{网课} (cours en ligne) pendant le Covid-19. Au Cameroun, WhatsApp groupes classe + Google Classroom ; en Chine, \zh{钉钉} (DingTalk), \zh{腾讯会议} (Tencent Meeting), \zh{学习强国} (Xuexi Qiangguo).
\item \textbf{Régulation / Santé :} Chine : loi 2021 limitant les jeux vidéo en ligne pour les mineurs à 3h/semaine (vendredi-samedi-dimanche 20h-21h) + systèmes anti-addiction (reconnaissance faciale). Cameroun : pas de loi similaire, mais campagnes \zh{护眼} (protection vue) et \zh{网络安全} (cybersécurité) dans les lycées (MINESUP/MINESEC).
\item \textbf{Vocabulaire comparatif :}
\begin{center}
\begin{tabularx}{\linewidth}{|X|X|X|}
\hline
\rowcolor{popblueL} Notion & Cameroun (Français / Pidgin) & Chine (Chinois / Pinyin) \\
\hline
Forfait data & Forfait / « Data » & \zh{流量包} (liúliàng bāo) \\
\hline
Recharger & Recharger / « Credit » & \zh{充值} (chōngzhí) \\
\hline
Payer par téléphone & Mobile Money / MoMo & \zh{扫码支付} (sǎomǎ zhīfù) — scanner QR code \\
\hline
Fake news & « Fake news » / « Bobard » & \zh{谣言} (yáoyán) / \zh{假新闻} (jiǎ xīnwén) \\
\hline
Cyber-harcèlement & Cyber-harcèlement & \zh{网络暴力} (wǎngluò bàolì) \\
\hline
\end{tabularx}
\end{center}
\end{itemize}
\textbf{Piste de débat oral (niveau Terminale) :} \zh{「网络让世界更近，还是让人更远？」} (« Internet rapproche-t-il le monde ou éloigne-t-il les gens ? ») — Utiliser \zh{因为……所以……}, \zh{但是}, \zh{比如} (par exemple), \zh{我认为} (je pense que).
\end{savaistu}

\courssub{Devoirs et prolongements}
\begin{exercice}[Devoirs obligatoires]
\begin{enumerate}
\item \textbf{Réécriture propre :} Recopiez votre journal \zh{我和网络} corrigé sur une feuille A5 (caractères soignés, pinyin au-dessus des mots nouveaux). À rendre au prochain cours pour affichage.
\item \textbf{Exercices de fixation (cahier) :}
  \begin{enumerate}
  \item Traduisez en chinois : « Je me connecte souvent avec mon ordinateur pour regarder des vidéos. » (\zh{我常用电脑上网看视频。})
  \item Traduisez : « Internet est pratique, mais c'est dangereux pour les yeux. » (\zh{网络很方便，但是对眼睛有危险。})
  \item Complétez avec \zh{因为} ou \zh{所以} : \zh{\_\_\_\_\_\_ 我想学好汉语，\_\_\_\_\_\_ 我每天看中文视频。} (\zh{因为\ldots 所以\ldots})
  \item Écrivez 3 phrases avec \zh{以后，我要\ldots} sur vos résolutions personnelles (études, sport, écrans).
  \end{enumerate}
\item \textbf{Recherche culturelle (optionnel, bonus) :} Trouvez le nombre d'utilisateurs de WeChat en 2024 (chiffre officiel Tencent) et comparez avec la population du Cameroun. Présentez en 3 phrases chinoises la semaine prochaine.
\end{enumerate}
\end{exercice}

\begin{corrige}[Devoirs — Éléments de correction]
\begin{enumerate}
\item \zh{我常用电脑上网看视频。} (Wǒ cháng yòng diànnǎo shàngwǎng kàn shìpín.)
\item \zh{网络很方便，但是对眼睛有危险。} (Wǎngluò hěn fāngbiàn, dànshì duì yǎnjing yǒu wēixiǎn.) — ou \zh{……但是上网太久对眼睛不好。}
\item \zh{因为我想学好汉语，所以我每天看中文视频。} (Yīnwèi wǒ xiǎng xué hǎo Hànyǔ, suǒyǐ wǒ měi tiān kàn Zhōngwén shìpín.)
\item Exemples : \zh{以后，我要少玩游戏，多复习。} / \zh{以后，我要每天跑步。} / \zh{以后，我要早睡早起。}
\end{enumerate}
\end{corrige}

\newpage

\coursec{Méthodes et exercices résolus}

\coursec{Leçon 34 — Vocabulaire et compréhension de texte : internet et les réseaux sociaux}

\courssub{Mise en route et découverte du thème}

\begin{experience}[Brainstorming : ma vie connectée]
En classe entière, le professeur projette ou écrit au tableau les mots-clés : \zh{手机} \emph{shǒujī}, \zh{上网} \emph{shàngwǎng}, \zh{微信} \emph{Wēixìn}, \zh{视频} \emph{shìpín}, \zh{发消息} \emph{fā xiāoxi}.
\begin{enumerate}
\item \textbf{Association libre} : chaque élève dit en français ce que ces mots lui évoquent (loisirs, études, famille, dangers).
\item \textbf{Mise en pinyin} : l'enseignant écrit le pinyin avec les tons au tableau ; la classe répète en chœur en exagérant les tons (3ᵉ ton \zh{网} wǎng : bas-montant ; 4ᵉ ton \zh{信} xìn : haut-bas sec).
\item \textbf{Contexte camerounais} : « Combien d'heures par jour passez-vous sur \zh{手机} ? Quelles applications utilisez-vous à Yaoundé, Douala, Buea ? (\zh{WhatsApp}, \zh{Facebook}, \zh{TikTok}, \zh{Orange Money}, \zh{MTN MoMo}) ».
\end{enumerate}
\emph{Objectif : activer le schéma mental, lier le lexique au vécu des élèves, préparer la lecture.}
\end{experience}

\begin{popfigure}
\begin{tikzpicture}[mindmap, grow cyclic, every node/.style=concept, concept color=popblue!60,
    level 1/.append style={level distance=4.5cm,sibling angle=90},
    level 2/.append style={level distance=3cm,sibling angle=45}]
\node[concept color=popblue] {Internet \& Réseaux}
    child[concept color=popgreen] { node[concept, align=center]{Matériel\\ \zh{手机 shǒujī}\\ \zh{电脑 diànnǎo}} }
    child[concept color=poporange] { node[concept, align=center]{Actions\\ \zh{上网 shàngwǎng}\\ \zh{发消息 fā xiāoxi}\\ \zh{聊天 liáotiān}\\ \zh{看视频 kàn shìpín}\\ \zh{下载 xiàzǎi}} }
    child[concept color=poppurple] { node[concept, align=center]{Contenu\\ \zh{视频 shìpín}\\ \zh{照片 zhàopiàn}\\ \zh{新闻 xīnwén}\\ \zh{消息 xiāoxi}} }
    child[concept color=poppink] { node[concept, align=center]{Plateformes\\ \zh{微信 Wēixìn}\\ \zh{朋友圈 péngyouquān}\\ \zh{网站 wǎngzhàn}\\ \zh{抖音 Dǒuyīn}} };
\end{tikzpicture}
\legende{Carte mentale du vocabulaire du thème (HSK 3-4).}
\end{popfigure}

\courssub{Texte support : lecture et compréhension globale}

\begin{textebox}[Texte original : Internet et les réseaux sociaux au quotidien]
现在，很多喀麦隆年轻人每天都上网。他们用手机看新闻、发消息、看视频。\\
Xiànzài, hěn duō Kāmàilóng niánqīngrén měi tiān dōu shàngwǎng. Tāmen yòng shǒujī kàn xīnwén, fā xiāoxi, kàn shìpín.\\
« Aujourd'hui, beaucoup de jeunes Camerounais se connectent chaque jour. Avec leur téléphone, ils lisent les informations, envoient des messages et regardent des vidéos. »\\[4pt]
社交媒体很有用：人们可以和朋友聊天，也可以学习汉语。但是，有些人上网时间太长，对身体不好，对眼睛也不好。\\
Shèjiāo méitǐ hěn yǒuyòng : rénmen kěyǐ hé péngyou liáotiān, yě kěyǐ xuéxí Hànyǔ. Dànshì, yǒuxiē rén shàngwǎng shíjiān tài cháng, duì shēntǐ bù hǎo, duì yǎnjing yě bù hǎo.\\
« Les réseaux sociaux sont très utiles : on peut discuter avec ses amis et aussi apprendre le chinois. Mais certaines personnes passent trop de temps en ligne, ce qui est mauvais pour le corps et pour les yeux. »\\[4pt]
在雅温得和杜阿拉，越来越多的学生用互联网来学习。我觉得网络很方便，可是我们应该少玩手机，多运动。\\
Zài Yǎwēndé hé Dù'ālā, yuè lái yuè duō de xuésheng yòng hùliánwǎng lái xuéxí. Wǒ juéde wǎngluò hěn fāngbiàn, kěshì wǒmen yīnggāi shǎo wán shǒujī, duō yùndòng.\\
« À Yaoundé et à Douala, de plus en plus d'élèves utilisent internet pour étudier. Je trouve le réseau très pratique, mais nous devrions moins jouer avec le téléphone et faire plus de sport. »
\end{textebox}

\begin{methode}[Lire et comprendre un court texte chinois (Bac)]
Pour réussir les questions de compréhension, suis cette démarche en 5 étapes :
\begin{enumerate}
\item \textbf{Lecture silencieuse globale} : lis une fois sans dictionnaire. Repère le thème (\zh{互联网} hùliánwǎng, \zh{社交媒体} shèjiāo méitǐ) et les mots connus.
\item \textbf{Repérage des mots-clés du thème} : souligne le vocabulaire cible : \zh{上网} shàngwǎng, \zh{手机} shǒujī, \zh{发消息} fā xiāoxi, \zh{视频} shìpín, \zh{聊天} liáotiān, \zh{学习} xuéxí, \zh{有用} yǒuyòng, \zh{方便} fāngbiàn, \zh{越来越多} yuè lái yuè duō.
\item \textbf{Découpage des phrases} : identifie le \textbf{Sujet} (souvent en début de phrase), les \textbf{compléments circonstanciels} (temps, lieu, instrument, destinataire) placés \textbf{avant le verbe}, puis le \textbf{Verbe} et son \textbf{complément d'objet}.
\item \textbf{Lecture des questions avant la relecture} : lis les questions, puis relis le texte en cherchant les passages correspondants. La réponse suit l'ordre du texte.
\item \textbf{Rédaction de la réponse} : réponds par une phrase complète en chinois, en reprenant les mots du texte, sans recopier le paragraphe entier.
\end{enumerate}
\cle{Réflexe Bac} : les mots de négation (\zh{不}, \zh{没有}, \zh{不好}) et les adverbes de fréquence/degré (\zh{常常}, \zh{有时候}, \zh{很少}, \zh{太}, \zh{越来越}) changent radicalement le sens : repère-les en priorité.
\end{methode}

\courssub{Glossaire du thème : vocabulaire en tableau}

\begin{center}
\begin{tabularx}{\linewidth}{|X|X|X|X|}
\hline
\rowcolor{popblueL} \textbf{Caractères} & \textbf{Pinyin} & \textbf{Nature} & \textbf{Traduction} \\
\hline
互联网 & hùliánwǎng & nom & internet (litt. « réseau interconnecté ») \\
\hline
社交媒体 & shèjiāo méitǐ & nom & réseaux sociaux, médias sociaux \\
\hline
手机 & shǒujī & nom & téléphone portable (smartphone) \\
\hline
电脑 & diànnǎo & nom & ordinateur \\
\hline
上网 & shàngwǎng & verbe (séparable) & se connecter, aller sur internet \\
\hline
发消息 & fā xiāoxi & verbe + objet & envoyer un message \\
\hline
看视频 & kàn shìpín & verbe + objet & regarder des vidéos \\
\hline
聊天 & liáotiān & verbe (séparable) & bavarder, discuter (en ligne) \\
\hline
下载 & xiàzǎi & verbe & télécharger \\
\hline
网站 & wǎngzhàn & nom & site web \\
\hline
视频 & shìpín & nom & vidéo \\
\hline
照片 & zhàopiàn & nom & photo \\
\hline
朋友圈 & péngyouquān & nom & cercle d'amis (fil d'actualité WeChat) \\
\hline
微信 & Wēixìn & nom propre & WeChat (appli dominante en Chine) \\
\hline
有用 & yǒuyòng & adj. & utile \\
\hline
方便 & fāngbiàn & adj. & pratique, commode \\
\hline
不好 & bù hǎo & adj. (nég.) & mauvais, pas bien \\
\hline
身体 & shēntǐ & nom & corps, santé physique \\
\hline
眼睛 & yǎnjing & nom & yeux \\
\hline
时间 & shíjiān & nom & temps \\
\hline
太长 & tài cháng & adj. + adv. & trop long \\
\hline
越来越多 & yuè lái yuè duō & structure & de plus en plus nombreux \\
\hline
感兴趣 & gǎn xìngqù & verbe statique & s'intéresser à (structure \zh{对}…\zh{感兴趣}) \\
\hline
雅温得 & Yǎwēndé & nom propre & Yaoundé (capitale du Cameroun) \\
\hline
杜阿拉 & Dù'ālā & nom propre & Douala (capitale économique) \\
\hline
应该 & yīnggāi & verbe modal & devoir, devrait \\
\hline
少 & shǎo & adv./verbe & moins, peu / réduire \\
\hline
多 & duō & adv./verbe & plus, beaucoup / augmenter \\
\hline
运动 & yùndòng & verbe/nom & faire du sport / sport \\
\hline
觉得 & juéde & verbe & penser, trouver (avis subjectif) \\
\hline
\end{tabularx}
\end{center}

\courssub{Radicaux (clés) et ordre des traits}

\begin{propriete}[Radicaux et classification des caractères clés]
Connaître le radical (clé) aide à mémoriser le sens et à retrouver le caractère dans un dictionnaire.
\begin{center}
\begin{tabularx}{\linewidth}{|X|X|X|X|}
\hline
\rowcolor{popblueL} \textbf{Caractère} & \textbf{Radical (clé)} & \textbf{Sens du radical} & \textbf{Indice mnémotechnique} \\
\hline
网 wǎng & \zh{纟} (sī, soie) & soie, fil, tissu & Un filet est fait de fils entrecroisés. Forme simplifiée de \zh{網}. \\
\hline
手 shǒu & \zh{手} (shǒu, main) & main, action manuelle & Clé de la main. Dans \zh{机} (machine), on trouve \zh{木} (bois) + phonétique. \\
\hline
机 jī & \zh{木} (mù, bois) & bois, arbre, matière & Machine = mécanisme en bois (anciennement) + phonétique \zh{幾}. \\
\hline
视 shì & \zh{见} (jiàn, voir) & voir, regard & Voir avec attention (radical \zh{见} en bas). \\
\hline
频 pín & \zh{页} (yè, page/tête) & tête, page & Fréquence = ce qui arrive souvent à la tête (phonétique \zh{品}). \\
\hline
交 jiāo & \zh{交} (jiāo, croiser) & croiser, échanger & Deux jambes qui se croisent. \\
\hline
媒 méi & \zh{女} (nǚ, femme) & femme, féminin & Médium, intermédiaire (phonétique \zh{每}). \\
\hline
信 xìn & \zh{亻} (rén, homme) & personne, humain & Message porté par un homme (phonétique \zh{言} parole). \\
\hline
息 xī & \zh{心} (xīn, cœur) & cœur, souffle, information & Information = souffle/nouvelle qui vient du cœur. \\
\hline
\end{tabularx}
\end{center}
\end{propriete}

\begin{popfigure}
\begin{tikzpicture}[
    strokebox/.style={draw=popdark, thick, rectangle, minimum size=1.2cm, inner sep=2pt},
    labelnode/.style={font=\small, text=popdark, anchor=north},
    ordernode/.style={font=\large\bfseries, text=poporange, anchor=center}
]
% Caractère 网 (6 traits)
\matrix (wang) [matrix of nodes, nodes={strokebox}, column sep=2mm, row sep=2mm, label=above:\textbf{网 wǎng (6 traits)}] {
|[ordernode]| 1 & |[ordernode]| 2 & |[ordernode]| 3 & |[ordernode]| 4 & |[ordernode]| 5 & |[ordernode]| 6 \\
\zh{一} & \zh{丨} & \zh{一} & \zh{一} & \zh{丨} & \zh{一} \\
};
% Caractère 信 (9 traits)
\matrix (xin) [matrix of nodes, nodes={strokebox}, column sep=2mm, row sep=2mm, below=1.5cm of wang, label=above:\textbf{信 xìn (9 traits)}] {
|[ordernode]| 1 & |[ordernode]| 2 & |[ordernode]| 3 & |[ordernode]| 4 & |[ordernode]| 5 & |[ordernode]| 6 & |[ordernode]| 7 & |[ordernode]| 8 & |[ordernode]| 9 \\
\zh{亻} & \zh{丶} & \zh{一} & \zh{一} & \zh{一} & \zh{丨} & \zh{丶} & \zh{丿} & \zh{乀} \\
};
\node[labelnode] at (wang-2-1.south) {横};
\node[labelnode] at (wang-2-2.south) {竖};
\node[labelnode] at (wang-2-3.south) {横};
\node[labelnode] at (wang-2-4.south) {横};
\node[labelnode] at (wang-2-5.south) {竖};
\node[labelnode] at (wang-2-6.south) {横 (fermeture)};
\node[labelnode] at (xin-2-1.south) {亻 (pied gauche)};
\node[labelnode] at (xin-2-2.south) {点};
\node[labelnode] at (xin-2-3.south) {横};
\node[labelnode] at (xin-2-4.south) {横};
\node[labelnode] at (xin-2-5.south) {横};
\node[labelnode] at (xin-2-6.south) {竖};
\node[labelnode] at (xin-2-7.south) {点};
\node[labelnode] at (xin-2-8.south) {撇};
\node[labelnode] at (xin-2-9.south) {捺};
\end{tikzpicture}
\legende{Ordre des traits pour \zh{网} (haut→bas, gauche→droite, horizontales avant verticales, fermeture finale) et \zh{信} (clé \zh{亻} d'abord, puis \zh{言} simplifié).}
\end{popfigure}

\courssub{Questions de compréhension du texte}

\begin{exercice}[Compréhension écrite : répondre en chinois]
Lis le texte support (page précédente) et réponds aux questions suivantes par des phrases complètes en caractères chinois.
\begin{enumerate}
\item 喀麦隆年轻人用什么上网？
\item 社交媒体有什么用？(Donne les deux usages mentionnés.)
\item 上网时间太长好不好？为什么？
\item 在雅温得和杜阿拉，谁用互联网来学习？
\item 根据最后一句，我们应该怎么做？(Deux actions.)
\end{enumerate}
\end{exercice}

\begin{corrige}[Exercice Compréhension écrite]
\textbf{1.} 他们用手机上网。 \emph{Tāmen yòng shǒujī shàngwǎng.} (Ils se connectent avec leur téléphone.)\\
\textbf{2.} 人们可以和朋友聊天，也可以学习汉语。 \emph{Rénmen kěyǐ hé péngyou liáotiān, yě kěyǐ xuéxí Hànyǔ.} (On peut discuter avec des amis et aussi apprendre le chinois.)\\
\textbf{3.} 不好，因为上网时间太长对身体不好，对眼睛也不好。 \emph{Bù hǎo, yīnwèi shàngwǎng shíjiān tài cháng duì shēntǐ bù hǎo, duì yǎnjing yě bù hǎo.} (Non, car passer trop de temps en ligne est mauvais pour le corps et les yeux.)\\
\textbf{4.} 越来越多的学生用互联网来学习。 \emph{Yuè lái yuè duō de xuésheng yòng hùliánwǎng lái xuéxí.} (De plus en plus d'élèves utilisent internet pour étudier.)\\
\textbf{5.} 我们应该少玩手机，多运动。 \emph{Wǒmen yīnggāi shǎo wán shǒujī, duō yùndòng.} (Nous devrions moins jouer avec le téléphone et faire plus de sport.)
\end{corrige}

\courssub{Production guidée et comparaison Cameroun -- Chine}

\begin{methode}[Rédiger un court paragraphe structuré (5-6 phrases)]
Modèle type pour l'expression écrite (Bac) :
\begin{enumerate}
\item \textbf{Phrase d'accroche (généralité)} : \zh{现在，…} (Aujourd'hui…) / \zh{随着科技的发展，…} (Avec le développement de la tech…).
\item \textbf{Usages concrets (2 phrases)} : \zh{人们用手机…} / \zh{我用微信…} / \zh{学生上网学习…}. Respecte l'ordre : Sujet + (Temps/Lieu/Instrument) + Verbe + Objet.
\item \textbf{Avantage / Utilité} : \zh{…很有用 / 很方便} ; structure \zh{可以…也可以…}.
\item \textbf{Limite / Problème (contraste)} : \zh{但是 / 可是} + \zh{上网时间太长对眼睛/身体不好}.
\item \textbf{Avis personnel / Conseil} : \zh{我觉得 / 我认为} + \zh{我们应该少…多…}.
\end{enumerate}
\cle{Connecteurs logiques} : \zh{也} (aussi), \zh{但是/可是} (mais), \zh{因为…所以…} (parce que… donc), \zh{越来越} (de plus en plus).
\end{methode}

\begin{exercice}[Expression écrite : ton avis]
Rédige un paragraphe de 5 à 6 phrases en chinois sur le thème : « Internet et les réseaux sociaux dans ma vie ».
\textbf{Mots imposés à utiliser} : \zh{上网}、\zh{发消息}、\zh{有用}、\zh{但是}、\zh{我觉得}、\zh{应该}.
\end{exercice}

\begin{corrige}[Exercice Expression écrite -- Modèle de réponse]
现在，我每天都上网。我用手机看新闻，也给朋友发消息。社交媒体很有用：我可以和同学聊天，也可以学习汉语。但是，上网时间太长对眼睛不好。我觉得网络很方便，可是我们应该少玩手机，多运动。\\
\emph{Xiànzài, wǒ měi tiān dōu shàngwǎng. Wǒ yòng shǒujī kàn xīnwén, yě gěi péngyou fā xiāoxi. Shèjiāo méitǐ hěn yǒuyòng: wǒ kěyǐ hé tóngxué liáotiān, yě kěyǐ xuéxí Hànyǔ. Dànshì, shàngwǎng shíjiān tài cháng duì yǎnjing bù hǎo. Wǒ juéde wǎngluò hěn fāngbiàn, kěshì wǒmen yīnggāi shǎo wán shǒujī, duō yùndòng.}\\
\textit{Traduction : Aujourd'hui, je me connecte tous les jours. Je lis les infos sur mon téléphone et j'envoie aussi des messages à mes amis. Les réseaux sociaux sont très utiles : je peux discuter avec mes camarades et aussi apprendre le chinois. Mais passer trop de temps en ligne est mauvais pour les yeux. Je trouve internet très pratique, mais nous devrions moins jouer avec le téléphone et faire plus de sport.}
\end{corrige}

\begin{savaistu}[Cameroun -- Chine : deux mondes connectés]
\begin{itemize}
\item \textbf{Super-applications} : En Chine, \zh{微信} (Wēixìn / WeChat) est un écosystème tout-en-un : messagerie, \zh{朋友圈} (Moments), paiement mobile (\zh{微信支付}), services publics, taxi, rendez-vous médicaux. Au Cameroun, on utilise plusieurs applis distinctes : \zh{WhatsApp}/\zh{Facebook} (social), \zh{Orange Money}/\zh{MTN MoMo} (paiement), navigateur web (info).
\item \textbf{Paiement mobile} : La Chine est quasi « sans espèces » (\zh{无现金社会}) via codes QR (\zh{二维码} èrwéimǎ). Au Cameroun, le Mobile Money (USSD \#150\# ou applis) domine l'inclusion financière, mais le cash reste roi dans l'informel.
\item \textbf{Accès internet} : Chine : 5G généralisée, fibre dans toutes les villes, \zh{网速} (wǎngsù) très haut. Cameroun : 4G dans les grandes villes (Yaoundé, Douala, Bafoussam, Garoua), 3G/Edge en zone rurale, fibre en déploiement (Camtel, MTN, Orange). Le coût du data (\zh{流量} liúliàng) reste élevé par rapport au revenu moyen.
\item \textbf{Régulation} : Chine : « Grand Firewall » (\zh{防火长城}), censure politique, réseau national (\zh{内网}) vs internet mondial. Cameroun : internet relativement ouvert, coupures ponctuelles lors de crises politiques (ex. 2017-2018 Anglophone regions), loi sur la cybercriminalité (2010).
\item \textbf{Usage éducatif} : Chine : \zh{网课} (wǎngkè) massifs (COVID), plateformes \zh{学习强国} (Xuéxí Qiángguó), \zh{哔哩哔哩} (Bilibili). Cameroun : développement de l'e-learning (Moodle, Google Classroom, WhatsApp groups pour cours), défi de la fracture numérique (équipement, électricité, data).
\end{itemize}
\cle{Pour l'oral (Bac)} : sois capable de dire : « En Chine, on paie tout avec le téléphone (\zh{扫码支付} sǎomǎ zhīfù), au Cameroun on utilise Orange Money. Les Chinois utilisent WeChat, nous on utilise WhatsApp. »
\end{savaistu}

\begin{experience}[Débat oral guidé : « Le téléphone portable à l'école : outil ou distraction ? »]
\begin{enumerate}
\item \textbf{Préparation (5 min)} : En binôme, préparez 3 arguments POUR (outil : dictionnaire \zh{词典} cídiǎn, apps éducatives \zh{学习软件} xuéxí ruǎnjiàn, communication parents) et 3 arguments CONTRE (distraction \zh{分心} fēnxīn, triche \zh{作弊} zuòbì, santé \zh{近视} jìnshì myopie, cyberharcèlement \zh{网络霸凌} wǎngluò bàlíng).
\item \textbf{Vocabulaire utile} : \zh{允许} yǔxǔ (autoriser), \zh{禁止} jìnzhǐ (interdire), \zh{管理} guǎnlǐ (gérer/encadrer), \zh{自律} zìlǜ (autodiscipline), \zh{利大于弊} lì dà yú bì (avantages > inconvénients).
\item \textbf{Débat (10 min)} : Chaque binôme présente sa position en chinois (1 min). La classe vote.
\item \textbf{Synthèse prof} : Rappel des structures \zh{虽然…但是…} (bien que… pourtant…), \zh{不仅…而且…} (non seulement… mais encore…).
\end{enumerate}
\end{experience}

\courssub{Entraînement : traduction version / thème}

\begin{methode}[Traduire : Version (Zh → Fr) et Thème (Fr → Zh)]
\begin{enumerate}
\item \textbf{Version} : 1) Identifie le Sujet. 2) Repère les compléments AVANT le verbe (\zh{用} instrument, \zh{在} lieu, \zh{给} destinataire, \zh{对} objet). 3) Traduis le verbe. 4) Replace les compléments APRÈS le verbe en français.
\item \textbf{Thème} : 1) Identifie le Sujet. 2) Prends les compléments français (avec qui, où, avec quoi, pour qui) et place-les AVANT le verbe chinois avec leur mot-outil (\zh{用}, \zh{在}, \zh{给}, \zh{对}). 3) Conjugue le verbe (pas de conjugaison, mais aspects \zh{了}, \zh{过}, \zh{着}). 4) Vérifie l'ordre : Sujet + Temps + Lieu + Instrument + Destinataire + Verbe + Objet + Complément de résultat/degré.
\item \textbf{Structures pièges} :
\begin{itemize}
\item \zh{对} + Nom + \zh{感兴趣 / 不好 / 有用} (bloc inséparable).
\item \zh{越来越} + Adj (de plus en plus…).
\item \zh{用} + Instrument + Verbe (pas *Verbe + \zh{用} Instrument).
\end{itemize}
\end{enumerate}
\end{methode}

\begin{exoresolu}[Traduction : Version et Thème corrigés]
Énoncé. a) Version : \zh{在杜阿拉，很多学生用手机上网学习汉语，因为很方便。}\\
b) Thème : « À Yaoundé, de plus en plus de jeunes s'intéressent aux vidéos courtes sur TikTok (抖音 Dǒuyīn). »
\tcblower
\textbf{Solution détaillée.}\\
\textbf{a) Version.}\\
Analyse : \zh{在杜阿拉} (Lieu) | \zh{很多学生} (Sujet) | \zh{用手机} (Instrument) | \zh{上网} (Verbe 1) | \zh{学习汉语} (Verbe 2 but) | \zh{因为很方便} (Cause).\\
Traduction : « À Douala, beaucoup d'élèves se connectent avec leur téléphone pour apprendre le chinois, car c'est très pratique. »\\
\textbf{b) Thème.}\\
Raisonnement : Lieu \zh{在雅温得} (tête). Sujet \zh{越来越多的年轻人}. Structure \zh{对}…\zh{感兴趣}. Objet \zh{抖音上的短视频} (vidéos courtes sur Douyin).\\
Réponse : \zh{在雅温得，越来越多的年轻人对抖音上的短视频感兴趣。}\\
\emph{Zài Yǎwēndé, yuè lái yuè duō de niánqīngrén duì Dǒuyīn shàng de duǎn shìpín gǎn xìngqù.}\\
Vérification : \zh{短视频} duǎn shìpín (vidéo courte), \zh{上} shàng (sur/dans), \zh{感兴趣} gǎn xìngqù (s'intéresser à).
\end{exoresolu}

\begin{exercice}[Traduction : à vous de jouer]
Traduis en chinois :
\begin{enumerate}
\item Je regarde souvent des vidéos sur mon téléphone.
\item Ce site web est très utile pour apprendre le chinois.
\item Nous ne devrions pas passer trop de temps sur les réseaux sociaux.
\item Ma mère m'envoie des photos sur WeChat tous les jours.
\end{enumerate}
\end{exercice}

\begin{corrige}[Exercice Traduction]
1. 我常常用手机看视频。 \emph{Wǒ chángcháng yòng shǒujī kàn shìpín.} (Instrument avant verbe.)\\
2. 这个网站学习汉语很有用。 \emph{Zhège wǎngzhàn xuéxí Hànyǔ hěn yǒuyòng.} (Sujet + Verbe/Objet comme thème + commentaire). Ou : \zh{这个网站对学习汉语很有用。} (Structure \zh{对}…\zh{有用}).\\
3. 我们不应该在社交媒体上花太多时间。 \emph{Wǒmen bù yīnggāi zài shèjiāo méitǐ shàng huā tài duō shíjiān.} (\zh{花时间} passer du temps, \zh{在}…\zh{上} sur). Ou plus simple : \zh{我们不应该上网太长时间。}\\
4. 我妈妈每天在微信给我发照片。 \emph{Wǒ māma měi tiān zài Wēixìn gěi wǒ fā zhàopiàn.} (Lieu \zh{在微信} avant verbe, Destinataire \zh{给我} avant verbe).
\end{corrige}

\courssub{Points de vigilance pour le Bac (Résumé)}

\begin{attention}[Les 5 erreurs qui coûtent des points au Bac -- Récapitulatif]
\begin{enumerate}
\item \textbf{Ordre des mots : calque du français.} \textbf{Faux} : \zh{我上网用手机} / \zh{我发消息给朋友} (si pas de verbe principal avant). \textbf{Vrai} : \zh{我用手机上网} / \zh{我给朋友发消息}. Règle d'or : \cle{Compléments (Temps, Lieu, Instrument, Destinataire) AVANT le verbe.}
\item \textbf{Caractères « jumeaux » (confusion visuelle).} \zh{网} wǎng (filet, clé \zh{纟}) ≠ \zh{往} wǎng (vers, clé \zh{彳}) ; \zh{聊} liáo (bavarder, clé \zh{耳} oreille) ≠ \zh{柳} liǔ (saule, clé \zh{木}) ; \zh{信} xìn (message, clé \zh{亻}) ≠ \zh{言} yán (parole, clé \zh{言}) ; \zh{视} shì (voir, clé \zh{见}) ≠ \zh{试} shì (essayer, clé \zh{言}).
\item \textbf{Tons négligés (pinyin oral/écrit).} \zh{wǎng} (3ᵉ, 网) vs \zh{wáng} (2ᵉ, 王) ; \zh{mǎi} (3ᵉ, 买 acheter) vs \zh{mài} (4ᵉ, 卖 vendre) ; \zh{shuì} (4ᵉ, 睡 dormir) vs \zh{shuǐ} (3ᵉ, 水 eau). Un ton faux = un mot faux = 0 point au mot.
\item \textbf{Classificateurs oubliés ou incorrects.} \zh{一部手机} (bù), \zh{一台电脑} (tái), \zh{一条消息} (tiáo), \zh{一张照片} (zhāng), \zh{一个视频} (gè), \zh{一个网站} (gè). Après \zh{这/那/每} + nombre : \zh{这两部手机}.
\item \textbf{Mots-outils mal placés (adverbes, prépositions).} \zh{也/都/只/就} se placent \textbf{avant} le verbe (\zh{我也去}, jamais \zh{我去也}). \zh{每天/常常} avant le verbe (\zh{我每天上网}). \zh{对} A \zh{感兴趣/不好} : A reste entre \zh{对} et l'adjectif (\zh{对身体不好}, jamais \zh{不好对身体}).
\end{enumerate}
\end{attention}

\begin{aretenir}[Check-list de révision Leçon 34]
\begin{itemize}
\item [ ] Je connais le sens, le pinyin (tons) et l'écriture des 25 mots du tableau de vocabulaire.
\item [ ] Je sais identifier les radicaux de \zh{网, 手, 机, 视, 信, 息, 交, 媒} et expliquer leur lien avec le sens.
\item [ ] Je sais répondre par une phrase complète en chinois aux 5 types de questions de compréhension (Qui, Quoi, Comment, Pourquoi, Avis).
\item [ ] Je sais rédiger un paragraphe de 6 phrases avec connecteurs (\zh{但是}, \zh{因为…所以…}, \zh{觉得}, \zh{应该}).
\item [ ] Je sais traduire une phrase Fr→Zh (Thème) en respectant l'ordre Sujet + Compléments + Verbe.
\item [ ] Je peux citer 3 différences concrètes entre l'usage d'internet au Cameroun et en Chine (WeChat vs WhatsApp/MoMo, paiement QR, régulation).
\end{itemize}
\end{aretenir}

\newpage

\coursec{Exercices}

\courssub{Je vérifie mes connaissances}

\begin{exercice}[QCM : le vocabulaire des réseaux sociaux]
Pour chaque question, entoure la bonne réponse.
\begin{enumerate}
\item \zh{朋友圈} (péngyouquān) signifie :
\begin{enumerate}
\item un groupe d'amis au lycée
\item le « cercle d'amis » (fil d'actualité) de WeChat
\item un jeu vidéo en ligne
\end{enumerate}
\item \zh{网友} (wǎngyǒu) signifie :
\begin{enumerate}
\item un ami rencontré sur internet
\item un site web
\item un mot de passe
\end{enumerate}
\item \zh{上网} (shàngwǎng) signifie :
\begin{enumerate}
\item éteindre son téléphone
\item aller sur internet
\item acheter un ordinateur
\end{enumerate}
\item \zh{方便} (fāngbiàn) est :
\begin{enumerate}
\item un verbe qui signifie « envoyer »
\item un adjectif qui signifie « pratique, commode »
\item un nom qui signifie « téléphone »
\end{enumerate}
\end{enumerate}
\end{exercice}

\begin{exercice}[Vrai ou faux ?]
Réponds par VRAI ou FAUX et justifie ta réponse par une phrase en français.
\begin{enumerate}
\item \zh{微信} (Wēixìn) est une application chinoise de messagerie très utilisée en Chine.
\item \zh{聊天} (liáotiān) signifie « dormir ».
\item \zh{发照片} (fā zhàopiàn) signifie « publier des photos ».
\item \zh{影响} (yǐngxiǎng) ne peut être qu'un nom, jamais un verbe.
\item \zh{安排} (ānpái) signifie « organiser, prévoir ».
\end{enumerate}
\end{exercice}

\begin{exercice}[Vocabulaire : trouve le mot]
Complète chaque phrase française par le mot chinois qui convient (en caractères et en pinyin).
\begin{enumerate}
\item Je passe deux heures par jour sur \ldots\ldots\ldots\ldots (internet).
\item Mon \ldots\ldots\ldots\ldots (téléphone portable) est dans mon sac.
\item Elle m'a envoyé un \ldots\ldots\ldots\ldots (message) hier soir.
\item Les jeunes aiment regarder des \ldots\ldots\ldots\ldots (vidéos) en ligne.
\item Trop d'écran, c'est mauvais pour la \ldots\ldots\ldots\ldots (santé).
\end{enumerate}
\end{exercice}

\begin{exercice}[Pinyin : associe les caractères à leur lecture]
Relie chaque mot en caractères à son pinyin.
\begin{center}
\begin{tabularx}{\linewidth}{|X|X|}
\hline
\rowcolor{popblueL} Caractères & Pinyin \\
\hline
1. 网络 & a. liáotiān \\
\hline
2. 聊天 & b. fāngbiàn \\
\hline
3. 方便 & c. wǎngluò \\
\hline
4. 视频 & d. ānpái \\
\hline
5. 安排 & e. shìpín \\
\hline
6. 网友 & f. wǎngyǒu \\
\hline
\end{tabularx}
\end{center}
\end{exercice}

\courssub{Je m'entraîne}

\begin{exercice}[Traduction français-chinois]
Traduis les phrases suivantes en chinois (caractères et pinyin).
\begin{enumerate}
\item Internet rend notre vie plus pratique, mais passer trop de temps en ligne n'est pas bon pour la santé.
\item J'utilise mon téléphone portable pour discuter avec mes amis.
\item Ma sœur publie souvent des photos sur son fil d'actualité.
\item À Yaoundé, de plus en plus de jeunes achètent des choses sur internet.
\end{enumerate}
\end{exercice}

\begin{exercice}[Traduction chinois-français]
Traduis les phrases suivantes en français.
\begin{enumerate}
\item \zh{很多年轻人每天都在朋友圈发照片和视频。}\\
\emph{Hěn duō niánqīngrén měitiān dōu zài péngyouquān fā zhàopiàn hé shìpín.}
\item \zh{我觉得上网聊天很方便，但是不能影响学习。}\\
\emph{Wǒ juéde shàngwǎng liáotiān hěn fāngbiàn, dànshì bù néng yǐngxiǎng xuéxí.}
\item \zh{他通过网络认识了很多中国朋友。}\\
\emph{Tā tōngguò wǎngluò rènshi le hěn duō Zhōngguó péngyou.}
\end{enumerate}
\end{exercice}

\begin{exercice}[Texte à trous]
Complète le texte suivant avec les mots de la banque (chaque mot ne sert qu'une fois) : \zh{上网、朋友圈、网友、方便、影响、安排} (shàngwǎng, péngyouquān, wǎngyǒu, fāngbiàn, yǐngxiǎng, ānpái).

\begin{textebox}[Texte à compléter]
现在，手机让我们的生活很 \ldots\ldots\ldots\ldots 。我每天用微信 \ldots\ldots\ldots\ldots ，看新闻，也在 \ldots\ldots\ldots\ldots 发照片。我在网上认识了一个 \ldots\ldots\ldots\ldots ，他住在上海。妈妈说，玩手机太多会 \ldots\ldots\ldots\ldots 学习，所以我 \ldots\ldots\ldots\ldots 好时间：晚上八点半以后不玩手机。\\
\emph{Xiànzài, shǒujī ràng wǒmen de shēnghuó hěn \ldots\ldots\ldots\ldots . Wǒ měitiān yòng Wēixìn \ldots\ldots\ldots\ldots , kàn xīnwén, yě zài \ldots\ldots\ldots\ldots fā zhàopiàn. Wǒ zài wǎngshang rènshi le yí ge \ldots\ldots\ldots\ldots , tā zhù zài Shànghǎi. Māma shuō, wán shǒujī tài duō huì \ldots\ldots\ldots\ldots xuéxí, suǒyǐ wǒ \ldots\ldots\ldots\ldots hǎo shíjiān : wǎnshang bā diǎn bàn yǐhòu bù wán shǒujī.}
\end{textebox}
\end{exercice}

\begin{exercice}[Remise en ordre et appariement]
\textbf{A.} Remets les mots dans le bon ordre pour faire une phrase correcte (écris la phrase en caractères).
\begin{enumerate}
\item \zh{在 / 我 / 照片 / 朋友圈 / 发 / 常常}
\item \zh{很 / 上网 / 方便 / 觉得 / 我 / 聊天}
\item \zh{影响 / 手机 / 太多 / 玩 / 会 / 健康}
\end{enumerate}
\textbf{B.} Relie chaque caractère à son radical (clé) et à sa traduction.
\begin{center}
\begin{tabularx}{\linewidth}{|X|X|X|}
\hline
\rowcolor{popblueL} Caractère & Radical & Traduction \\
\hline
1. 网 & a. 耳 (oreille) & i. pratique \\
\hline
2. 聊 & b. 罒 (filet) & ii. organiser \\
\hline
3. 便 & c. 扌 (main) & iii. ombre ; influence \\
\hline
4. 排 & d. 亻(homme) & iv. filet ; réseau \\
\hline
5. 影 & e. 彡 (ornement) & v. bavarder \\
\hline
\end{tabularx}
\end{center}
\end{exercice}

\courssub{Je me prépare au Bac}

\begin{exercice}[Type Bac — Compréhension écrite]
Lis le texte suivant, puis réponds aux questions \textbf{en chinois, par des phrases complètes}.

\begin{textebox}[微信和中国年轻人]
现在，微信是中国年轻人最喜欢的手机应用。每天早上，很多年轻人一起床就看微信：他们看朋友发来的消息，也看朋友圈里的照片和视频。微信很方便，不但可以聊天，还可以付钱、买东西、叫出租车。有些年轻人还在网上认识了新朋友。\\
但是，也有问题。有的年轻人每天花三四个小时玩手机，这影响了他们的学习和健康。一位老师说：「手机是好东西，可是我们要学会安排时间。」现在，很多学校告诉学生：上课不要看手机，晚上不要太晚睡。年轻人应该好好利用网络，不能让网络利用自己。\\
\emph{Xiànzài, Wēixìn shì Zhōngguó niánqīngrén zuì xǐhuan de shǒujī yìngyòng. Měitiān zǎoshang, hěn duō niánqīngrén yì qǐchuáng jiù kàn Wēixìn : tāmen kàn péngyou fā lái de xiāoxi, yě kàn péngyouquān lǐ de zhàopiàn hé shìpín. Wēixìn hěn fāngbiàn, bùdàn kěyǐ liáotiān, hái kěyǐ fùqián, mǎi dōngxi, jiào chūzūchē. Yǒuxiē niánqīngrén hái zài wǎngshang rènshi le xīn péngyou.}\\\emph{Dànshì, yě yǒu wèntí. Yǒu de niánqīngrén měitiān huā sān sì ge xiǎoshí wán shǒujī, zhè yǐngxiǎng le tāmen de xuéxí hé jiànkāng. Yí wèi lǎoshī shuō : « Shǒujī shì hǎo dōngxi, kěshì wǒmen yào xuéhuì ānpái shíjiān. » Xiànzài, hěn duō xuéxiào gàosu xuésheng : shàngkè bú yào kàn shǒujī, wǎnshang bú yào tài wǎn shuì. Niánqīngrén yīnggāi hǎohao lìyòng wǎngluò, bù néng ràng wǎngluò lìyòng zìjǐ.}\\
Traduction : Aujourd'hui, WeChat est l'application préférée des jeunes Chinois. Chaque matin, beaucoup de jeunes regardent WeChat dès le réveil : ils lisent les messages envoyés par leurs amis et regardent les photos et vidéos du fil d'actualité. WeChat est très pratique : on peut non seulement discuter, mais aussi payer, acheter des choses et appeler un taxi. Certains jeunes se font aussi de nouveaux amis en ligne.\\
Mais il y a aussi des problèmes. Certains jeunes passent trois ou quatre heures par jour sur leur téléphone, ce qui nuit à leurs études et à leur santé. Un professeur dit : « Le téléphone est une bonne chose, mais nous devons apprendre à organiser notre temps. » Aujourd'hui, beaucoup d'écoles disent aux élèves : ne regardez pas votre téléphone en classe, ne vous couchez pas trop tard. Les jeunes doivent bien utiliser internet, et ne pas se laisser utiliser par internet.
\end{textebox}

Questions :
\begin{enumerate}
\item \zh{中国年轻人早上起来以后常常做什么？}
\item \zh{微信除了聊天，还可以做什么？说出两个。}
\item \zh{玩手机太多有什么不好？}
\item \zh{老师对学生有什么建议？}
\item \zh{根据课文，年轻人应该怎么使用网络？}
\end{enumerate}
Questions de langue :
\begin{enumerate}
\item Relève dans le texte deux mots du glossaire de la leçon et traduis-les en français.
\item Dans la phrase \zh{这影响了他们的学习和健康}, quel est le rôle de \zh{了} ?
\end{enumerate}
\end{exercice}

\begin{exercice}[Type Bac — Thème et version]
\textbf{A. Thème.} Traduis en chinois (caractères obligatoires, pinyin en dessous) :
\begin{enumerate}
\item Internet rend notre vie plus pratique, mais passer trop de temps en ligne n'est pas bon pour la santé.
\item Mon ami camerounais utilise WeChat tous les jours pour parler avec ses amis chinois.
\item Les réseaux sociaux influencent la vie des jeunes : il faut bien organiser son temps.
\end{enumerate}
\textbf{B. Version.} Traduis en français :
\begin{enumerate}
\item \zh{很多年轻人每天都在朋友圈发照片和视频。}\\
\emph{Hěn duō niánqīngrén měitiān dōu zài péngyouquān fā zhàopiàn hé shìpín.}
\item \zh{网络让我们的生活更方便，也让我们认识世界各地的朋友。}\\
\emph{Wǎngluò ràng wǒmen de shēnghuó gèng fāngbiàn, yě ràng wǒmen rènshi shìjiè gèdì de péngyou.}
\end{enumerate}
\end{exercice}

\begin{exercice}[Type Bac — Expression écrite guidée]
\textbf{Sujet :} \zh{手机和我的生活} (Le téléphone portable et ma vie).

Rédige un court texte en chinois de \textbf{6 à 8 phrases} (environ 80 à 120 caractères). Consignes obligatoires :
\begin{itemize}
\item utilise au moins \textbf{5 mots} du glossaire de la leçon (par exemple : \zh{上网、朋友圈、网友、方便、影响、安排、聊天、视频、健康}) ;
\item utilise au moins \textbf{un connecteur} : \zh{因为……所以……} (parce que... donc...), \zh{虽然……但是……} (bien que... mais...) ou \zh{不但……而且……} (non seulement... mais aussi...) ;
\item termine par une phrase qui donne ton avis personnel avec \zh{我觉得……} (je pense que...).
\end{itemize}
Tu peux t'inspirer de la situation suivante : tu es élève en Terminale à Douala ; tu utilises ton téléphone pour étudier, discuter avec tes amis et suivre l'actualité, mais tes parents pensent que tu passes trop de temps en ligne.
\end{exercice}

\newpage

\coursec{Corrigés des exercices}

\begin{corrige}[Exercice 1 — QCM : le vocabulaire des réseaux sociaux]
\begin{enumerate}
\item Réponse \textbf{b}. \zh{朋友圈} (péngyouquān), littéralement « cercle d'amis », désigne le fil d'actualité de WeChat où l'on publie photos et messages. \zh{朋友} = ami, \zh{圈} = cercle.
\item Réponse \textbf{a}. \zh{网友} (wǎngyǒu) = \zh{网} (internet) + \zh{友} (ami) : un ami rencontré sur internet.
\item Réponse \textbf{b}. \zh{上网} (shàngwǎng) = \zh{上} (monter, accéder) + \zh{网} (réseau) : aller sur internet, se connecter.
\item Réponse \textbf{b}. \zh{方便} (fāngbiàn) est un adjectif : « pratique, commode ». Exemple : \zh{微信很方便。} (Wēixìn hěn fāngbiàn.) « WeChat est très pratique. »
\end{enumerate}
\textbf{Point de vigilance :} ne pas confondre \zh{上网} (aller sur internet) et \zh{网上} (wǎngshang, « sur internet », complément de lieu : \zh{我在网上买东西} « j'achète des choses sur internet »).
\end{corrige}

\begin{corrige}[Exercice 2 — Vrai ou faux ?]
\begin{enumerate}
\item \textbf{VRAI.} \zh{微信} (Wēixìn) est l'application de messagerie la plus utilisée en Chine ; elle sert aussi à payer et à publier des photos.
\item \textbf{FAUX.} \zh{聊天} (liáotiān) signifie « bavarder, discuter » ; « dormir » se dit \zh{睡觉} (shuìjiào).
\item \textbf{VRAI.} \zh{发} (fā) signifie « envoyer, publier » et \zh{照片} (zhàopiàn) « photo » : \zh{发照片} = publier/envoyer des photos.
\item \textbf{FAUX.} \zh{影响} (yǐngxiǎng) peut être un nom (« une influence ») ou un verbe (« influencer, nuire à ») : \zh{玩手机太多会影响学习。} « Jouer trop avec le téléphone nuit aux études. »
\item \textbf{VRAI.} \zh{安排} (ānpái) signifie « organiser, prévoir, planifier » : \zh{安排时间} « organiser son temps ».
\end{enumerate}
\end{corrige}

\begin{corrige}[Exercice 3 — Vocabulaire : trouve le mot]
\begin{enumerate}
\item Je passe deux heures par jour sur \textbf{\zh{网络}} (wǎngluò) — internet. (On accepte aussi \zh{网上} wǎngshang.)
\item Mon \textbf{\zh{手机}} (shǒujī) — téléphone portable — est dans mon sac.
\item Elle m'a envoyé un \textbf{\zh{消息}} (xiāoxi) — message — hier soir.
\item Les jeunes aiment regarder des \textbf{\zh{视频}} (shìpín) — vidéos — en ligne.
\item Trop d'écran, c'est mauvais pour la \textbf{\zh{健康}} (jiànkāng) — santé.
\end{enumerate}
\textbf{Point de vigilance :} attention aux tons : \zh{手机} shǒujī (3\textsuperscript{e} + 1\textsuperscript{er} ton), \zh{视频} shìpín (4\textsuperscript{e} + 2\textsuperscript{e} ton). Ne pas confondre \zh{消息} (message, nouvelle) et \zh{信息} (xìnxī, information).
\end{corrige}

\begin{corrige}[Exercice 4 — Pinyin : associe les caractères à leur lecture]
\begin{center}
\begin{tabularx}{\linewidth}{|X|X|X|}
\hline
\rowcolor{popblueL} Caractères & Pinyin & Traduction \\
\hline
1. 网络 & c. wǎngluò & internet, réseau \\
\hline
2. 聊天 & a. liáotiān & bavarder, discuter \\
\hline
3. 方便 & b. fāngbiàn & pratique, commode \\
\hline
4. 视频 & e. shìpín & vidéo \\
\hline
5. 安排 & d. ānpái & organiser, prévoir \\
\hline
6. 网友 & f. wǎngyǒu & ami rencontré sur internet \\
\hline
\end{tabularx}
\end{center}
\textbf{Astuce de mémorisation :} les mots commençant par \zh{网} (wǎng, réseau) se prononcent tous « wǎng... » : \zh{网络} wǎngluò, \zh{网友} wǎngyǒu, \zh{上网} shàngwǎng. Le radical 罒 (filet), visible en haut de \zh{网}, aide à reconnaître la famille du « réseau ».
\end{corrige}

\begin{corrige}[Exercice 5 — Traduction français-chinois]
\begin{enumerate}
\item \zh{网络让我们的生活更方便，但是上网时间太长对健康不好。}\\
\emph{Wǎngluò ràng wǒmen de shēnghuó gèng fāngbiàn, dànshì shàngwǎng shíjiān tài cháng duì jiànkāng bù hǎo.}\\
Structure : \zh{让} + personne + adjectif (« rendre... ») ; \zh{对……不好} « ne pas être bon pour... ».
\item \zh{我用手机和朋友聊天。}\\
\emph{Wǒ yòng shǒujī hé péngyou liáotiān.}\\
Structure : Sujet + \zh{用} + outil + \zh{和} + personne + verbe.
\item \zh{我姐姐常常在朋友圈发照片。}\\
\emph{Wǒ jiějie chángcháng zài péngyouquān fā zhàopiàn.}\\
Le complément de lieu \zh{在朋友圈} se place avant le verbe \zh{发}.
\item \zh{在雅温得，越来越多的年轻人在网上买东西。}\\
\emph{Zài Yǎwēndé, yuè lái yuè duō de niánqīngrén zài wǎngshang mǎi dōngxi.}\\
\zh{越来越多} (yuè lái yuè duō) = « de plus en plus ».
\end{enumerate}
\textbf{Points de vigilance :} en chinois, le complément de lieu introduit par \zh{在} se place \emph{avant} le verbe (contrairement au français) ; \zh{和} (avec) précède aussi le verbe : \zh{和朋友聊天}, jamais \zh{聊天和朋友}.
\end{corrige}

\begin{corrige}[Exercice 6 — Traduction chinois-français]
\begin{enumerate}
\item \zh{很多年轻人每天都在朋友圈发照片和视频。}\\
« Beaucoup de jeunes publient chaque jour des photos et des vidéos sur leur fil d'actualité. »\\
Note : \zh{都} (dōu, « tous ») renforce \zh{每天} ; il se place avant le verbe.
\item \zh{我觉得上网聊天很方便，但是不能影响学习。}\\
« Je trouve que discuter en ligne est très pratique, mais cela ne doit pas nuire aux études. »\\
Note : \zh{觉得} (juéde) = « penser, trouver » ; \zh{影响} est ici un verbe.
\item \zh{他通过网络认识了很多中国朋友。}\\
« Il s'est fait beaucoup d'amis chinois grâce à internet. »\\
Note : \zh{通过} (tōngguò) = « par l'intermédiaire de » ; \zh{认识了} marque le résultat accompli de la rencontre.
\end{enumerate}
\end{corrige}

\begin{corrige}[Exercice 7 — Texte à trous]
Texte complété :
\begin{enumerate}
\item \zh{方便} (fāngbiàn) : \zh{手机让我们的生活很方便。} « Le téléphone rend notre vie très pratique. » (adjectif après \zh{很})
\item \zh{上网} (shàngwǎng) : \zh{我每天用微信上网，看新闻} « Chaque jour, j'utilise WeChat pour aller sur internet et lire les nouvelles. »
\item \zh{朋友圈} (péngyouquān) : \zh{也在朋友圈发照片} « et je publie aussi des photos sur mon fil d'actualité. » (\zh{在} + lieu + verbe)
\item \zh{网友} (wǎngyǒu) : \zh{我在网上认识了一个网友} « Je me suis fait un ami sur internet. » (classificateur \zh{个})
\item \zh{影响} (yǐngxiǎng) : \zh{玩手机太多会影响学习} « trop jouer avec le téléphone nuit aux études. » (verbe après \zh{会})
\item \zh{安排} (ānpái) : \zh{所以我安排好时间} « donc j'organise bien mon temps. »
\end{enumerate}
\textbf{Méthode :} repérer d'abord la nature du mot attendu (adjectif après \zh{很}, verbe après \zh{会}, nom après un classificateur), puis vérifier le sens global de la phrase.
\end{corrige}

\begin{corrige}[Exercice 8 — Remise en ordre et appariement]
\textbf{A. Remise en ordre}
\begin{enumerate}
\item \zh{我常常在朋友圈发照片。}\\
\emph{Wǒ chángcháng zài péngyouquān fā zhàopiàn.} « Je publie souvent des photos sur mon fil d'actualité. »\\
Ordre : Sujet + adverbe de fréquence \zh{常常} + complément de lieu \zh{在朋友圈} + verbe + complément.
\item \zh{我觉得上网聊天很方便。}\\
\emph{Wǒ juéde shàngwǎng liáotiān hěn fāngbiàn.} « Je trouve que discuter en ligne est très pratique. »\\
\zh{上网聊天} est le sujet de la subordonnée ; l'adjectif \zh{方便} est précédé de \zh{很}.
\item \zh{玩太多手机会影响健康。}\\
\emph{Wán tài duō shǒujī huì yǐngxiǎng jiànkāng.} « Jouer trop avec le téléphone nuit à la santé. »\\
(On accepte aussi \zh{玩手机太多会影响健康。}) Le sujet est la proposition verbale \zh{玩太多手机} ; \zh{会} précède le verbe \zh{影响}.
\end{enumerate}
\textbf{B. Appariement caractère — radical — traduction}
\begin{center}
\begin{tabularx}{\linewidth}{|X|X|X|}
\hline
\rowcolor{popblueL} Caractère & Radical & Traduction \\
\hline
1. 网 & b. 罒 (filet) & iv. filet ; réseau \\
\hline
2. 聊 & a. 耳 (oreille) & v. bavarder \\
\hline
3. 便 & d. 亻(homme) & i. pratique \\
\hline
4. 排 & c. 扌 (main) & ii. organiser \\
\hline
5. 影 & e. 彡 (ornement) & iii. ombre ; influence \\
\hline
\end{tabularx}
\end{center}
\textbf{Remarque mnémotechnique :} \zh{聊} contient la clé de l'oreille 耳 : bavarder, c'est écouter et parler ; \zh{排} contient la clé de la main 扌 : organiser, c'est « mettre en rang avec les mains ».
\end{corrige}

\begin{corrige}[Exercice 9 — Type Bac : compréhension écrite]
\textbf{Barème indicatif (sur 20) :} questions de compréhension 5 × 2 points = 10 ; questions de langue 2 × 2 points = 4 ; qualité de la langue (caractères, grammaire) = 6.

\textbf{Réponses attendues (en chinois, phrases complètes) :}
\begin{enumerate}
\item \zh{中国年轻人早上起来以后常常看微信：他们看朋友发来的消息，也看朋友圈里的照片和视频。}\\
\emph{Zhōngguó niánqīngrén zǎoshang qǐlái yǐhòu chángcháng kàn Wēixìn : tāmen kàn péngyou fā lái de xiāoxi, yě kàn péngyouquān lǐ de zhàopiàn hé shìpín.}\\
« Après s'être levés, les jeunes Chinois regardent souvent WeChat : les messages de leurs amis et les photos et vidéos du fil d'actualité. »
\item \zh{微信除了聊天，还可以付钱、买东西、叫出租车。}\\
\emph{Wēixìn chúle liáotiān, hái kěyǐ fùqián, mǎi dōngxi, jiào chūzūchē.}\\
« En plus de discuter, on peut aussi payer, acheter des choses, appeler un taxi. » (deux des trois réponses suffisent)
\item \zh{玩手机太多会影响他们的学习和健康。}\\
\emph{Wán shǒujī tài duō huì yǐngxiǎng tāmen de xuéxí hé jiànkāng.}\\
« Trop jouer avec le téléphone nuit à leurs études et à leur santé. »
\item \zh{老师建议学生学会安排时间。}\\
\emph{Lǎoshī jiànyì xuésheng xuéhuì ānpái shíjiān.}\\
« Le professeur conseille aux élèves d'apprendre à organiser leur temps. »
\item \zh{年轻人应该好好利用网络，不能让网络利用自己。}\\
\emph{Niánqīngrén yīnggāi hǎohao lìyòng wǎngluò, bù néng ràng wǎngluò lìyòng zìjǐ.}\\
« Les jeunes doivent bien utiliser internet et ne pas se laisser utiliser par internet. »
\end{enumerate}
\textbf{Questions de langue :}
\begin{enumerate}
\item Exemples : \zh{朋友圈} (péngyouquān) = « fil d'actualité (cercle d'amis) » ; \zh{方便} (fāngbiàn) = « pratique » ; \zh{影响} (yǐngxiǎng) = « influencer, nuire à » ; \zh{安排} (ānpái) = « organiser ». (Deux mots suffisent.)
\item Dans \zh{这影响了他们的学习和健康}, le \zh{了} placé après le verbe \zh{影响} est la \textbf{marque de l'accompli} : il indique que l'action a produit un résultat réel (« cela a nui à... »). Ce n'est pas un marqueur de passé au sens français, mais un indice que l'action est réalisée.
\end{enumerate}
\textbf{Points de vigilance :} répondre par des phrases complètes (sujet + verbe), recopier les caractères sans faute, ne pas traduire mot à mot le français dans la réponse.
\end{corrige}

\begin{corrige}[Exercice 10 — Type Bac : thème et version]
\textbf{Barème indicatif (sur 20) :} thème 3 × 4 points = 12 ; version 2 × 3 points = 6 ; présentation (pinyin, caractères soignés) = 2.

\textbf{A. Thème — propositions de traduction :}
\begin{enumerate}
\item \zh{网络让我们的生活更方便，但是上网时间太长对健康不好。}\\
\emph{Wǎngluò ràng wǒmen de shēnghuó gèng fāngbiàn, dànshì shàngwǎng shíjiān tài cháng duì jiànkāng bù hǎo.}
\item \zh{我的喀麦隆朋友每天用微信和他的中国朋友聊天。}\\
\emph{Wǒ de Kāmàilóng péngyou měitiān yòng Wēixìn hé tā de Zhōngguó péngyou liáotiān.}
\item \zh{社交网络影响年轻人的生活，所以我们要好好安排时间。}\\
\emph{Shèjiāo wǎngluò yǐngxiǎng niánqīngrén de shēnghuó, suǒyǐ wǒmen yào hǎohao ānpái shíjiān.}
\end{enumerate}
\textbf{B. Version :}
\begin{enumerate}
\item \zh{很多年轻人每天都在朋友圈发照片和视频。}\\
« Beaucoup de jeunes publient chaque jour des photos et des vidéos sur leur fil d'actualité. »
\item \zh{网络让我们的生活更方便，也让我们认识世界各地的朋友。}\\
« Internet rend notre vie plus pratique et nous permet aussi de connaître des amis du monde entier. »\\
Note : \zh{世界各地} (shìjiè gèdì) = « de partout dans le monde » ; la structure \zh{让……也……} coordonne deux résultats.
\end{enumerate}
\textbf{Points de vigilance :} en thème, placer les compléments de moyen (\zh{用微信}) et de lieu (\zh{在朋友圈}) \emph{avant} le verbe ; ne pas oublier \zh{的} dans \zh{我们的生活} ; en version, rendre \zh{都} par « tous / chaque » selon le contexte.
\end{corrige}

\begin{corrige}[Exercice 11 — Type Bac : expression écrite guidée]
\textbf{Barème indicatif (sur 20) :} respect des consignes (6 à 8 phrases, 5 mots du glossaire, 1 connecteur, avis final avec \zh{我觉得}) = 6 ; correction grammaticale = 6 ; exactitude des caractères = 4 ; richesse et cohérence = 4.

\textbf{Proposition de rédaction modèle :}

\zh{手机和我的生活}\\
\zh{我是杜阿拉的一名中学生。手机对我很重要：我每天用它上网、学习，也用微信和朋友聊天。周末，我常常在朋友圈发照片，有时候还看视频。因为手机很方便，所以我能学到很多东西。虽然手机很有用，但是我父母说我玩得太多了。他们担心手机会影响我的学习和健康。现在我安排好时间：上课的时候不看手机，晚上十点以后也不玩。我觉得手机是好东西，但是我们要好好利用它，不能让它利用我们。}\\
\emph{Wǒ shì Dù'ālā de yí míng zhōngxuéshēng. Shǒujī duì wǒ hěn zhòngyào : wǒ měitiān yòng tā shàngwǎng, xuéxí, yě yòng Wēixìn hé péngyou liáotiān. Zhōumò, wǒ chángcháng zài péngyouquān fā zhàopiàn, yǒu shíhou hái kàn shìpín. Yīnwèi shǒujī hěn fāngbiàn, suǒyǐ wǒ néng xué dào hěn duō dōngxi. Suīrán shǒujī hěn yǒuyòng, dànshì wǒ fùmǔ shuō wǒ wán de tài duō le. Tāmen dānxīn shǒujī huì yǐngxiǎng wǒ de xuéxí hé jiànkāng. Xiànzài wǒ ānpái hǎo shíjiān : shàngkè de shíhou bú kàn shǒujī, wǎnshang shí diǎn yǐhòu yě bù wán. Wǒ juéde shǒujī shì hǎo dōngxi, dànshì wǒmen yào hǎohao lìyòng tā, bù néng ràng tā lìyòng wǒmen.}

Traduction : « Je suis un élève de Douala. Le téléphone est très important pour moi : je m'en sers chaque jour pour aller sur internet, étudier, et aussi pour discuter avec mes amis sur WeChat. Le week-end, je publie souvent des photos sur mon fil d'actualité et je regarde parfois des vidéos. Parce que le téléphone est très pratique, je peux apprendre beaucoup de choses. Bien que le téléphone soit très utile, mes parents disent que j'y passe trop de temps. Ils craignent qu'il ne nuise à mes études et à ma santé. Maintenant, j'organise bien mon temps : je ne regarde pas mon téléphone en classe, et après dix heures du soir je n'y touche plus. Je pense que le téléphone est une bonne chose, mais nous devons bien l'utiliser, sans nous laisser utiliser par lui. »

\textbf{Vérification des consignes :}
\begin{itemize}
\item Mots du glossaire employés (au moins 5) : \zh{上网、聊天、朋友圈、视频、方便、影响、安排、健康} — 8 mots.
\item Connecteurs : \zh{因为……所以……} et \zh{虽然……但是……} — deux connecteurs.
\item Avis personnel final : \zh{我觉得手机是好东西……}.
\item Longueur : 8 phrases, environ 150 caractères (légèrement au-dessus du minimum, ce qui est valorisé au Bac si la langue reste correcte).
\end{itemize}
\textbf{Points de vigilance :}
\begin{itemize}
\item Ne pas écrire \zh{我很重要手机} : la structure est \zh{手机对我很重要} (« le téléphone, pour moi, est important »).
\item \zh{玩得太多了} : avec \zh{得} après le verbe pour introduire le complément de degré ; ne pas confondre avec \zh{的}.
\item Le complément de temps \zh{晚上十点以后} se place avant le verbe.
\item Éviter les phrases trop longues sans ponctuation ; une phrase = une idée.
\end{itemize}
\end{corrige}

\newpage

\coursec{Fiche bilan}

\begin{popfigure}
\begin{tikzpicture}[mindmap, grow cyclic, every node/.style={concept}, text=white,
root concept/.append style={concept color=popblue, font=\large\bfseries},
level 1/.append style={level distance=3.4cm, sibling angle=60},
level 1 concept/.append style={font=\small}]
\node[root concept, align=center]{互联网\\ \emph{hùliánwǎng}\\ Internet}
  child[concept color=poporange]{ node[align=center]{词汇\\ Lexique\\ 网络 · 手机 · 网站} }
  child[concept color=popgreen]{ node[align=center]{社交网络\\ Réseaux sociaux\\ 微信 · 微博 · 朋友圈} }
  child[concept color=poppurple]{ node[align=center]{语法\\ Grammaire\\ 用 + 名\\ 越来越 + adj.} }
  child[concept color=popgold]{ node[align=center]{汉字\\ Caractères\\ 网 · 络 · 信 · 微} }
  child[concept color=poppink]{ node[align=center]{理解\\ Compréhension\\ lire · répondre} }
  child[concept color=popgreen!70!black]{ node[align=center]{文化\\ Culture\\ Cameroun / Chine} };
\end{tikzpicture}
\legende{Carte mentale de la leçon : internet et les réseaux sociaux.}
\end{popfigure}

\begin{bilan}
\textbf{1. Lexique clé à connaître par cœur}
\begin{itemize}
  \item \zh{互联网} \emph{hùliánwǎng} : internet ; \zh{上网} \emph{shàngwǎng} : aller sur internet ; \zh{网络} \emph{wǎngluò} : réseau.
  \item \zh{手机} \emph{shǒujī} : téléphone portable ; \zh{电脑} \emph{diànnǎo} : ordinateur ; \zh{网站} \emph{wǎngzhàn} : site web.
  \item \zh{社交网络} \emph{shèjiāo wǎngluò} : réseaux sociaux ; \zh{微信} \emph{Wēixìn} : WeChat ; \zh{朋友圈} \emph{péngyouquān} : fil d'actualité (Moments).
  \item \zh{发消息} \emph{fā xiāoxi} : envoyer un message ; \zh{照片} \emph{zhàopiàn} : photo ; \zh{视频} \emph{shìpín} : vidéo.
  \item \zh{信息} \emph{xìnxī} : information ; \zh{方便} \emph{fāngbiàn} : pratique ; \zh{危险} \emph{wēixiǎn} : dangereux ; \zh{注意} \emph{zhùyì} : faire attention.
\end{itemize}

\textbf{2. Structures de grammaire à maîtriser}
\begin{center}
\begin{tabularx}{\linewidth}{|X|X|X|}
\hline
\rowcolor{popblueL} Structure & Exemple (caractères + pinyin) & Traduction \\
\hline
用 + nom + verbe (instrument) & \zh{我用手机上网。} \emph{Wǒ yòng shǒujī shàngwǎng.} & J'utilise mon portable pour aller sur internet. \\
\hline
越来越 + adjectif & \zh{上网的人越来越多。} \emph{Shàngwǎng de rén yuèláiyuè duō.} & De plus en plus de gens vont sur internet. \\
\hline
在 + lieu + verbe & \zh{他在网上买东西。} \emph{Tā zài wǎngshang mǎi dōngxi.} & Il achète des choses en ligne. \\
\hline
可以 / 不可以 + verbe & \zh{我们不可以在网上说坏话。} \emph{Wǒmen bù kěyǐ zài wǎngshang shuō huàihuà.} & Nous ne pouvons pas dire du mal en ligne. \\
\hline
\end{tabularx}
\end{center}

\textbf{3. Écriture des caractères}
\begin{itemize}
  \item \zh{网} (wǎng, « filet/réseau ») : clé \zh{冂} (enceinte) ; on trace d'abord le cadre extérieur, puis l'intérieur, de haut en bas.
  \item \zh{信} (xìn, « lettre/message ») : clé \zh{亻} (homme, à gauche) ; gauche avant droite.
  \item \zh{微} (wēi) : clé \zh{彳} (pas, à gauche) ; écrire de gauche à droite, en trois parties.
\end{itemize}

\textbf{4. Méthodes express}
\begin{itemize}
  \item Lire le texte deux fois : d'abord pour l'idée générale, ensuite en soulignant les mots du thème.
  \item Pour répondre aux questions : repérer dans le texte la phrase qui contient la réponse, puis reformuler simplement.
  \item Pour retenir le vocabulaire : grouper par familles (appareils / actions / réseaux sociaux / avantages et dangers).
\end{itemize}
\end{bilan}

\begin{aretenir}[Checklist avant le Bac]
\begin{itemize}
  \item $\square$ Je sais lire et prononcer le texte support avec les bons tons.
  \item $\square$ Je connais le sens de \zh{互联网}, \zh{网络}, \zh{手机}, \zh{网站}, \zh{社交网络}.
  \item $\square$ Je sais écrire \zh{网}, \zh{信}, \zh{微} avec le bon ordre des traits.
  \item $\square$ Je connais les clés \zh{冂}, \zh{亻} et \zh{彳}.
  \item $\square$ Je sais construire une phrase avec \zh{用} + instrument (\zh{我用手机发消息。}).
  \item $\square$ Je sais employer \zh{越来越} + adjectif (\zh{越来越方便。}).
  \item $\square$ Je place correctement le complément de lieu avec \zh{在} avant le verbe.
  \item $\square$ Je sais dire les avantages et les dangers d'internet (\zh{方便} / \zh{危险}, \zh{注意}).
  \item $\square$ Je sais répondre aux questions de compréhension en phrases complètes.
  \item $\square$ Je peux comparer l'usage d'internet au Cameroun et en Chine (WeChat, WhatsApp).
  \item $\square$ Je mets la ponctuation chinoise correcte ：，。 ？
  \item $\square$ Je révise le pinyin avec les tons sur les voyelles, jamais avec des chiffres.
\end{itemize}
\end{aretenir}

\findecours

\end{document}
